\documentclass[11pt,a4paper]{article}
\usepackage{fontspec}
\setmainfont{Latin Modern Roman}
\usepackage[margin=23mm]{geometry}
\usepackage{amsmath,amssymb,amsthm,mathtools,mathrsfs}
\usepackage{longtable,booktabs}
\usepackage[unicode,hidelinks]{hyperref}
\providecommand{\C}{\mathbb C}
\providecommand{\R}{\mathbb R}
\newtheorem{theorem}{Theorem}[section]
\newtheorem{lemma}[theorem]{Lemma}
\newtheorem{proposition}[theorem]{Proposition}
\newtheorem{corollary}[theorem]{Corollary}
\newtheorem{definition}[theorem]{Definition}
\newtheorem{remark}[theorem]{Remark}
\allowdisplaybreaks
\emergencystretch=3em
\makeatletter
\renewcommand\@pnumwidth{2.5em}
\renewcommand\@tocrmarg{3.5em}
\makeatother
\hypersetup{pdftitle={Original Mixed Cohomology: Periods, Kernels and Arithmetic Control},pdfauthor={Split-Zero research programme}}
\title{Original Mixed Cohomology\\Periods, Kernels and Arithmetic Control}
\author{Split-Zero research programme}
\date{14 September 2026}
\begin{document}
\maketitle
\paragraph{Complete proof edition.}
The chapters below retain the original periods, amplitude units,
primary multiplicities, cyclic maps, source masses and four endpoint
metrics. Each chapter supplies its full calculation and proof. Earlier
finite bounds remain recorded in their original form; subsequent
chapters give their proved refinements. In particular KAF retains the
factorial denominator of OKV and improves the absolute kernel bound,
while AMT evaluates the source differences left in OMV.
\paragraph{Sources and scope.}
The complete preceding Gamma and observation proofs, including their
historical analytic providers, are supplied in the source closure.
The original all-multiplicity cyclic maps remain in OCS. The evaluated
five-member quadric orbit, exact kernel dimension, and its individual
volume conclusions state their simple-quartet period domain explicitly.
The full all-multiplicity operator estimate is given in KAF.
OCF, OKM and OKA calculate the actual kernel frame, its original
metrics and the boundary action. OMG proves the full all-multiplicity
Segre--Veronese geometry and retains the actual cyclic stabilizer.
\tableofcontents
\clearpage

\clearpage
% BEGIN COMPLETE OCS
% New original-observation calculation; the sealed cut22 is not edited.
\section{The literal cyclic sectors and the original mixed kernel}
\label{sec:OCS}

This calculation evaluates the cyclic projector already present in
the original period observation. Its coefficients retain the original
period integrals, amplitude unit, and all primary nilpotents. The
result includes an exact finite matrix for the actual rank, as well
as unconditional rank bounds obtained from the original symmetric
tensor geometry.

\subsection{Original data and the Fourier basis of the literal matrix}
Use the following data of OPG1--5, including the same contours and
their orientations:
\[
 \begin{gathered}
 \alpha_{\epsilon,\eta}=\tfrac12+\epsilon\delta+i\eta\gamma,
 \quad\epsilon,\eta\in\{+1,-1\},\quad0<\delta<\tfrac12,\quad\gamma>2,\\
 h(z_1)=\prod_{\epsilon,\eta}(z_1-\alpha_{\epsilon,\eta})^m,
 \quad n=4m,\quad N=n+1,\quad E_h=\mathbb C[z_1]/(h),\\
 k=4l+1,\quad l,m\geq1,\quad c=k/2,
 \quad e=1+k(m-1),\quad q=e(k+1)^2,\\
 \beta_{a,b}=c+(2a-k)\delta+i(2b-k)\gamma,
 \quad\chi(S)=\prod_{a,b=0}^k(S-\beta_{a,b})^e,
 \quad E=\mathbb C[S]/(\chi).
 \end{gathered}\tag{OCS1}
\]
Here $A_h$ and $A$ are multiplication by $z_1$ and $S$ on their
respective algebras, and $Y=(A-cI)/i$. The original amplitude
unit is $U=M_{j_hv_h}$, with $v_h=(2\xi)/h$ and every original
$v_h(\alpha)\ne0$. For the specified nonzero $u$ and its branch,
$\Phi_h'=h$, $\Phi_h(0)=0$, and
\[
 \Gamma_j=-\ell_0+\ell_j,\qquad
 \arg\ell_j=\frac{\arg u+\pi+2\pi j}{N},\qquad
 (\Pi[P])_j=\int_{\Gamma_j}e^{\Phi_h(z_1)/u}
                       \operatorname{rem}_hP(z_1)\,dz_1.
 \tag{OCS2}
\]
The representative has degree less than $n$. The convergence and
invertibility of this original period map are the complete period
theorems retained in OPG4. No period value or unit derivative is
chosen in the present calculation. The original injection and
observation are
\[
\begin{gathered}
 \iota[P]=[P(z_1+\cdots+z_k)],\quad
 \mathcal V=\Pi^{\otimes k}U^{\otimes k}\iota,
 \\ P_k=\frac1N\sum_{a=0}^{N-1}(C^{-a})^{\otimes k},
 \\ L=P_k\mathcal V=\jmath\Lambda,
 \quad B=\operatorname{im}L,\quad K=\ker\Lambda.
\end{gathered}\tag{OCS3}
\]
The map $\jmath:B\hookrightarrow(\mathbb C^n)^{\otimes k}$
is the original image inclusion. The injection $\iota$ has precisely
the defining polynomial in OCS1: on each ordered primary tensor
the sum of its $k$ independent nilpotents has index
$e=1+k(m-1)$, its top coefficient being
$(k(m-1))!/((m-1)!)^k\ne0$, and its distinct eigenvalues are
the displayed $\beta_{a,b}$. Thus $\mathcal V$ is injective.

On a period-coordinate column $x\in\mathbb C^{N-1}$ the literal
matrix in OPG4--5 is
\[
 (Cx)_j=x_{j+1}-x_1\quad(1\leq j<N-1),\qquad
 (Cx)_{N-1}=-x_1.
 \tag{OCS4}
\]
Set $\zeta=\exp(2\pi i/N)$ and, for $1\leq t\leq N-1$, let
\[
 (f_t)_j=\zeta^{jt}-1,\qquad
 F=[f_1\ \cdots\ f_{N-1}],\qquad
 (F^{-1}x)_t=\frac1N\sum_{j=1}^{N-1}\zeta^{-tj}x_j.
 \tag{OCS5}
\]
The last expression is the exact inverse. Indeed
$\sum_{j=1}^{N-1}\zeta^{-tj}(\zeta^{uj}-1)=N\mathbf1_{t=u}$
for $t,u\ne0$ modulo $N$, by the finite geometric sum.
The first $N-2$ rows of OCS4 give
$\zeta^{(j+1)t}-\zeta^t=\zeta^t(\zeta^{jt}-1)$, and the
last row gives $1-\zeta^t=\zeta^t(\zeta^{(N-1)t}-1)$.
Consequently $Cf_t=\zeta^tf_t$, proving the orientation of
every inverse power in OCS3, as well as $C^N=I$.
The original Euclidean Gram of this basis is
\[
 F^*F=N(I_{N-1}+\mathbf1\mathbf1^*).
 \tag{OCS6}
\]
To verify it, expand
$\sum_j(\zeta^{-jt}-1)(\zeta^{ju}-1)$ and use the same
geometric sums; its value is $N(\mathbf1_{t=u}+1)$.
Thus the complete coefficient and metric transformations are
both specified.

\subsection{The ordered-to-symmetric map and the exact sector projector}
Let
\[
 \mathcal A_k=\{\nu\in\mathbb Z_{\geq0}^{N-1}:\sum_t\nu_t=k\},
 \quad\operatorname{wt}(\nu)=\sum_{t=1}^{N-1}t\nu_t\pmod N,
 \quad\mathcal I_{k,\tau}=\{\nu\in\mathcal A_k:
                                      \operatorname{wt}(\nu)=\tau\}.
 \tag{OCS7}
\]
For each $\nu$ define the following unscaled symmetric tensor:
\[
 S_\nu=\sum_{\substack{(t_1,\ldots,t_k)\in\{1,\ldots,N-1\}^k\\
                   \#\{j:t_j=t\}=\nu_t\ \text{for each }t}}
                    f_{t_1}\otimes\cdots\otimes f_{t_k}.
 \tag{OCS8}
\]
Every distinct ordered word occurs once. Since the ordered tensors
in the $f_t$ form a basis, their disjoint word supports show that
the $S_\nu$ are linearly independent. A permutation-invariant tensor
has equal coefficients on each word orbit, so they span exactly
$\operatorname{Sym}^k(\mathbb C^{N-1})$, realized as that invariant
subspace of the full ordered tensor product. In particular, for
$x=(x_1,\ldots,x_{N-1})$,
\[
 (Fx)^{\otimes k}=\sum_{\nu\in\mathcal A_k}x^\nu S_\nu,
 \quad
 (C^{-a})^{\otimes k}S_\nu=
                         \zeta^{-a\operatorname{wt}(\nu)}S_\nu,
 \quad
 P_k(Fx)^{\otimes k}=\sum_{\nu\in\mathcal I_{k,0}}x^\nu S_\nu.
 \tag{OCS9}
\]
The finite average in OCS3 is exactly one on weight zero and
zero on every other weight, by its geometric sum. No multinomial
coefficient is omitted: the multiplicity $k!/\prod_t\nu_t!$
is already the number of summands in OCS8.

Write $\mathcal F_k:\mathbb C^{\mathcal A_k}\to
\operatorname{Sym}^k(\mathbb C^{N-1})$ for the basis isomorphism
$e_\nu\mapsto S_\nu$. Let $J_0$ be the coordinate inclusion of
$\mathbb C^{\mathcal I_{k,0}}$ into $\mathbb C^{\mathcal A_k}$,
$\pi_0$ the coordinate restriction in the opposite direction,
and let $\mathcal F_0=\mathcal F_kJ_0$, with its target the original
ordered tensor space. The exact projector factorization is
\[
 P_k\mathcal F_k=\mathcal F_0\pi_0.
 \tag{OCS10}
\]
All images of $\mathcal V$ are symmetric: polynomial evaluation
at a sum is invariant under every factor permutation, and the
same map $\Pi U$ is applied in all factors. Therefore OCS10
applies on the entire original image $\mathcal V(E)$.

For completeness, the full inherited metric is retained in these
coordinates. Put $G=F^*F$ as in OCS6. The Gram of the basis OCS8 is
\[
 \mathcal G_{\nu,\mu}=
 \sum_{\substack{a_{tu}\geq0\\\sum_u a_{tu}=\nu_t\ ,\
                                    \sum_t a_{tu}=\mu_u}}
       \frac{k!}{\prod_{t,u}a_{tu}!}\prod_{t,u}G_{tu}^{a_{tu}}.
 \tag{OCS11}
\]
Indeed the coefficient of $\bar x^\nu y^\mu$ in
$(\bar x^TGy)^k$ is the right side by the multinomial theorem.
Taking the inner product of the two pure tensors in OCS9 gives
the same coefficient as $\langle S_\nu,S_\mu\rangle$.
This proves OCS11, including its off-diagonal entries; it is
positive definite because OCS8 is a basis. On the full symmetric
coordinates the projector matrix is $J_0\pi_0$, and its adjoint
for this exact Gram is
$\mathcal G^{-1}(J_0\pi_0)^*\mathcal G$.
Thus the Fourier-sector computation makes no unproved
orthogonality or norm assertion about the original cyclic average.

\subsection{Every primary period coefficient of the projected curve}
For each original root $\alpha$ let $E_{\alpha,a}$, $0\leq a<m$,
be the degree-less-than-$4m$ primary representatives of OPG10.
Explicitly, if $h_\alpha=h/(z_1-\alpha)^m$, put
\[
 E_{\alpha,a}(z_1)=\operatorname{rem}_h\left[
 h_\alpha(z_1)\sum_{j=0}^{m-1}
       \frac{(h_\alpha^{-1})^{(j)}(\alpha)}{j!}
                  (z_1-\alpha)^{j+a}\right],\qquad
 b_{j;\alpha,a}=\int_{\Gamma_j}e^{\Phi_h(z_1)/u}
                                    E_{\alpha,a}(z_1)\,dz_1.
 \tag{OCS12}
\]
The four primary factors are coprime. The Taylor inverse gives
constant one on its own factor and zero on the other three;
thus the classes $[E_{\alpha,a}]$ form the full primary basis
with all $m$ nilpotent degrees at every root.
Define the exact finite Fourier and unit coefficients
\[
 \widehat b_{t;\alpha,a}=
     \frac1N\sum_{j=1}^{N-1}\zeta^{-tj}b_{j;\alpha,a},
 \qquad
 f_{t;\alpha,d}=\frac{i^{-d}}{d!}
       \sum_{a=d}^{m-1}\widehat b_{t;\alpha,a}
                         \frac{v_h^{(a-d)}(\alpha)}{(a-d)!},
 \quad0\leq d<m.
 \tag{OCS13}
\]
Let $\omega_\alpha=(\alpha-1/2)/i$ and set
\[
 \mathbf v(w)=\Pi U\exp\bigl(w(A_h-\tfrac12I)/i\bigr)[1],
 \quad x(w)=F^{-1}\mathbf v(w),\qquad
 x_t(w)=\sum_{\alpha}e^{\omega_\alpha w}
                          \sum_{d=0}^{m-1}f_{t;\alpha,d}w^d.
 \tag{OCS14}
\]
To verify the final equality, in the primary algebra at $\alpha$
multiply the Taylor series of $v_h$ by the terminating series of
$\exp(w(z_1-\alpha)/i)$. The coefficient on
$[E_{\alpha,a}]$ is
$e^{\omega_\alpha w}\sum_{d=0}^a
v_h^{(a-d)}(\alpha)(w/i)^d/((a-d)!d!)$.
Apply the original polynomial period map and the inverse in OCS5;
this gives precisely OCS13--14. In particular, the unreduced
entire function has passed through its full primary remainder
before any period is integrated.

For $\nu\in\mathcal I_{k,0}$, $0\leq a,b\leq k$ and
$0\leq D<e$, define $c_{\nu;a,b,D}$ by the following finite
multinomial rule. Sum over integers $n_{t,\alpha,d}\geq0$ satisfying
\[
 \begin{gathered}
 \sum_{\alpha,d}n_{t,\alpha,d}=\nu_t\quad\text{for each }t,\qquad
 \sum_{t,\alpha,d}d\,n_{t,\alpha,d}=D,\\
 \sum_{t,d,\alpha:\operatorname{Re}\alpha=1/2+\delta}
                                 n_{t,\alpha,d}=a,\qquad
 \sum_{t,d,\alpha:\operatorname{Im}\alpha=+\gamma}
                                 n_{t,\alpha,d}=b.
 \end{gathered}
\]
For these indices put
\[
 c_{\nu;a,b,D}=
 \sum_{\text{stated indices}}
   \left(\prod_{t=1}^{N-1}
             \frac{\nu_t!}{\prod_{\alpha,d}n_{t,\alpha,d}!}\right)
        \prod_{t,\alpha,d}f_{t;\alpha,d}^{\,n_{t,\alpha,d}}.
 \tag{OCS15}
\]
Empty index sets give zero and zero exponents give factors one.
Every original primary coefficient is present in this finite rule.
The ordinary multinomial expansion of OCS14 now proves
\[
 x(w)^\nu=\sum_{a,b=0}^k e^{\omega_{a,b}w}
                   \sum_{D=0}^{e-1}c_{\nu;a,b,D}w^D,
 \qquad
 \omega_{a,b}=(\beta_{a,b}-c)/i
                 =(2b-k)\gamma-i(2a-k)\delta.
 \tag{OCS16}
\]
Indeed $D\leq k(m-1)=e-1$ and the counts of each original
sign give exactly the displayed exponent. The sums with different
$(a,b)$ are distinct since their real and imaginary parts are
independently distinct.

\subsection{The actual finite observation matrix and its complete kernel}
Define the primary coefficient isomorphism
$\mathcal H:\mathbb C^q\to E$ by sending coordinate $(a,b,D)$
to $\varepsilon_{a,b}(S-\beta_{a,b})^D$, where
$\varepsilon_{a,b}$ is obtained from $\chi$ by the same truncated
Taylor-inverse construction as OCS12, now of order $e$.
The proof by the coprime factors is identical and shows that this
is a basis of all $q$ primary directions. Set
\[
 \mathfrak A_{\nu;(a,b,D)}=i^D D!\,c_{\nu;a,b,D},
 \qquad
 \mathfrak A:\mathbb C^q\longrightarrow
                            \mathbb C^{\mathcal I_{k,0}}.
 \tag{OCS17}
\]
Then, on the original spaces and with the original ordered target,
\[
 \boxed{\ L\mathcal H=\mathcal F_0\mathfrak A,\qquad
 K=\mathcal H(\ker\mathfrak A),\qquad
 \operatorname{rank}\Lambda=\operatorname{rank}\mathfrak A. }
 \tag{OCS18}
\]
Here is a proof giving every factorial and phase. In the primary
sum algebra,
\[
 e^{wY}[1]=\sum_{a,b}e^{\omega_{a,b}w}
          \sum_{D=0}^{e-1}\frac{i^{-D}w^D}{D!}
                  \varepsilon_{a,b}(S-\beta_{a,b})^D.
\]
The exact original tensor identity of OPG8 gives
$\mathcal V e^{wY}[1]=\mathbf v(w)^{\otimes k}$: the centre
$c=k/2$ splits into the $k$ original centres $1/2$, while
$U$ commutes with multiplication by $z_1$.
Applying the literal projector and then OCS9 and OCS16 gives
the right-hand curve with coefficients
$\mathcal F_0c_{a,b,D}$.
Exponential polynomials $w^De^{\omega_{a,b}w}$ with these distinct
exponents are linearly independent. To prove this directly,
apply $\prod_{(a',b')\ne(a,b)}(\partial_w-\omega_{a',b'})^e$
to any zero linear combination. It kills all the other exponents.
On the selected polynomial factor, each remaining operator is
$\partial_w+(\omega_{a,b}-\omega_{a',b'})$, invertible on the
polynomials of degree less than $e$ because its constant term
is nonzero and differentiation is nilpotent there. The surviving
polynomial must therefore vanish. Doing this for every pair proves
the asserted independence. Equating coefficients now gives
$L\mathcal H e_{a,b,D}=\mathcal F_0(i^DD!c_{a,b,D})$.
This proves the first equality in OCS18 on a basis. The other
two follow from injectivity of $\mathcal F_0$ and $\jmath$.

In particular the isomorphism from the actual sector image to
the original boundary is
\[
 \mathcal Q:\operatorname{im}\mathfrak A\xrightarrow{\ \sim\ }B,
 \qquad \mathcal Qx=\jmath^{-1}\mathcal F_0x,
 \qquad
 E/K\xrightarrow{\ \sim\ }\operatorname{im}\mathfrak A,
 \quad[z]\longmapsto\mathfrak A\mathcal H^{-1}z.
 \tag{OCS19}
\]
Both inverses exist on precisely their displayed images by OCS18.
These are full coefficient maps; the inherited Gram of
$\mathcal F_0$ is $J_0^*\mathcal GJ_0$ from OCS11.

For a completely explicit finite rank prescription, define
\[
 \Delta_r(\mathfrak A)=
   \sum_{\substack{I\subseteq\mathcal I_{k,0},\ J\subseteq\{(a,b,D)\}\\
                   |I|=|J|=r}}|\det\mathfrak A_{I,J}|^2,
 \qquad\Delta_0=1.
 \tag{OCS20}
\]
Every entry of every minor is the finite expression OCS12--17
in the actual periods and amplitude derivatives. Gaussian
elimination proves that a matrix has rank at least $r$ exactly
when one of its $r$ by $r$ minors is nonzero: pivot rows and
columns supply such a minor, and a nonzero minor supplies $r$
independent columns. Hence the actual rank is exactly
$\max\{r:\Delta_r(\mathfrak A)>0\}$. This evaluates its
finite deciding data without prescribing generic values or
asserting that any unexamined maximal minor is nonzero.

The Gamma vectors used in the existing endpoint increments have
the equally explicit coordinates
\[
 \jmath\lambda_d=\mathcal F_0\widehat\lambda_d,
 \qquad
 (\widehat\lambda_d)_\nu=
 \frac1{\sqrt{h_d^{(s)}}}\sum_{a,b=0}^k\sum_{D=0}^{e-1}
               c_{\nu;a,b,D}p_d^{(D)}(\omega_{a,b}),
 \quad h_d^{(s)}=(2\pi)^{s/2}d!(s/2)_d,
 \quad s\in\{1,k\}.
 \tag{OCS21}
\]
Taylor expansion of $p_d(Y)[1]$ in the sum primary basis gives
the coefficient $i^{-D}p_d^{(D)}(\omega_{a,b})/D!$.
Multiplication by OCS17 cancels exactly these factorials and
phases, proving OCS21. This also proves that the generating
formula of OPG9 becomes
$(1+z^2)^{-s/4}\mathcal F_0
 (x(\arctan z)^\nu)_{\nu\in\mathcal I_{k,0}}$,
with the original scalar exponent and both existing sources.

\subsection{Exact invariant dimensions and unconditional forced losses}
Put $d_{N,k,\tau}=|\mathcal I_{k,\tau}|$.
By OCS7--10 this is the dimension of the full symmetric sector,
and its exact finite character formula is
\[
 d_{N,k,0}=[T^k]\frac1N
       \sum_{h\mid N}\varphi(h)\frac{1-T}{(1-T^h)^{N/h}}.
 \tag{OCS22}
\]
To prove it, the trace generating series for $C^a$ on symmetric
tensors is
$\prod_{t=1}^{N-1}(1-T\zeta^{at})^{-1}$, by summing the
monomial basis OCS8. If $\zeta^a$ has order $h$, the list
$\zeta^{at}$ with $0\leq t<N$ contains every $h$th root
$N/h$ times. Their product is
$\prod_{t=0}^{N-1}(1-T\zeta^{at})=(1-T^h)^{N/h}$.
Removing the $t=0$ factor gives the series
$(1-T)/(1-T^h)^{N/h}$. There are exactly $\varphi(h)$
indices $a$ of that order, so averaging the traces of the
original finite projector proves OCS22. The inverse powers
give the same average by $a\mapsto-a$.
In ordinary finite binomial terms, for $k\geq1$ this is
\[
 d_{N,k,0}=\frac1N\sum_{h\mid N}\varphi(h)
 \left\{
 \mathbf1_{h\mid k}\binom{k/h+N/h-1}{N/h-1}
 -\mathbf1_{h\mid(k-1)}
                  \binom{(k-1)/h+N/h-1}{N/h-1}
 \right\}.
 \tag{OCS23}
\]
The binomial expansion of $(1-T^h)^{-N/h}$ follows by
counting weak compositions, proving the coefficient extraction
used here. It follows from the actual matrix OCS17 that
\[
 \operatorname{rank}\Lambda\leq\min(q,d_{N,k,0}),
 \qquad \dim K\geq\max(0,q-d_{N,k,0}).
 \tag{OCS24}
\]
These inequalities apply to the original map, since its full
image factorization was proved in OCS18; the ambient dimension
has not been equated with its actual rank.
There is also an unconditional nonzero observation for the
original $k\geq5$. Each entire function $x_t(w)$ in OCS14 is
nonzero. Otherwise its nonzero coordinate covector would
annihilate every derivative at zero of
$\Pi U\exp(w(A_h-1/2)/i)[1]$. Those derivatives span the full
space, since $[1]$ is cyclic for multiplication by $z_1$ and
$\Pi U$ is invertible. A product of finitely many nonzero entire
functions is nonzero. The zero sector is nonempty: take the
weights $1,1,N-2$ and then $(k-3)/2$ pairs $1,N-1$.
Their sum is a multiple of $N$. Its monomial curve in OCS9
is nonzero, so $\operatorname{rank}\Lambda\geq1$.

For $m=1$, $N=5$, and every $k\geq0$, write
\[
\begin{gathered}
 D_k=\binom{k+3}{3},\qquad
 \epsilon_k=\begin{cases}1&k\equiv0\pmod5,\\-1&k\equiv1\pmod5,\\
                         0&\text{otherwise}.
             \end{cases}\\
 d_{5,k,0}=\frac{D_k+4\epsilon_k}{5},\quad
 d_{5,k,\tau}=\frac{D_k-\epsilon_k}{5}\quad(\tau\ne0).
\end{gathered}\tag{OCS25}
\]
For each of the four nonidentity powers of $C$, the trace
series is $(1-T)/(1-T^5)$, whose coefficient is $\epsilon_k$.
The weight-$\tau$ projector introduces the Fourier multiplier;
the sum of its four nontrivial fifth roots is $4$ when
$\tau=0$ and $-1$ otherwise. This proves both formulas.
For the requested original orders the exact counts and resulting
kernel lower bounds are
\[
\begin{array}{c|r|r|r|r}
 k&D_k&d_{5,k,0}&q=(k+1)^2&\max(0,q-d_{5,k,0})\\\hline
 5&56&12&36&24\\
 9&220&44&100&56\\
 13&560&112&196&84\\
 17&1140&228&324&96\\
 21&2024&404&484&80\\
 25&3276&656&676&20\\
 29&4960&992&900&0
\end{array}\tag{OCS26}
\]
These are integer evaluations of OCS25. The strict inequality
$d_{5,k,0}<q$ holds for $k=4l+1$, $l\geq1$, exactly for the
first six orders in the table. They exhaust $k\leq25$ in this
progression. For $k\geq29$,
\[
 D_k-5(k+1)^2=\frac{(k+1)(k^2-25k-24)}6\geq460>4,
\]
where $k^2-25k-24$ is positive and increasing from $k=29$.
Even the most negative correction $4\epsilon_k\geq-4$
therefore gives $d_{5,k,0}>q$. This proves the threshold
without extrapolating the finite table.

\subsection{The actual simple-quartet quadric and a stronger rank bound}
Continue with $m=1$. Order the four roots as
$(++),(+-),(-+),(--)$ and let
$\mathcal E_{\rm prim}:\mathbb C^4\to E_h$ be the primary
basis isomorphism whose columns are the four $[E_{\alpha,0}]$.
Define the actual invertible matrix and the original exponential
quartet by
\[
 \mathsf H=F^{-1}\Pi U\mathcal E_{\rm prim},\qquad
 t_{\epsilon,\eta}(w)=e^{(\eta\gamma-i\epsilon\delta)w},
 \qquad x(w)=\mathsf Ht(w).
 \tag{OCS27}
\]
Every column of $\mathsf H$ is OCS13 with $d=0$, namely the
Fourier transform of its original primary period times its
actual nonzero $v_h(\alpha)$. Each displayed map is invertible,
so $\mathsf H$ is invertible without a genericity assumption.
Put $\mathcal R_k=\mathbb C[x_1,x_2,x_3,x_4]_k$, the homogeneous
polynomials of degree $k$, and define
\[
 Q_0(t)=t_{++}t_{--}-t_{+-}t_{-+},\qquad
 Q(x)=Q_0(\mathsf H^{-1}x).
 \tag{OCS28}
\]
This is the literal transformed quadratic relation of the
original curve; its coefficients are determined by the periods
and unit through OCS27. In particular $Q\ne0$.

The kernel of homogeneous restriction to that curve is exactly
\[
 \ker\{\mathcal R_k\to\operatorname{Hol}(\mathbb C),
                         f\mapsto f(x(w))\}
            =Q\mathcal R_{k-2}\quad(k\geq2).
 \tag{OCS29}
\]
Here is a complete verification. First work in the $t$ coordinates.
A degree-$k$ monomial has exponent
$(2b-k)\gamma-i(2a-k)\delta$, where $a$ is its count of
first sign $+$ and $b$ its count of second sign $+$.
These $(k+1)^2$ complex exponents are distinct: equality of
their real parts forces $b$ equal and equality of their
imaginary parts forces $a$ equal. Equivalently, the two
frequencies $2\gamma$ and $-2i\delta$ admit no nonzero
integer relation, by their real and imaginary parts. Thus
the exponentials have no extra polynomial relation of this
homogeneous type; their linear independence follows, for
example, by differentiating $0,\ldots,(k+1)^2-1$ times and
using the Vandermonde determinant of the distinct exponents.
Every pair $(a,b)$ occurs: the count $r_{++}$ may be any integer
between $\max(0,a+b-k)$ and $\min(a,b)$, a nonempty interval,
and the other three counts are $a-r_{++}$, $b-r_{++}$,
and $k-a-b+r_{++}$.

Within any fixed pair $(a,b)$, two monomial exponent vectors
differ by an integer multiple of $(1,-1,-1,1)$. After their
common monomial is factored out, their difference is therefore
the difference of equal powers of $t_{++}t_{--}$ and
$t_{+-}t_{-+}$. The identity
$A^r-B^r=(A-B)\sum_{j=0}^{r-1}A^{r-1-j}B^j$ shows that
this difference is divisible by $Q_0$. Any polynomial
vanishing on the curve has coefficient sum zero within
each pair $(a,b)$, by the proved exponential independence.
It is thus a sum of these monomial differences and is
divisible by $Q_0$. Its quotient is homogeneous of degree
$k-2$. Conversely $Q_0(t(w))=0$, since each of its two
products is one. This proves the assertion in $t$
coordinates, and the invertible substitution OCS27 proves
OCS29. In particular
$\dim(\mathcal R_k/Q\mathcal R_{k-2})=(k+1)^2=q$;
no information about the original curve has been replaced
by a freely chosen quadric.

Let $\mathcal I_k\subseteq\mathcal R_k$ be the span of the
weight-zero monomials $x^\nu$, using the weights $1,2,3,4$
from the exact Fourier basis. The coordinate functions of
the original projected curve in OCS9 are precisely these
monomials restricted to $x(w)$. The span of the vector
curve $Le^{wY}[1]$ equals $\operatorname{im}L$: its
derivatives at zero contain $LY^d[1]$, $0\leq d<q$, and
$Y^d[1]$ is a basis because $Y$ is an invertible affine
change of multiplication by $S$ in the degree filtration.
The rank of that vector curve equals the dimension of
the span of its coordinate functions, by the equality of
row and column ranks of its finite exponential coefficient
matrix. Hence OCS29 proves the exact equality
\[
 \operatorname{rank}\Lambda
   =d_{5,k,0}-\dim(\mathcal I_k\cap Q\mathcal R_{k-2}).
 \tag{OCS30}
\]

Fix any multiplicative monomial order, for instance
lexicographic order, and write $x^\mu$ for the leading
monomial of the actual $Q$, with $|\mu|=2$. Choose a
row-echelon basis of the finite-dimensional intersection
in OCS30 in descending monomial order. Its leading
monomials are distinct, all have weight zero, and all
are divisible by $x^\mu$: the leading monomial of
$Qf$ is the product of those of $Q$ and $f$ because
the order is multiplicative and the coefficient field
has no zero divisors. Thus each such leading monomial
is $x^\mu x^\eta$ for a distinct degree-$k-2$ monomial
of weight $-\operatorname{wt}(\mu)$. Counting this set
and using OCS30 gives
\[
 \boxed{\quad
 d_{5,k,0}-d_{5,k-2,-\operatorname{wt}(\mu)}
 \ \leq\ \operatorname{rank}\Lambda
 \ \leq\ \min(q,d_{5,k,0}).\quad}
 \tag{OCS31}
\]
The left side depends only on the actual leading monomial
of the actual period quadric, and remains valid whatever
its other coefficients are. A uniform consequence retaining
the strongest count independent of that monomial is
\[
 r_k^-:=d_{5,k,0}-\max_{\tau\in\mathbb Z/5}d_{5,k-2,\tau}
       =\left\lfloor\frac{(k+1)^2}{5}\right\rfloor
                                  +\mathbf1_{5\mid k}
       \leq\operatorname{rank}\Lambda.
 \tag{OCS32}
\]
For the equality, $D_k-D_{k-2}=(k+1)^2$, as direct
subtraction of the two cubic binomials shows. Formula
OCS25 then gives
\[
 r_k^-=\frac{(k+1)^2+4\epsilon_k
                   -\max(4\epsilon_{k-2},-\epsilon_{k-2})}{5}.
\]
For $k$ congruent respectively to $0,1,2,3,4$ modulo five,
its numerator correction is respectively $4,-4,-4,-1,0$.
The corresponding residues of $(k+1)^2$ are $1,4,4,1,0$.
These five cases prove OCS32 for every $k\geq2$.
For the original orders in OCS26 the complete unconditional
intervals are therefore
\[
\begin{array}{c|r|r|r|r}
 k&r_k^-&\min(q,d_{5,k,0})&\max(0,q-d_{5,k,0})&q-r_k^-\\\hline
 5&8&12&24&28\\
 9&20&44&56&80\\
 13&39&112&84&157\\
 17&64&228&96&260\\
 21&96&404&80&388\\
 25&136&656&20&540\\
 29&180&900&0&720
\end{array}\tag{OCS33}
\]
The middle two rank columns bound $\operatorname{rank}\Lambda$;
the last two columns bound $\dim K$ from below and above,
respectively. No endpoint of either interval has been assigned
as the actual value without its original minors.

\subsection{A second exact finite matrix for the simple-quartet kernel}
The quadric calculation also gives the actual kernel a
polynomial quotient model. Write $Q(x)=\sum_{|\mu|=2}q_\mu x^\mu$.
Let $T_Q$ be the full multiplication matrix
$\mathcal R_{k-2}\to\mathcal R_k$ in the monomial bases:
\[
 (T_Q)_{\nu,\eta}=q_{\nu-\eta},
 \quad |\nu|=k,\quad|\eta|=k-2,
 \qquad q_{\nu-\eta}=0\text{ if an entry of }\nu-\eta<0.
\]
Let $T_{\rm out}$ be its rows of nonzero cyclic weight.
Every coefficient $q_\mu$ is explicitly fixed by
$Q_0(\mathsf H^{-1}x)$, and the inverse matrix may be
computed as $\operatorname{adj}\mathsf H/\det\mathsf H$,
with nonzero denominator proved in OCS27. Multiplication
by $Q$ is injective. The condition that $Qf$ belongs
to $\mathcal I_k$ is exactly $T_{\rm out}f=0$.
Thus OCS30 gives the further exact rank identities
\[
 \begin{split}
 \operatorname{rank}\Lambda
   &=d_{5,k,0}-D_{k-2}+\operatorname{rank}T_{\rm out},\\
 \dim K&=D_k-d_{5,k,0}-\operatorname{rank}T_{\rm out}.
 \end{split}\tag{OCS34}
\]
The ranks of this concrete matrix are decided by its finite
minors, by the same elimination proof as OCS20. This is
an alternative exact calculation from the actual periods,
not an assumption about the quadric's coefficients.

To give the typed map behind the kernel equality, let
$E^*$ denote the complex linear dual. Define
$\Theta:\mathcal R_k\to E^*$ as follows: $\Theta(f)$ is
the unique covector such that
\[
 \Theta(f)(e^{wY}[1])=f(x(w))\quad\text{for all }w.
 \tag{OCS35}
\]
For $m=1$ every function on the right is a linear
combination of the $q$ distinct exponentials in OCS16.
The primary expansion of $e^{wY}[1]$ gives exactly those
exponentials as a basis of covector functions. This proves
existence and uniqueness. The proof of OCS29 proves
surjectivity and $\ker\Theta=Q\mathcal R_{k-2}$.
Moreover $\Theta(\mathcal I_k)=\operatorname{im}L^\vee$,
where $L^\vee$ denotes the complex linear dual map:
the monomial coefficient functionals in the basis OCS8
pull back along $L$ to exactly OCS35 for $x^\nu$ of
weight zero. These functionals span the full dual of
the actual image, since every linear functional on a
subspace of a finite-dimensional space extends to the
whole space. Restriction of $E^*$ to $K$ has kernel
$\operatorname{im}L^\vee$: a covector vanishing on $K$
factors uniquely through $E/K\simeq\operatorname{im}L$.
Thus the exact maps give
\[
 \boxed{\quad
 \operatorname{coker}T_{\rm out}
 \simeq\mathcal R_k/(\mathcal I_k+Q\mathcal R_{k-2})
 \xrightarrow{\ \Theta|_K\ }K^*.
 \quad}\tag{OCS36}
\]
For the first isomorphism, use the direct monomial splitting
of $\mathcal R_k$ into its zero and nonzero sectors; after
quotienting by $\mathcal I_k$, multiplication by $Q$
becomes precisely $T_{\rm out}$. For the second, the
surjectivity and kernel were just proved. This identifies
the actual lost directions by a full quotient map.

\subsection{All primary action directions remain in the mixed calculation}
For arbitrary original $m\geq1$, retain the full primary
matrix $\mathfrak A$ in OCS17. In the basis $\mathcal H$
the original sum action is exactly
\[
 J e_{a,b,D}=\omega_{a,b}e_{a,b,D}
       +\begin{cases}i^{-1}e_{a,b,D+1}&D<e-1,\\0&D=e-1,
         \end{cases}
 \qquad Y\mathcal H=\mathcal HJ.
 \tag{OCS37}
\]
This follows by multiplying
$\varepsilon_{a,b}(S-\beta_{a,b})^D$ by
$(S-c)/i=\omega_{a,b}+(S-\beta_{a,b})/i$ in its original
primary algebra. Consequently the actual possible action
leaving the cyclic kernel is the fully specified map
\[
 \ker\mathfrak A\longrightarrow\operatorname{im}\mathfrak A,
 \qquad x\longmapsto\mathfrak A Jx,
 \qquad
 \Lambda Y\mathcal Hx=\mathcal Q\mathfrak A Jx.
 \tag{OCS38}
\]
The codomain assertion follows from $\mathfrak A Jx$ being
a value of the same matrix $\mathfrak A$ on $Jx$.
For the original inclusion $I:K\hookrightarrow E$, this
is the OPG22 defect transported through the exact kernel
isomorphism in OCS18. In particular none of the discarded
sectors has been declared invariant under $Y$.

The part of the kernel undetected by all original iterates
is exactly computed by the stacked matrix
\[
 \mathfrak O=\begin{bmatrix}
          \mathfrak A\\\mathfrak AJ\\\vdots\\\mathfrak AJ^{q-1}
          \end{bmatrix},\qquad
 \mathcal K_\infty=\mathcal H\ker\mathfrak O,
 \quad
 \dim\mathcal K_\infty=q-\operatorname{rank}\mathfrak O,
 \quad
 \dim(K/\mathcal K_\infty)=
        \operatorname{rank}\mathfrak O-\operatorname{rank}\mathfrak A.
 \tag{OCS39}
\]
Indeed $\mathfrak A J^a x$ is the exact sector-coordinate
value of $\Lambda Y^a\mathcal Hx$ by OCS18 and OCS37.
The polynomial $\chi(c+iT)$ annihilates $J$ and has
degree $q$ with leading coefficient $i^q$. It expresses
every higher power in the span of the first $q$, so the
stacked kernel is precisely the intersection over all
iterates and is $J$-invariant. Rank-nullity proves both
dimension formulas. Its quotient identifies the directions
that the original one-step observation loses and subsequent
original iterates detect.

Equations OCS17--21, OCS34--39 supply actual finite maps
and deciding minors for the original observation and its
kernel. The forced kernel losses in OCS26 and the rank
intervals OCS33 concern this unchanged projector and these
unchanged periods. The endpoint volume calculation still
uses the full positive sums of the actual vectors
$\lambda_d$ in OPG16--18 and their exact inherited
coordinates OCS11 and OCS19. No volume growth estimate
is deduced from these dimension bounds alone.

% END COMPLETE OCS

\clearpage
% BEGIN COMPLETE OQC
% Derived from the literal original period quadric and cyclic projector.
\section{The original quadric orbit and an explicit invariant lifting map}
\label{sec:OQC}

This calculation uses the original simple-quartet construction, with
$m=1$, $k=4l+1$, $l\geq1$, $q=(k+1)^2$, $0<\delta<1/2$, and
$\gamma>2$. The period parameter, its branch, the polynomial contours,
the complete amplitude unit, and the observation are those of OPG and
OCS. The degree-eight polynomial below is the product of four actual
orbit equations. Its degree and all of its coefficients are derived
from that orbit; no source order is selected to match an estimate.

\subsection{The actual coefficient ring and the original observation}
Order the original roots as $(++),(+-),(-+),(--)$, with
$\alpha_{\epsilon,\eta}=1/2+\epsilon\delta+i\eta\gamma$.
Let $\mathcal E_{\rm prim}$ have columns the four original primary
idempotents in $E_h=\mathbb C[z_1]/(h)$. For the original matrices
$\Pi,U,F$ in OCS2, OCS5, OCS27 put
\[
\begin{gathered}
 \mathsf H=F^{-1}\Pi U\mathcal E_{\rm prim},\qquad
 t_{\epsilon,\eta}(w)=e^{(\eta\gamma-i\epsilon\delta)w},\\
 x(w)=\mathsf Ht(w),\qquad
 Q(x)=Q_0(\mathsf H^{-1}x),\quad
 Q_0(t)=t_{++}t_{--}-t_{+-}t_{-+}.
\end{gathered}\tag{OQC1}
\]
Every matrix in $\mathsf H$ is the specified original isomorphism;
in particular its nonzero amplitude factors are the actual
$v_h(\alpha)$, and its period entries are the original convergent
polynomial integrals. Thus $\det\mathsf H\ne0$ by the included
period-isomorphism and primary-basis proofs. Define
\[
 \mathscr P_d=\mathbb C[x_1,x_2,x_3,x_4]_d,\qquad
 R_d=\mathscr P_d/(Q\mathscr P_{d-2}),\qquad
 \mathscr P_d=R_d=0\quad(d<0).
 \tag{OQC2}
\]
The matrix of $Q_0(t)=t^TM_0t$ has entries
$(M_0)_{14}=(M_0)_{41}=1/2$ and
$(M_0)_{23}=(M_0)_{32}=-1/2$, all others zero. Hence
\[
 M=\mathsf H^{-T}M_0\mathsf H^{-1},\qquad
 \det M=\frac{1}{16(\det\mathsf H)^2}\ne0.
 \tag{OQC3}
\]
In particular $Q$ is irreducible. Indeed, a reducible homogeneous
quadratic over $\mathbb C$ is a product of two linear forms, whose
symmetric coefficient matrix is a sum of two rank-one matrices and
has rank at most two. This contradicts OQC3. The polynomial ring
over a field is a unique factorization domain, by the repeated
one-variable Gauss lemma, so its irreducible element $Q$ is prime.
Consequently $R=\mathbb C[x_1,x_2,x_3,x_4]/(Q)$ is an integral domain.

For clarity, the complete homogeneous restriction calculation has
the precise kernel $Q\mathscr P_{d-2}$. In the $t$ coordinates a
degree-$d$ monomial restricts to
\[
 \exp\bigl(((2b-d)\gamma-i(2a-d)\delta)w\bigr),
       \qquad 0\leq a,b\leq d.
 \tag{OQC4}
\]
Every pair occurs: the number of $(++)$ factors can range from
$\max(0,a+b-d)$ to $\min(a,b)$. The exponents are distinct by
their real and imaginary parts. Their exponentials are linearly
independent, since differentiating orders zero through $(d+1)^2-1$
gives a Vandermonde matrix with nonzero determinant. Two monomials
giving the same pair have exponent vectors differing by a multiple
of $(1,-1,-1,1)$. After their common monomial is removed, their
difference is divisible by $t_{++}t_{--}-t_{+-}t_{-+}$ through
$A^r-B^r=(A-B)\sum_{j=0}^{r-1}A^{r-1-j}B^j$. A vanishing
polynomial has coefficient sum zero within each pair by the
independence just proved, and is therefore a sum of these differences.
Conversely $Q_0(t(w))=0$. Transport by $\mathsf H$ proves the
stated kernel and
\[
                  \dim R_d=(d+1)^2\quad(d\geq0).
 \tag{OQC5}
\]

Write $\zeta=e^{2\pi i/5}$ and
$D=\operatorname{diag}(\zeta,\zeta^2,\zeta^3,\zeta^4)$.
The literal Fourier calculation gives $CF=FD$. On polynomials
define the original induced action and its average by
\[
 (g f)(x)=f(D^{-1}x),\qquad g^5=1,\qquad
 \mathsf R f=\frac15\sum_{a=0}^4g^af,
 \qquad \mathscr I_d=\mathscr P_d^{\langle g\rangle}.
 \tag{OQC6}
\]
Thus $\mathscr I_d$ consists of the weight-zero monomials, since
$gx^\nu=\zeta^{-\sum j\nu_j}x^\nu$. This is the dual of the
original average of $(C^{-a})^{\otimes d}$; both the inverse powers
and the factor $1/5$ are retained.

For the original algebra $E_d$ of the $d$-fold sum, define
$\Theta_d:R_d\xrightarrow{\sim}E_d^\vee$ by
\[
 \Theta_d([f])(e^{wY_d}[1])=f(x(w)),\qquad
       Y_d=(A_d-dI/2)/i.
 \tag{OQC7}
\]
Here $\vee$ denotes the complex-linear dual. The primary expansion
of the vector on the left gives all distinct exponentials OQC4;
therefore existence and uniqueness of this covector and bijectivity
follow from the preceding complete restriction calculation. The
original observation $\Lambda_d:E_d\to B_d$ obeys
\[
 J_d:=\operatorname{im}(\mathscr I_d\to R_d),\qquad
 \Theta_d(J_d)=\operatorname{im}\Lambda_d^\vee,
 \qquad \dim J_d=\operatorname{rank}\Lambda_d,
 \qquad R_d/J_d\xrightarrow{\sim}(\ker\Lambda_d)^\vee.
 \tag{OQC8}
\]
Indeed the coordinate functions of the projected pure tensor in
OCS9 are precisely $x(w)^\nu$ of weight zero; their pulled-back
covectors span the dual of its actual image. A covector vanishing
on $\ker\Lambda_d$ factors through that image, which proves the
last isomorphism as well. OQC8 specifies the exact original image;
it has not been replaced by the whole invariant tensor space.

\subsection{The complete orbit decision from the actual coefficients}
Put $Q_a=g^aQ$ for $0\leq a<5$. The orbit of the line $\mathbb C Q$
has size one or five because five is prime. Moreover a size-one
orbit has $gQ=Q$, with no other possible scalar. To prove this, write
$gQ=\lambda Q$. Then $\lambda^5=1$, while OQC3 and the determinant
of $D$, which is $\zeta^{10}=1$, imply
\[
 \det(D^{-T}MD^{-1})=\det M=\lambda^4\det M.
 \tag{OQC9}
\]
Thus $\lambda^4=1$ and $\lambda^5=1$, forcing $\lambda=1$.
If any nonidentity power fixed the line, its powers generate the
same cyclic group, so the preceding argument applies. This proves
the entire orbit classification.

There is an exact coefficient test, requiring no genericity premise.
Write the actual $Q=\sum_{|\mu|=2}q_\mu x^\mu$ and put
\[
 \mathfrak e(Q)=\sum_{\substack{|\mu|=2\\
                      \sum j\mu_j\not\equiv0\ (5)}}|q_\mu|^2.
 \tag{OQC10}
\]
All coefficients here are the entries of OQC3, including the full
inverse of the original period matrix. Since the monomials are a
basis and each summand is nonnegative, $\mathfrak e(Q)=0$ exactly
when $Q$ is invariant. The only degree-two weight-zero monomials
are $x_1x_4$ and $x_2x_3$. Hence the size-one case is exactly
$Q=a x_1x_4+b x_2x_3$ with $ab\ne0$; nonvanishing of the two
coefficients follows from OQC3. Positive $\mathfrak e(Q)$ gives
the five distinct irreducible orbit equations. These are exhaustive
outcomes of a test on the given periods, not assumptions on them.

In the size-one outcome, multiplication by $Q$ preserves every
weight. If $Qf$ is invariant, then $Q(gf-f)=0$ in the polynomial
ring, and $gf=f$. Therefore
\[
 \ker(\mathscr I_k\to R_k)=Q\mathscr I_{k-2},\qquad
 \operatorname{rank}\Lambda_k=d_{5,k,0}-d_{5,k-2,0},
 \quad \dim\ker\Lambda_k=q-d_{5,k,0}+d_{5,k-2,0}.
 \tag{OQC11}
\]
This also proves that $J_k=R_k^{\langle g\rangle}$: an invariant
class can be lifted to a polynomial, and averaging that lift gives
an invariant representative. The count itself follows from
\[
 d_{5,k,0}=\frac{\binom{k+3}{3}+4\epsilon_k}{5},\qquad
 \epsilon_k=\mathbf1_{k\equiv0\ (5)}-\mathbf1_{k\equiv1\ (5)}.
 \tag{OQC12}
\]
For a proof, the identity element has symmetric-power generating
series $(1-z)^{-4}$; each of the other four powers has series
$\prod_{j=1}^4(1-z\zeta^j)^{-1}=(1-z)/(1-z^5)$.
Averaging their coefficients gives OQC12. Subtraction and
$\binom{k+3}{3}-\binom{k+1}{3}=(k+1)^2$ evaluate the rank in
OQC11, for every $k\geq2$, as
\[
 \left\lfloor\frac{(k+1)^2}{5}\right\rfloor
                 +\mathbf1_{k\equiv0\text{ or }3\ (5)}.
 \tag{OQC13}
\]
The five numerator corrections before division by five are
$4,-4,-4,4,0$, respectively; the residues of $(k+1)^2$ are
$1,4,4,1,0$. This verifies every case in OQC13.

\subsection{An explicit invariant lift for the five-member orbit}
For the size-five outcome determined by OQC10, retain the full
actual products
\[
 \mathcal H_Q=\prod_{a=0}^4Q_a\quad(\deg\mathcal H_Q=10),
 \qquad \mathcal A_Q=\prod_{a=1}^4Q_a
                                      \quad(\deg\mathcal A_Q=8).
 \tag{OQC14}
\]
The action permutes all five factors cyclically, so
$g\mathcal H_Q=\mathcal H_Q$. None of the four factors of
$\mathcal A_Q$ is a scalar multiple of $Q$. Since each is
irreducible and $Q$ is prime, multiplication by $[\mathcal A_Q]$
is injective on the integral domain $R$.

For every homogeneous $f\in\mathscr P_d$ the exact trace identity is
\[
 \left[5\mathsf R(\mathcal A_Q f)\right]_{(Q)}
                          =[\mathcal A_Q f]_{(Q)}.
 \tag{OQC15}
\]
Indeed $g^a\mathcal A_Q$ is the product of all orbit factors
except $Q_a$. For $a\ne0$ it contains $Q_0=Q$. Those four
terms vanish in the displayed quotient, and the $a=0$ term is
exactly the right side. This proves OQC15 with its original
average and full factors, including every period-dependent
coefficient.

The lift also has a well-defined typed quotient map. Let
$\mathscr U=\mathbb C[x_1,x_2,x_3,x_4]/(\mathcal H_Q)$;
the invariant equation $\mathcal H_Q$ makes the cyclic action
well defined on this quotient. Then
\[
 \begin{gathered}
 \mathfrak L_d:R_d\longrightarrow
                      (\mathscr U_{d+8})^{\langle g\rangle},\qquad
 [f]_{(Q)}\longmapsto[5\mathsf R(\mathcal A_Q f)]_{(\mathcal H_Q)},\\
 \operatorname{res}\mathfrak L_d([f])=[\mathcal A_Q f]_{(Q)},
 \qquad \operatorname{res}:\mathscr U\to R.
 \end{gathered}\tag{OQC16}
\]
If $f$ changes by $Qb$, its proposed image changes by
$5\mathsf R(\mathcal H_Qb)=5\mathcal H_Q\mathsf Rb$, hence
by zero in $\mathscr U$. This proves well-definedness. The
image is invariant because it has an invariant polynomial
representative; OQC15 proves the second line. Injectivity follows
from the injectivity of multiplication by $[\mathcal A_Q]$ on
$R$. All maps in OQC16 are thus established, including their kernels.

The actual original invariant image consequently contains
\[
 [\mathcal A_Q]R_{k-8}\ \subseteq\ J_k\ \subseteq\ R_k.
 \tag{OQC17}
\]
The first inclusion follows by taking the invariant polynomial
representatives OQC15, and uses no dimension inference. Therefore,
for every $k\geq8$ in this outcome,
\[
 \operatorname{rank}\Lambda_k\geq(k-7)^2,\qquad
 \dim\ker\Lambda_k\leq(k+1)^2-(k-7)^2=16k-48.
 \tag{OQC18}
\]
These inequalities follow from the proved injection and OQC5.
They may be combined with the unconditional actual-matrix bounds
in OCS24, OCS31--33 by taking the stronger lower and upper bounds.
For the original degrees $k=9,13,17,21,25,29$, the displayed upper
bounds on the original kernel are respectively
$96,160,224,288,352,416$; at the first two degrees the earlier
unconditional bounds may be stronger. No such upper bound is
assigned to the size-one outcome, whose exact rank is OQC11.
Together OQC10--18 compute the complete orbit alternatives from
the original coefficients. Evaluating $\mathfrak e(Q)$ in the
original period family is a separate calculation on those explicit
coefficients; no outcome is postulated here.

\subsection{The induced original dual map and its complete action defect}
At $k\geq9$ the degree $d=k-8$ is positive, and the original
$d$-fold sum algebra is defined using the same $h,\Pi,U$.
For the size-five outcome, OQC7, OQC8 and OQC17 give the uniquely
typed injective map
\[
 \begin{gathered}
 \mathfrak J_k:E_{k-8}^\vee\hookrightarrow B_k^\vee,\\
 \Lambda_k^\vee\mathfrak J_k
        =\Theta_k\,M_{[\mathcal A_Q]}\,\Theta_{k-8}^{-1},\qquad
 \mathfrak J_k^\vee:B_k\twoheadrightarrow E_{k-8}.
 \end{gathered}\tag{OQC19}
\]
The right side of the second line has image in
$\operatorname{im}\Lambda_k^\vee$ by OQC17, and
$\Lambda_k^\vee$ is injective because $\Lambda_k$ is onto its
actual image. This gives existence and uniqueness. Each of the
three right-hand maps is injective, so $\mathfrak J_k$ is
injective. Finite-dimensional duality then proves the stated
surjection, using the canonical bidual identifications. The new
degree $k-8$ records contraction against the degree-eight
polynomial derived in OQC14. It does not replace either original
Gamma source of order $1$ or $k$.

The surviving action and its defect are fully determined. Put
\[
 \Omega=\operatorname{diag}(\gamma-i\delta,-\gamma-i\delta,
                              \gamma+i\delta,-\gamma+i\delta),
 \quad G=\mathsf H\Omega\mathsf H^{-1},\qquad
 (\mathcal Df)(x)=\sum_{j=1}^4(Gx)_j\partial_jf(x).
 \tag{OQC20}
\]
Here $x'(w)=Gx(w)$ by OQC1. In the original $t$ coordinates,
the weights of each product in $Q_0$ sum to zero, so
$\mathcal DQ=0$. Thus this derivation acts on each $R_d$.
Differentiating OQC7 proves exactly
\[
          \Theta_d\mathcal D=Y_d^\vee\Theta_d.
 \tag{OQC21}
\]
Both sides evaluated at $e^{wY_d}[1]$ are the derivative of
$f(x(w))$, and these vectors span $E_d$ by the distinct
exponentials; hence equality holds as maps, not only on a curve.
Leibniz's rule gives the complete defect
\[
 \begin{split}
 Y_k^\vee\Lambda_k^\vee\mathfrak J_k
       -\Lambda_k^\vee\mathfrak J_kY_{k-8}^\vee
 &=\Theta_k M_{[\mathcal D\mathcal A_Q]}\Theta_{k-8}^{-1},\\
 \mathcal D\mathcal A_Q
 &=\sum_{a=1}^4(\mathcal DQ_a)
                  \prod_{\substack{1\leq b\leq4\\b\ne a}}Q_b.
 \end{split}\tag{OQC22}
\]
The codomain in the first line is $E_k^\vee$. Its image has not
been asserted to lie in $\operatorname{im}\Lambda_k^\vee$.
Restriction to $\ker\Lambda_k$ supplies the exact remaining
map into $(\ker\Lambda_k)^\vee$, through OQC8. In particular
OQC19 is not used as an unproved intertwining of boundary actions.
OQC22 retains every term that can leave the observed dual subspace.

Finally, the failure of the cyclic average to commute with the
original action is explicitly calculable before any quotient:
\[
 ([\mathcal D,g^a]f)(x)
   =(\nabla f)(D^{-a}x)^T
                   (D^{-a}G-GD^{-a})x,
 \qquad
 [\mathcal D,\mathsf R]=\frac15\sum_{a=0}^4[\mathcal D,g^a].
 \tag{OQC23}
\]
The chain rule applied separately to $f(D^{-a}x)$ and to
$(\mathcal Df)(D^{-a}x)$ proves both identities. Thus all
action comparisons in the invariant lift have explicit original
coefficient maps. The rank estimates in OQC18 concern exactly
the kernel used in the original Gamma and arithmetic quotient
Grams, because those constructions use the same $\Lambda_k$.
An estimate for its volume still requires those original metrics;
no volume conclusion has been inferred solely from its dimension.

% END COMPLETE OQC

\clearpage
% BEGIN COMPLETE OQX
% Exact completion of the orbit rank calculation. No period outcome is assumed.
\section{Exact original rank from the complete quadric orbit}
\label{sec:OQX}

Retain every original object and map in OQC1--23. In particular
$m=1$, $k=4l+1$, $l\geq1$, $q=(k+1)^2$,
$\mathscr P=\mathbb C[x_1,x_2,x_3,x_4]$, and
$g f(x)=f(D^{-1}x)$ for the original Fourier matrix
$D=\operatorname{diag}(\zeta,\zeta^2,\zeta^3,\zeta^4)$.
The polynomial $Q=Q_0(\mathsf H^{-1}x)$ has its actual original
period and amplitude coefficients. Write
$\mathscr I=\mathscr P^{\langle g\rangle}$,
$R=\mathscr P/(Q)$ and $J=\operatorname{im}(\mathscr I\to R)$,
with the grading and the actual observation identification of OQC8.
For negative indices put $\mathscr I_d=0$ and $d_{5,d,0}=0$.

The orbit decision is the exhaustive test OQC10. The size-one
case already has the exact rank OQC11--13. The following calculation
completes the other case using the same original coefficients. It
does not assert that a period parameter belongs to that case before
its coefficient test has been evaluated.

\subsection{Every invariant vanishing polynomial contains the full orbit}

In the size-five case, let $Q_a=g^aQ$ and retain exactly
$\mathcal H_Q=\prod_{a=0}^4Q_a$ from OQC14. Each $Q_a$ is
irreducible, because an invertible linear substitution preserves
factorizations and $Q$ is irreducible by OQC3. The five factors are
pairwise nonassociate by the proved orbit classification, and each
is prime in the original polynomial ring. Consequently
\[
 \boxed{\quad
 \mathscr I\cap Q\mathscr P=\mathcal H_Q\mathscr I,
 \qquad
 \ker(\mathscr I_d\longrightarrow R_d)
                       =\mathcal H_Q\mathscr I_{d-10}.
 \quad}\tag{OQX1}
\]
Here is the complete divisibility proof. If $f\in\mathscr I$ and
$f=Qb$, then for every $a=0,\ldots,4$ the exact equality
$g^af=f$ gives $f=Q_a g^ab$. Thus every one of the five original
prime factors divides $f$. Distinctness of the factors implies
that their product divides $f$: after factoring out $Q_0$,
primality of $Q_1$ and the fact that $Q_1$ does not divide $Q_0$
force $Q_1$ to divide the remaining quotient. Repeating this
argument with $Q_2,Q_3,Q_4$ proves
$f=\mathcal H_Q b_0$. The action permutes the five exact factors
cyclically, so $g\mathcal H_Q=\mathcal H_Q$. Applying $g$ to
this last equality and using $gf=f$ gives
$\mathcal H_Q(gb_0-b_0)=0$. The polynomial ring is an integral
domain and $\mathcal H_Q\ne0$, whence $gb_0=b_0$.
This proves the first inclusion in OQX1. Conversely an invariant
$b_0$ gives an invariant $\mathcal H_Qb_0$ divisible by $Q$,
which proves the reverse inclusion. The product is homogeneous
of degree ten. If its product with $b_0$ is homogeneous of degree
$d$, decomposition of $b_0$ into homogeneous parts shows that
only its part of degree $d-10$ can be nonzero: each other part
has zero product with $\mathcal H_Q$ and hence is zero.
This proves the graded kernel formula, including $d<10$.

\subsection{The exact invariant quotient and the restriction map}

Let $\mathscr U=\mathscr P/(\mathcal H_Q)$ as in OQC16. The
following are isomorphisms of graded complex algebras with their
specified maps:
\[
 \boxed{\quad
 \mathscr I/(\mathcal H_Q\mathscr I)
 \xrightarrow{\ [f]\mapsto[f]_{(\mathcal H_Q)}\ }
       \mathscr U^{\langle g\rangle}
 \xrightarrow{\ \operatorname{res}\ } J\subseteq R.
 \quad}\tag{OQX2}
\]
The action on $\mathscr U$ is defined because $\mathcal H_Q$ is
invariant. If $[f]_{(\mathcal H_Q)}$ is invariant, averaging its
representatives gives
$[\mathsf R f]_{(\mathcal H_Q)}=[f]_{(\mathcal H_Q)}$, since
$[g^af]=[f]$ for every $a$. Thus the first map is onto.
Its kernel before passage to the displayed domain is
$\mathscr I\cap\mathcal H_Q\mathscr P
=\mathcal H_Q\mathscr I$: the reverse inclusion is immediate,
and for $f=\mathcal H_Qb\in\mathscr I$, cancellation in
$\mathcal H_Q(gb-b)=0$ proves $b\in\mathscr I$.
This proves that the first map is injective as well.

The full quotient restriction
$\operatorname{res}:\mathscr U\to R$ is well defined because
$(\mathcal H_Q)\subseteq(Q)$; its kernel is exactly
$Q\mathscr P/\mathcal H_Q\mathscr P$.
For its invariant subspace, choose the invariant representative
just constructed. If its restriction vanishes, that representative
belongs to $\mathscr I\cap Q\mathscr P$, which is
$\mathcal H_Q\mathscr I$ by OQX1. Its class in $\mathscr U$ is
therefore zero. The invariant restriction is injective. Its image
is exactly $J$, since an invariant representative restricts to
$J$ by definition and every element of $J$ has such a representative.
Both maps are induced by polynomial quotient maps, so they preserve
products, sums, complex scalars and degrees. No action of $g$ on
$R$ is asserted in the size-five case. In particular the second
map retains precisely the invariant image of the original
observation, with its complete restriction kernel already computed.

\subsection{All original degrees and the exact kernel dimension}

The map $\Theta_k$ and the original observation in OQC8 identify
$\dim J_k$ with $\operatorname{rank}\Lambda_k$ and identify
$R_k/J_k$ with $(\ker\Lambda_k)^\vee$. Multiplication by the
nonzero polynomial $\mathcal H_Q$ is injective in $\mathscr P$.
Taking dimensions in OQX1 thus gives, for every nonnegative degree
$d$, the exact invariant restriction rank
$\dim J_d=d_{5,d,0}-d_{5,d-10,0}$. At the original degrees this
evaluates to
\[
 \boxed{\quad
 \operatorname{rank}\Lambda_k
      =d_{5,k,0}-d_{5,k-10,0}
      =k^2-6k+17,\qquad
 \dim\ker\Lambda_k=8k-16,
 \quad k=4l+1,\ l\geq1,
 \quad}\tag{OQX3}
\]
in the size-five outcome of the actual coefficient test.
To verify every constant and degree range, for $k\geq10$ the
complete character formula OQC12 gives
\[
 d_{5,k,0}-d_{5,k-10,0}
 =\frac15\left\{\binom{k+3}{3}-\binom{k-7}{3}
                         +4(\epsilon_k-\epsilon_{k-10})\right\}.
\]
The two $\epsilon$ values are equal because their indices differ
by ten. Expanding the remaining cubics gives
\[
 \binom{k+3}{3}-\binom{k-7}{3}
 =\frac{(k^3+6k^2+11k+6)
             -(k^3-24k^2+191k-504)}6
 =5(k^2-6k+17).
\]
The only original degrees below ten are five and nine. Their
negative-degree terms in OQX1 vanish, while the same exact
character formula gives
\[
 d_{5,5,0}=(56+4)/5=12=5^2-6\cdot5+17,
 \qquad
 d_{5,9,0}=220/5=44=9^2-6\cdot9+17.
\]
This proves the first equality of OQX3 for every original $l\geq1$.
Since $\dim E_k=\dim R_k=(k+1)^2$, rank-nullity gives
$(k+1)^2-(k^2-6k+17)=8k-16$, proving the second equality.
For the original degrees already tabulated in OCS26, the exact
five-orbit values are
\[
 \begin{array}{c|r|r}
 k&\operatorname{rank}\Lambda_k&\dim\ker\Lambda_k\\\hline
 5&12&24\\
 9&44&56\\
 13&108&88\\
 17&204&120\\
 21&332&152\\
 25&492&184\\
 29&684&216
 \end{array}
\]
These are exact values for this exhaustively classified orbit
outcome, replacing the less precise bounds OQC18 wherever that
outcome occurs. The full invariant lift and its action defect
OQC15--23 remain valid with all original coefficients. The
size-one outcome retains its different exact values OQC11--13.
No choice between the two outcomes has been made here without
evaluating the original period coefficient test, and no estimate
for a metric volume has been inferred from its dimension alone.

% END COMPLETE OQX

\clearpage
% BEGIN COMPLETE PQO
\section{The actual period quadric and its cyclic continuation}
\label{sec:PQO}

This calculation uses the original simple quartet, $m=1$, and the
literal period matrix, amplitude unit and cyclic matrix of BRU and
OPG. Its formulas evaluate the transported quadric from those data.
All source orders elsewhere in the programme remain unchanged.

\subsection{Original roots, entire unit and centred phase}
Order the roots and keep their actual values as
\[
\begin{gathered}
 \alpha_1=\tfrac12+\delta+i\gamma,\quad
 \alpha_2=\tfrac12+\delta-i\gamma,\quad
 \alpha_3=\tfrac12-\delta+i\gamma,\quad
 \alpha_4=\tfrac12-\delta-i\gamma,\\
 0<\delta<\tfrac12,\quad\gamma>2,\quad
 h(s)=\prod_{j=1}^4(s-\alpha_j),\quad g=2\xi,\quad v_h=g/h.
\end{gathered}\tag{PQO1}
\]
The four zero orders of $g$ are complete and equal to one.
Thus $v_h$ is entire and has nonzero value at each displayed root.
Put $w=s-1/2$, and define the following original constants:
\[
\begin{gathered}
 z_+=\delta+i\gamma,\quad z_-=\delta-i\gamma,\quad
 \lambda=z_+^2,\quad\mu=z_-^2,\\
 A=2(\gamma^2-\delta^2),\qquad B=(\delta^2+\gamma^2)^2,\qquad
 \mathfrak a=\Phi_h(1/2)=\frac1{160}+\frac A{24}+\frac B2,\\
 h(1/2+w)=w^4+Aw^2+B,\qquad
 \Phi_h(1/2+w)=\mathfrak a+\frac{w^5}{5}+\frac{Aw^3}{3}+Bw.
\end{gathered}\tag{PQO2}
\]
The product of the two quadratics
$((w-\delta)^2+\gamma^2)((w+\delta)^2+\gamma^2)$ proves the
quartic identity. Integrating it and using $\Phi_h(0)=0$
proves the phase and constant in PQO2. In particular the translation
has not removed the original constant $\mathfrak a$.

The theta integral IPM.2 has real coefficients, real theta kernel,
and an expression unchanged by $s\mapsto1-s$. It therefore gives
$g(\bar s)=\overline{g(s)}$ and $g(1-s)=g(s)$ directly; the same
two identities hold for $h$. Taking their entire quotient gives
\[
 v_h(\alpha_1)=v_h(\alpha_4)=a,\qquad
 v_h(\alpha_2)=v_h(\alpha_3)=\bar a,\qquad a\ne0.
 \tag{PQO3}
\]
These are the full actual unit values. Their magnitude and phase
are not assigned. Set $t_+=a^{-2}$ and $t_-=\bar a^{-2}$.
For $j=0,1,2,3$ retain
\[
 \sigma_j=t_+\lambda^j-t_-\mu^j,\qquad
 \sigma_2=-A\sigma_1-B\sigma_0,\qquad
 \sigma_3=(A^2-B)\sigma_1+AB\sigma_0.
 \tag{PQO4}
\]
The recurrence follows because $\lambda+\mu=-A$ and
$\lambda\mu=B$. Each $\sigma_j$ is purely imaginary.
Furthermore $\sigma_0$ and $\sigma_1$ cannot both vanish:
the first would give $t_+=t_-\ne0$, and the second would then
give $t_+(\lambda-\mu)=0$, whereas
$\lambda-\mu=4i\delta\gamma\ne0$.

\subsection{The original periods as an absolutely convergent series}
Fix $u\ne0$ on its original logarithm branch. Let
$\ell_j(r)=r\exp(i(\arg u+\pi+2\pi j)/5)$,
$\Gamma_j=-\ell_0+\ell_j$, $1\leq j\leq4$.
Let $\Pi$ integrate the degree-less-than-four representative
against $e^{\Phi_h(s)/u}\,ds$ on these original oriented contours.
Write $F_{ja}=\zeta^{ja}-1$, $\zeta=e^{2\pi i/5}$,
$1\leq j,a\leq4$. Direct finite geometric sums prove
\[
 (F^{-1})_{aj}=\frac15\zeta^{-ja},\qquad
 CF=F D,\qquad D=\operatorname{diag}(\zeta,\zeta^2,\zeta^3,\zeta^4),
 \tag{PQO5}
\]
where $(Cx)_j=x_{j+1}-x_1$ for $j<4$ and $(Cx)_4=-x_1$.
For example multiplying the asserted inverse by $F$ gives
$\frac15\sum_{j=1}^4\zeta^{-ja}(\zeta^{jb}-1)=\delta_{ab}$;
substitution into each row gives $CF=FD$.

Let $T$ be the exact increasing-coefficient translation matrix
$T_{rb}=\binom br2^{-(b-r)}$ for $0\leq r\leq b\leq3$, and
zero otherwise. Thus it sends coefficients of $p(s)$ to those
of $p(1/2+w)$, with determinant one. Define the centred
Fourier-period matrix
\[
 P_{ar}(u)=\frac15\sum_{j=1}^4\zeta^{-ja}
 \int_{\Gamma_j}e^{(w^5/5+Aw^3/3+Bw)/u}w^r\,dw,
 \quad 1\leq a\leq4,\quad0\leq r\leq3.
 \tag{PQO6}
\]
The contours in this formula start at $w=0$. They give precisely
the translated original period by
\[
 F^{-1}\Pi=e^{\mathfrak a/u}P(u)T.
 \tag{PQO7}
\]
To justify this exact contour translation, the substitution
$w=s-1/2$ first gives rays starting at $-1/2$. Join each shifted
ray at a large radius to the corresponding unshifted ray. On the
joining segment their arguments differ by $O(1/r)$, and the
leading real phase is $-r^5/(5|u|)+O(r^4)$; all lower phase
terms and polynomial factors are dominated by this negative
leading term. Those outer joining integrals tend to zero.
The two common-origin joining integrals cancel in
$-\ell_0+\ell_j$ with their opposite orientations. Cauchy's
theorem for the entire integrand proves equality in the limit.
Substitution of the polynomial gives exactly the matrix $T$.

For every integer $L\geq1$ define, on the same branch,
\[
 G_L(u)=u^{L/5}5^{L/5-1}e^{\pi iL/5}\Gamma(L/5).
 \tag{PQO8}
\]
Integration of the leading phase on the two rays gives
$\int_{\Gamma_j}e^{w^5/(5u)}w^{L-1}dw
 =G_L(u)(\zeta^{jL}-1)$ by the real substitution
$r^5/(5|u|)$ on each ray, including its phase in $dw$.
Expanding the remaining two exponentials now gives the exact
entries
\[
\boxed{
 P_{ar}(u)=
 \sum_{\substack{p,q\geq0\\r+3p+q+1\equiv a\pmod5}}
 \frac{A^pB^q}{3^p p!q!u^{p+q}}\,
                 G_{r+3p+q+1}(u).
}\tag{PQO9}
\]
Terms with $r+3p+q+1\equiv0$ have zero two-ray integral
in every column and contribute zero before the Fourier map.
To prove every exchange, the sum of the absolute values of the
unintegrated series on either ray is bounded by a constant times
\[
 R^r\exp\left(-\frac{R^5}{5|u|}
             +\frac{|A|R^3}{3|u|}+\frac{|B|R}{|u|}\right).
\]
Here $R$ is the radial coordinate and $r$ is the fixed column
degree. This integrable majorant
proves absolute termwise integration. It is uniform on compact
coefficient and nonzero-$u$ sets on the chosen branch. The finite
Fourier sum in PQO5 then selects exactly the stated residue class.
Thus PQO9 is an actual period formula, with all coefficients and
all terms retained.

For additional finite computation, let $I_d$ denote the vector
of the four centred ray moments of $w^d$. Integration of the
derivative of $e^{(w^5/5+Aw^3/3+Bw)/u}w^d$, with the same
vanishing outer endpoints and cancelling inner endpoints, proves
\[
 I_{d+4}+A I_{d+2}+B I_d=-ud I_{d-1}\quad(d\geq0),
 \qquad u\partial_B I_d=I_{d+1},\quad
 3u\partial_A I_d=I_{d+3}.
 \tag{PQO10}
\]
The first right side is zero at $d=0$. Coefficient derivatives
are justified by the same compact domination. These are the
original polynomial period recurrences, and they compute every
higher moment from the first four and the actual coefficients.

The determinant is especially explicit in these centred
coordinates:
\[
 \det\Pi=-(2\pi)^2u^2e^{4\mathfrak a/u},\qquad
 \det F=-25\sqrt5,\qquad
 \boxed{\det P(u)=\frac{(2\pi)^2}{25\sqrt5}\,u^2.}
 \tag{PQO11}
\]
For the first equality use the full PDE1--13 determinant
calculation at $n=4$: the phase is $e^{15\pi i}=-1$ and
the critical-value sum is $4\mathfrak a$. This critical sum
also follows directly by pairing $w$ with $-w$ in the odd
centred phase. For $F$, subtract row zero from the other rows
in the five-by-five Fourier Vandermonde and expand its first
column. Its phase is $e^{13\pi i}=-1$ and its magnitude
$5^{5/2}$, by the finite Fourier orthogonality identity.
The last equality follows by taking determinants in PQO7 and
using $\det T=1$. The Gamma grid-product proof PGCC1--6 retains
the full evaluated constant in the first equality. In particular
every inverse below exists on the full stated domain.

\subsection{The actual transformed quadric and every deciding coefficient}
Let $V$ be the centred evaluation matrix with rows
$(1,z_j,z_j^2,z_j^3)$, where
$(z_1,z_2,z_3,z_4)=(z_+,z_-,-z_-,-z_+)$. Its inverse sends
the four primary values to their unique degree-less-than-four
polynomial. Thus the actual matrix from primary coordinates
to the Fourier-period space is
\[
 B_{\rm per}=F^{-1}\Pi UJ_{\rm prim}
       =e^{\mathfrak a/u}P(u)V^{-1}
                         \operatorname{diag}(a,\bar a,\bar a,a).
 \tag{PQO12}
\]
This follows either by applying all maps to the four CRT
idempotents or by evaluating the full unit at each root.
It agrees with OCS and OPG in the given original orientation.

For $Q_0(t)=t_1t_4-t_2t_3$, define
$Q_u(x)=Q_0(B_{\rm per}^{-1}x)$. On the unchanged original
period coordinate $\mathbf y\in\mathbb C^4$ the quadratic is
$Q_u^{\rm per}(\mathbf y)=Q_0((\Pi UJ_{\rm prim})^{-1}\mathbf y)
 =Q_u(F^{-1}\mathbf y)$, because $FB_{\rm per}=\Pi UJ_{\rm prim}$.
Thus the exact isomorphism $\mathbf y=Fx$ carries the following
quadric to the original period space. If
$p(w)=p_0+p_1w+p_2w^2+p_3w^3$, multiplication gives
\[
 p(w)p(-w)=p_0^2+(2p_0p_2-p_1^2)w^2
                     +(p_2^2-2p_1p_3)w^4-p_3^2w^6.
\]
Consequently the exact symmetric matrix in the centred
coefficient frame is
\[
 S=\begin{pmatrix}
 \sigma_0&0&\sigma_1&0\\
 0&-\sigma_1&0&-\sigma_2\\
 \sigma_1&0&\sigma_2&0\\
 0&-\sigma_2&0&-\sigma_3
 \end{pmatrix},\qquad
 \boxed{Q_u(x)=x^{\mathsf T}H(u)x,
 \quad H(u)=e^{-2\mathfrak a/u}P(u)^{-\mathsf T}S P(u)^{-1}.}
 \tag{PQO13}
\]
Indeed applying the preceding product formula at $z_+$ and
$z_-$ with multipliers $t_+$ and $-t_-$ gives exactly $S$.
This retains both actual unit values through every $\sigma_j$.
The two parity blocks have determinants
\[
 \sigma_0\sigma_2-\sigma_1^2
          =-t_+t_-(\lambda-\mu)^2,\qquad
 \sigma_1\sigma_3-\sigma_2^2
          =-t_+t_-B(\lambda-\mu)^2.
\]
Thus $S$ and $H(u)$ are nondegenerate, with the exact determinant
\[
 \det H(u)=e^{-8\mathfrak a/u}
       \frac{B(t_+t_-)^2(\lambda-\mu)^4}{(\det P(u))^2}\ne0.
 \tag{PQO14}
\]

Write $N(u)=\operatorname{adj}P(u)$, with entries $N_{ra}$,
$0\leq r\leq3$, $1\leq a\leq4$, and retain
$\Delta(u)=\det P(u)$ from PQO11. Then every coefficient is
given by the explicit minor formula
\[
\begin{aligned}
 H_{ab}(u)=\frac{e^{-2\mathfrak a/u}}{\Delta(u)^2}
 \{&\sigma_0N_{0a}N_{0b}
 +\sigma_1(N_{0a}N_{2b}+N_{2a}N_{0b}-N_{1a}N_{1b})\\
 &+\sigma_2(N_{2a}N_{2b}-N_{1a}N_{3b}-N_{3a}N_{1b})
                         -\sigma_3N_{3a}N_{3b}\}.
\end{aligned}\tag{PQO15}
\]
This is matrix multiplication in PQO13 using
$P^{-1}=N/\Delta$. Each $N_{ra}$ is the signed original
three-by-three period minor. Its entries are given by the
absolutely convergent series PQO9 or by the original integrals
PQO6. The polynomial coefficient of $x_a^2$ is $H_{aa}$;
that of $x_ax_b$, $a<b$, is $2H_{ab}$. Both factors are
retained in this specified quadratic form.

The exact eight deciding entries are
\[
 \mathcal E(u)=
 (H_{11},H_{22},H_{33},H_{44},H_{12},H_{13},H_{24},H_{34}).
 \tag{PQO16}
\]
The cyclic action on quadrics is the pullback
$Q_u\mapsto Q_u(D^{-1}x)$; its symmetric matrix is
$D^{-\mathsf T}H(u)D^{-1}$. Since $D^5=I$ and $\det D=1$,
the line spanned by a nondegenerate form can be stable only
with character one. In fact if its multiplier is $\eta$, then
$\eta^5=1$ by the group action, while taking determinants
gives $\eta^4=1$ because $\det H(u)\ne0$. Thus $\eta=1$.
Entrywise, invariance is precisely
$(\zeta^{-(a+b)}-1)H_{ab}=0$. The only pairs with
$a+b\equiv0$ are $(1,4)$ and $(2,3)$. Therefore
\[
\boxed{\begin{gathered}
 \text{The quadric is invariant exactly on }
       Z=\{u:\mathcal E(u)=0\};\\
 u\in Z\ \Longrightarrow\
 Q_u(x)=2H_{14}(u)x_1x_4+2H_{23}(u)x_2x_3,
                 \quad H_{14}(u)H_{23}(u)\ne0;\\
 \text{at every }u\notin Z\text{ its orbit consists of five distinct quadrics.}
\end{gathered}}\tag{PQO17}
\]
Every assertion follows from the entry test and the rank-four
determinant. The orbit size divides five; it is one exactly
when the line is stable. Thus no two-member orbit occurs.
Moreover $F^{-1}C^{-1}=D^{-1}F^{-1}$ by PQO5, so this is
exactly the orbit of $Q_u^{\rm per}(C^{-j}\mathbf y)$ in the
original period coordinates, with the same five indexed maps.
PQO15--17 determine the exact exceptional locus by original
period minors and original amplitude values, without assigning
generic values to them.

\subsection{The cyclic action is the actual original continuation in \texorpdfstring{$u$}{u}}
Continue the original logarithm of $u$ by $2\pi i$. Its rays
advance from $\ell_j$ to $\ell_{j+1}$, and their oriented
differences advance by the exact original matrix $C$.
For the last row $\ell_5=\ell_0$, giving $-\Gamma_1$.
Equivalently every summand of PQO9 acquires the phase
\[
 e^{2\pi i((r+3p+q+1)/5-p-q)}
       =\zeta^{r+3p+q+1}=\zeta^a.
\]
The original constant $e^{\mathfrak a/u}$ is single-valued
under this continuation. Hence
\[
\begin{gathered}
 P(\log u+2\pi i)=DP(\log u),\quad
 B_{\rm per}(\log u+2\pi i)=DB_{\rm per}(\log u),\\
 H(\log u+2\pi i)=D^{-\mathsf T}H(\log u)D^{-1}.
\end{gathered}\tag{PQO18}
\]
The contour argument and the convergent series prove the same
map. Thus PQO17 describes the actual period continuation of
this original quadric. It introduces no freely chosen cyclic
action on its coefficients.

\subsection{The large-parameter geometry and a positive-ray obstruction}
Let $\varepsilon=u^{-1/5}$ on the same branch, and set
$D_\varepsilon=\operatorname{diag}(\varepsilon^{-1},
\varepsilon^{-2},\varepsilon^{-3},\varepsilon^{-4})$.
The exact series PQO9 gives
\[
 P(u)=\mathscr P(A\varepsilon^2,B\varepsilon^4)D_\varepsilon,
 \qquad
 \mathscr P(0,0)=\operatorname{diag}(g_1,g_2,g_3,g_4),
 \quad g_a=5^{a/5-1}e^{\pi ia/5}\Gamma(a/5).
 \tag{PQO19}
\]
The entries of $\mathscr P$ are entire in the two displayed
arguments, by the same absolute integrable series majorant.
The same determinant proof PQO11 applies to every complex pair
of lower odd-phase coefficients, since the centred phase remains
odd and its critical values pair to zero with full multiplicities.
Thus their determinant is the nonzero constant
$(2\pi)^2/(25\sqrt5)$, by PQO11 and $\det D_\varepsilon=u^2$.
Their inverse is therefore holomorphic everywhere, given by its
adjugate divided by that precise constant. Near $\varepsilon=0$
it is $\operatorname{diag}(g_a^{-1})+O(\varepsilon^2)$.
This bound follows from its convergent power series after
substituting $(A\varepsilon^2,B\varepsilon^4)$; it holds on
a full small complex disc, with the fixed original $A,B$.

If $\sigma_0\ne0$, PQO13 and PQO19 give
\[
 e^{2\mathfrak a/u}H_{11}(u)
              =\frac{\sigma_0}{g_1^2}\varepsilon^2
                                     +O(\varepsilon^4).
\]
If $\sigma_0=0$, then $\sigma_1\ne0$ by PQO4, and they give
\[
 e^{2\mathfrak a/u}H_{22}(u)
              =-\frac{\sigma_1}{g_2^2}\varepsilon^4
                                     +O(\varepsilon^6).
 \tag{PQO20}
\]
For verification of the orders, insert
$P^{-1}=D_\varepsilon^{-1}\mathscr P^{-1}$ into PQO13.
The middle matrix has entries $S_{rs}\varepsilon^{r+s+2}$.
Its first possible degree is two, from $S_{00}$, and its
next possible degree is four, from $S_{02},S_{20},S_{11}$.
The constant term of $\mathscr P^{-1}$ is diagonal, so the
indicated leading entry is exactly the displayed nonzero
coefficient. Every further contribution has the stated higher
degree. This also proves the expansion when some other
$\sigma_j$ vanish.

Consequently for every fixed original quartet and its full
actual amplitude there exists a finite radius outside which
the actual orbit has five distinct quadrics. This conclusion
holds on every original logarithm branch: the estimate is on
the full complex $\varepsilon$ disc, and its leading coefficient
is nonzero. The set $Z$ of PQO17 is a discrete subset of the
logarithmic $u$ surface. Indeed the decisive entry just proved
is holomorphic on that connected surface and is not identically
zero; its zeros are isolated by its local power series, and
$Z$ is contained in that zero set. The explicit finite-radius
bound from the companion period-connection calculation
strengthens this argument while retaining the same actual
coefficients.

Moreover PQO18 multiplies each deciding entry by its nonzero
root-of-unity factor, so the condition $\mathcal E(u)=0$ is
unchanged by every full logarithm turn. Therefore $Z$ is the
inverse image under $\log u\mapsto u$ of a well-defined set
$Z_{\rm abs}\subset\mathbb C^*$. On a sufficiently small disc
about any nonzero $u$, choose one logarithm branch. The preceding
isolated-zero proof shows that $Z_{\rm abs}$ is discrete and
closed on that disc. The explicit radius in the companion
calculation gives $|u|<R_*$ throughout $Z_{\rm abs}$.
Consequently its only possible accumulation point in the finite
$u$-plane is zero. In particular each compact annulus about
zero contains only finitely many such points: an infinite
subset of that compact set would have a nonzero accumulation
point, contradicting the proved local isolated-zero property.

For positive real $u$ retain the actual lift
$\log u=\ln u+2\pi i\ell_{\rm br}$, $\ell_{\rm br}\in\mathbb Z$.
The exact series has an additional real structure. Every
surviving index in PQO9 is $a+5L$, and its phase at the lift
$\ell_{\rm br}=0$ is $e^{\pi ia/5}(-1)^L$. Its remaining
coefficients are real. The actual lift multiplies row $a$ by
$\zeta^{\ell_{\rm br}a}$, by PQO18. Absolute convergence
therefore proves
\[
 P(u)=E_{\ell_{\rm br}} R(u),\qquad
 E_{\ell_{\rm br}}=D^{\ell_{\rm br}}
 \operatorname{diag}(e^{\pi i/5},e^{2\pi i/5},
                       e^{3\pi i/5},e^{4\pi i/5}),
 \qquad R(u)\text{ real and invertible}.
 \tag{PQO21}
\]
Let $M$ be multiplication by $w$ in the increasing centred
polynomial frame, so
\[
 M=\begin{pmatrix}0&0&0&-B\\1&0&0&0\\0&1&0&-A\\0&0&1&0\end{pmatrix},
 \quad c_3=\frac{t_+-t_-}{4i\delta\gamma},\quad
 c_1=\frac{t_++t_-}{2}-(\delta^2-\gamma^2)c_3.
 \tag{PQO22}
\]
Both $c_1,c_3$ are real by PQO3, and the polynomial
$c_1+c_3 w^2$ takes values $t_+,t_-,t_-,t_+$ at the four
original roots. Hence $c_1M+c_3M^3$ is precisely multiplication
by $w(j_hv_h)^{-2}$, with its entire unit retained.

The companion period-connection proof establishes the exact
identity between cyclic invariance of $Q_u$ and
\[
 [P(u)(c_1M+c_3M^3)P(u)^{-1},D]=0.
 \tag{PQO23}
\]
It proves this by the flat residue/reflection form on the
original polynomial periods and retains all its constants.
Using that proved identity, PQO21 gives a further precise
obstruction on the positive ray. Since $D$ has distinct
diagonal entries, a matrix commuting with it is diagonal.
The matrix in PQO23 is $E_{\ell_{\rm br}}$ times a real matrix
times $E_{\ell_{\rm br}}^{-1}$;
if diagonal, its diagonal entries must therefore be real.
Its eigenvalues are exactly
\[
 \frac{\delta+i\gamma}{a^2},\quad
 \frac{\delta-i\gamma}{\bar a^2},\quad
 -\frac{\delta-i\gamma}{\bar a^2},\quad
 -\frac{\delta+i\gamma}{a^2},
 \tag{PQO24}
\]
as follows by evaluation at the four distinct roots of $h$.
Thus every positive-real-$u$ point in $Z$ satisfies the explicit
original-unit equation
\[
 \operatorname{Im}\left(\frac{\delta+i\gamma}{a^2}\right)=0.
 \tag{PQO25}
\]
If that imaginary part is nonzero for the original unit, the
positive ray contains no point of $Z$. When it is zero, the
four displayed eigenvalues are the two nonzero real numbers
$r,-r$, each twice, and the complete period-minor tests
PQO15--17 still determine $Z$. The equation is not assumed
to hold or fail. Together the exact minor tests, the actual
period continuation, the proved large-parameter result and
this real-structure obstruction give control of the original
quadric family with every amplitude value and phase present.

% END COMPLETE PQO

\clearpage
% BEGIN COMPLETE RPC
% Independent mathematical provider. Original periods, amplitude and branch retained.
\section{The original period quadric as a residue commutator}
\label{sec:actual-period-residue-commutator}

Retain $0<\delta<1/2$, $\gamma>2$, the ordered simple quartet
$\alpha_1=1/2+\delta+i\gamma$, $\alpha_2=1/2+\delta-i\gamma$,
$\alpha_3=1/2-\delta+i\gamma$, $\alpha_4=1/2-\delta-i\gamma$,
and $h(s)=\prod_a(s-\alpha_a)$. Put
\[
 w=s-\tfrac12,\quad r_a=\alpha_a-\tfrac12,\quad
 \beta=2(\gamma^2-\delta^2),\quad \eta=(\delta^2+\gamma^2)^2,
 \quad H(w)=w^4+\beta w^2+\eta.
 \tag{RPC1}
\]
Thus $r_4=-r_1$, $r_3=-r_2$, $r_2=\overline{r_1}$ and
$r_1^2-r_2^2=4i\delta\gamma$. No root or original coordinate is
discarded by this translation. Its coefficient isomorphism is
$J_c:\mathbb C[w]_{<4}\to E_h$, $J_cp=[p(s-1/2)]$; its inverse
is substitution $s=w+1/2$ in the unique original representative.
The primitive with its original constant is
\[
 \Phi_h(s)=A+\phi(w),\quad
 \phi(w)=\frac{w^5}{5}+\frac{\beta w^3}{3}+\eta w,\quad
 A=\Phi_h(1/2)=\frac1{160}+\frac\beta{24}+\frac\eta2.
 \tag{RPC2}
\]
The last equality is forced by $\Phi_h(0)=0$ and follows by evaluating
the displayed odd polynomial at $w=-1/2$.

Let $u\ne0$ carry the specified original branch of its logarithm.
Let $\Pi$ be the original period map on the rays and orientations
of PDE2, and let $F_{ja}=\zeta^{ja}-1$, $\zeta=e^{2\pi i/5}$,
$1\leq j,a\leq4$. Its inverse has entries
$(F^{-1})_{aj}=\zeta^{-ja}/5$: multiplying the two matrices and using
$\sum_{j=1}^4\zeta^{jb}=4$ for $5\mid b$ and $-1$ otherwise proves
this entry by entry. The literal centered coefficient-to-Fourier map is
\[
 P=F^{-1}\Pi J_c,\qquad P_0=e^{-A/u}P.
 \tag{RPC3}
\]
These are invertible by the complete original period determinant
PDE3--13. The contours for $P_0$ may be taken to be the rays based at
$w=0$ with the original five angles. Indeed translating the old rays
gives rays based at $-1/2$. Join corresponding finite endpoints and
close each contour at radius $R$. The integrand is entire. On the
bounded-length outer joining curves its exponent is
$-R^5/(5|u|)+O(R^4)$, uniformly along the joining curves, so the
outer integrals tend to zero. The common finite joining path occurs
with opposite coefficients in $-\ell_0+\ell_j$ and cancels exactly.
Cauchy's integral theorem therefore proves this contour identity,
including its signs and the factor $e^{A/u}$.

\subsection{A flat alternating pairing with all constants}

On $\mathbb C[w]/(H)$ define the bilinear, rather than sesquilinear,
pairing
\[
 \mathcal J(p,q)=\sum_{H(r)=0}\frac{p(r)q(-r)}{H'(r)}.
 \tag{RPC4}
\]
The roots are simple at the original packet. The equivalent polynomial
formula is the coefficient of $w^3$ in the remainder of $p(w)q(-w)$
modulo $H$. To prove it, use Lagrange interpolation: the coefficient
of $w^3$ in the idempotent at $r$ is $1/H'(r)$. This polynomial
formula extends the pairing to every complex value of $\beta,\eta$,
including repeated roots, without a division by a repeated-root value.
In the increasing frame $1,w,w^2,w^3$, multiplication by $w$ and
this pairing are the exact matrices
\[
 M=\begin{pmatrix}0&0&0&-\eta\\1&0&0&0\\0&1&0&-\beta\\0&0&1&0\end{pmatrix},
 \qquad
 J=\begin{pmatrix}0&0&0&-1\\0&0&1&0\\0&-1&0&\beta\\1&0&-\beta&0\end{pmatrix}.
 \tag{RPC5}
\]
For the last matrix, the coefficient functional is zero through degree
two, one in degree three, zero in degree four and six, and $-\beta$
in degree five. Inserting the sign $(-1)^j$ from the second argument
gives every displayed entry. Thus $J^{\mathsf T}=-J$, $\det J=1$,
and direct multiplication gives $M^{\mathsf T}J+JM=0$.

Allow the two original odd-phase coefficients to vary complexly,
solely to calculate the existing period map. The ray domination
$e^{-r^5/(5|u|)+O(r^3+r)}$ on compact parameter sets justifies all
coefficient derivatives and integrations by parts below. Ordinary
division of $w^{j+3}$ by $H$, followed by differential reduction,
gives
\[
 \partial_\eta P_0=\frac1uP_0M,\qquad
 \partial_\beta P_0=\frac1{3u}P_0(M^3-uN),\qquad
 N=\begin{pmatrix}0&0&1&0\\0&0&0&2\\0&0&0&0\\0&0&0&0\end{pmatrix}.
 \tag{RPC6}
\]
For clarity, the quotient derivatives for columns $j=0,1,2,3$ are
$0,0,1,2w$, respectively; hence this is the full correction matrix.
The two outer boundary terms vanish, and the two values at the common
origin cancel. No differential remainder was replaced by an arithmetic
remainder in RPC6.

There is also a direct nonvanishing proof on this entire two-parameter
family. The displayed matrices satisfy $\operatorname{Tr}M=0$,
$\operatorname{Tr}M^3=0$ and $\operatorname{Tr}N=0$. The first and
third follow from their diagonals; the second follows by multiplying
RPC5 three times, or from the vanishing third Newton sum of the even
monic polynomial $H$. Multilinearity of the determinant in its columns
then gives $\partial_\eta\det P_0=\partial_\beta\det P_0=0$ from
RPC6, without first assuming an inverse exists. At the origin the
literal Gamma ray integrals give the diagonal matrix $G_u$ defined in
RPC8, whose determinant is nonzero. Hence $P_0$ is invertible at every
point of $\mathbb C^2$, justifying the inverses throughout this family.

The matrix identity $M^{\mathsf T}J+JM=0$ also gives
$(M^3)^{\mathsf T}J+JM^3=0$. Direct multiplication of RPC5--6 gives
\[
 \partial_\beta J+\frac13(N^{\mathsf T}J+JN)=0.
 \tag{RPC7}
\]
In particular $\partial_\beta J$ has entries $(2,3)=1$, $(3,2)=-1$
when indices begin at zero, while $N^{\mathsf T}J+JN$ has the
opposite threefold entries. Differentiate the inverse matrices in
$P_0^{-\mathsf T}JP_0^{-1}$. RPC6--7 show that both its $\beta$ and
its $\eta$ derivatives vanish. This matrix is consequently constant
on the connected coefficient space $\mathbb C^2$.

Define the unchanged positive-Gamma ray constants
\[
 g_a=5^{a/5-1}e^{\pi ia/5}\Gamma(a/5),\qquad
 G=\operatorname{diag}(g_1,g_2,g_3,g_4),\quad
 G_u=\operatorname{diag}(u^{a/5}g_a)_{a=1}^4.
 \tag{RPC8}
\]
At $\beta=\eta=0$, literal substitution in each original ray integral
gives $P_0=G_u$. Every $g_a$ is nonzero, since its Gamma factor is
the integral of a strictly positive function. We have therefore proved
the complete flat-pairing identity
\[
 \begin{split}
 \Omega_0&=P_0^{-\mathsf T}JP_0^{-1}
       =G_u^{-\mathsf T}J\big|_{\beta=0}G_u^{-1},\\
 (\Omega_0)_{14}&=-\frac1{u g_1g_4},\qquad
 (\Omega_0)_{23}=\frac1{u g_2g_3},\qquad
 (\Omega_0)_{ba}=-(\Omega_0)_{ab},\\
 \Omega&=P^{-\mathsf T}JP^{-1}=e^{-2A/u}\Omega_0.
 \end{split}
 \tag{RPC9}
\]
All entries not listed are zero; the indices in RPC9 begin at one.
This formulation retains the Gamma constants directly and requires
no unevaluated normalization of a period or source mass. For
$D=\operatorname{diag}(\zeta,\zeta^2,\zeta^3,\zeta^4)$, the two
paired sums of weights are five. It follows entry by entry that
\[
 D^{\mathsf T}\Omega D=\Omega,
 \qquad D^{\mathsf T}\Omega_0D=\Omega_0.
 \tag{RPC10}
\]

\subsection{The full arithmetic unit and the exact quadratic form}

Retain the original $v_h=2\xi/h$ and its complete simple-root unit
values. Reflection and conjugation of this original function give
\[
 v_h(\alpha_1)=v_h(\alpha_4)=a,\qquad
 v_h(\alpha_2)=v_h(\alpha_3)=\overline a,\qquad a\ne0.
 \tag{RPC11}
\]
Indeed both numerator and denominator are invariant under $s\mapsto1-s$
and conjugate under $s\mapsto\overline s$, first away from the roots
and then by the removable values. The complete zero orders give
nonzero values at these simple roots. Write
\[
 c_3=\frac{a^{-2}-\overline a^{-2}}{4i\delta\gamma},\qquad
 c_1=\frac{a^{-2}+\overline a^{-2}}2-(\delta^2-\gamma^2)c_3,
 \quad \mathfrak r(w)=c_1w+c_3w^3.
 \tag{RPC12}
\]
Both coefficients are real, and $(c_1,c_3)\ne(0,0)$. Directly,
$c_1+c_3r_1^2=a^{-2}$ and $c_1+c_3r_2^2=\overline a^{-2}$;
the same equalities hold at the respective opposite roots. Thus
$\mathfrak r(M)=M\,v_h(M+1/2)^{-2}$ in the original finite algebra, with the
entire function evaluated by its complete original arithmetic unit.
RPC12 is an exact interpolation of this specific multiplier.

For $Q_0(t)=t_1t_4-t_2t_3$, let $J_{\rm prim}$ be the original
Lagrange coordinate isomorphism and set
\[
 B=F^{-1}\Pi UJ_{\rm prim},\qquad Q(x)=Q_0(B^{-1}x),\qquad
 \mathcal H=P\mathfrak r(M)P^{-1}.
 \tag{RPC13}
\]
If $p=P^{-1}x$, the primary components of $B^{-1}x$ are exactly
$p(r_1)/a,p(r_2)/\overline a,p(r_3)/\overline a,p(r_4)/a$.
Since $H'(r)=4r(r^2+\gamma^2-\delta^2)$, its values give
$r/H'(r)=1/(8i\delta\gamma)$ on roots one and four, and
$r/H'(r)=-1/(8i\delta\gamma)$ on roots two and three. Inserting
these four values into RPC4 proves the exact identity
\[
 \boxed{\quad Q(x)=4i\delta\gamma\,
       p^{\mathsf T}\mathfrak r(M)^{\mathsf T}Jp
       =4i\delta\gamma\,x^{\mathsf T}\mathcal H^{\mathsf T}\Omega x.\quad}
 \tag{RPC14}
\]
There is no conjugate transpose in this quadratic identity. Because
$\mathfrak r$ is odd, $\mathfrak r(M)^{\mathsf T}J+J\mathfrak r(M)=0$; hence the matrices in
RPC14 are symmetric. This also proves
$\mathcal H^{\mathsf T}\Omega+\Omega\mathcal H=0$.
The four eigenvalues of $\mathfrak r(M)$ are exactly
$r_1/a^2,\overline{r_1/a^2},-\overline{r_1/a^2},-r_1/a^2$.
They are nonzero, although they need not be distinct. Both $B$ and
$Q_0$ are invertible as linear and quadratic matrices, respectively,
so $Q$ always has rank four, including the possible eigenvalue collisions.

By RPC10, substituting $Dx$ into RPC14 transports $\mathcal H$ to
$D^{-1}\mathcal H D$. The resulting symmetric quadratic matrices
are equal if and only if the corresponding $\mathcal H$ matrices
are equal, because $\Omega$ is invertible. Consequently
\[
 \boxed{\quad Q(Dx)=Q(x)\quad\Longleftrightarrow\quad
 [\mathcal H,D]=0\quad\Longleftrightarrow\quad
 \mathcal H_{ab}=0\text{ for every }a\ne b.\quad}
 \tag{RPC15}
\]
This is an exact test in the actual original period matrix and unit.
The final equivalence uses the four distinct diagonal entries of $D$.
It makes no assumption that their off-diagonal tests are independent.

For completeness, a rank-four quadratic semi-invariant of cyclic
weight $j\ne0$ is impossible: its coefficient in row $j$ could pair
only with the missing weight zero, so that row is zero. Thus any
semi-invariant rank-four quadratic must have weight zero. Since five
is prime, the projective orbit of $Q$ under $D$ has size one or five,
and size one is exactly RPC15. Every failure of any displayed
off-diagonal test proves size five.

\subsection{An explicit original-parameter radius forcing orbit five}

Here we prove that RPC15 fails for all sufficiently large $|u|$,
with a completely specified finite radius for the given original
quartet and unit. The argument evaluates the original periods; it
does not choose or insert a replacement amplitude.

Set $b=\beta u^{-2/5}$, $c=\eta u^{-4/5}$, and let $T(b,c)$
be the Fourier period matrix for
$z^5/5+bz^3/3+cz$ at $u=1$, in columns $1,z,z^2,z^3$ and on
the five rays of angles $(\pi+2\pi j)/5$. Substitution
$w=u^{1/5}z$, with exactly the original branch, gives
\[
 P_0=T(b,c)S_u,\quad S_u=\operatorname{diag}(u^{a/5})_{a=1}^4,
 \quad S_uMS_u^{-1}=u^{1/5}M(b,c),\quad T(0,0)=G.
 \tag{RPC16}
\]
Every column includes its original differential factor. The matrix
$M(b,c)$ is RPC5 with $\beta,\eta$ replaced by $b,c$; in particular
$N_0:=M(0,0)$ has its three subdiagonal entries equal to one.

The following positive constants are explicit finite real integrals
or finite expressions in the original Gamma constants:
\[
 \begin{split}
 E_j&=\frac85\int_0^\infty r^j(r^3/3+r)
                    e^{-r^5/5+r^3/3+r}\,dr\quad(0\le j\le3),\\
 C_T&=2\left(\sum_{j=0}^3E_j^2\right)^{1/2},\qquad
 g_+=\max_a|g_a|,\quad g_-^{-1}=\max_a|g_a|^{-1},\\
 \kappa&=\max_{1\le a<4}|g_{a+1}/g_a|,\qquad
 K_0=GN_0G^{-1},\\
 C_K&=4C_Tg_-^{-1}+2g_+g_-^{-1}
                       +2g_+g_-^{-2}C_T,\\
 C_{K,3}&=C_K(3\kappa^2+3\kappa C_K+C_K^2),\qquad
 \epsilon_0=\min\{1,(2C_Tg_-^{-1})^{-1}\}.
 \end{split}
 \tag{RPC17}
\]
Here $g_-^{-2}$ means $(\max_a|g_a|^{-1})^2$. The negative fifth
degree exponent proves convergence of every integral. All norms in
the following estimates are the Euclidean operator norm, with the
Frobenius norm used explicitly as an upper bound where stated.

Write $\epsilon=|b|+|c|$. When $|b|,|c|\le1$, on each ray
\[
 |e^{bz^3/3+cz}-1|
 \le \epsilon(r^3/3+r)e^{r^3/3+r}.
\]
Each contour has two rays, and each entry of $F^{-1}$ has absolute
value $1/5$ with four summands. The entry in column $j$ of
$T-G$ is therefore at most $E_j\epsilon$. Summing the squares
over all four rows proves
\[
 \|T-G\|\le\|T-G\|_{\rm F}\le C_T\epsilon.
 \tag{RPC18}
\]
For $\epsilon\le\epsilon_0$, the convergent inverse series for
$I+G^{-1}(T-G)$ gives
$\|T^{-1}\|\le2g_-^{-1}$ and
$\|T^{-1}-G^{-1}\|\le2g_-^{-2}C_T\epsilon$.
Also $\|M(b,c)-N_0\|\le\epsilon$, $\|M(b,c)\|\le2$,
and $\|N_0\|=1$. Subtract the products
$K:=T M(b,c)T^{-1}$ and $K_0$ in three terms, changing the left
factor, then the middle factor, then the inverse factor. The preceding
inequalities give, with the complete RPC17 constants,
\[
 \|K-K_0\|\le C_K\epsilon,\quad
 \|K\|\le\kappa+C_K\epsilon,\quad
 \|K^3-K_0^3\|\le C_{K,3}\epsilon.
 \tag{RPC19}
\]
For the last assertion use the exact noncommutative identity
$K^3-K_0^3=(K-K_0)K^2+K_0(K-K_0)K+K_0^2(K-K_0)$.
Its norm is at most $C_K\epsilon$ times
$(\kappa+C_K)^2+\kappa(\kappa+C_K)+\kappa^2$, as claimed.

The original multiplication and full unit in RPC13 now give exactly
\[
 \mathcal H=c_1u^{1/5}K+c_3u^{3/5}K^3,
 \qquad (K_0^3)_{41}=g_4/g_1,\quad (K_0)_{21}=g_2/g_1.
 \tag{RPC20}
\]
The scalar $e^{A/u}$ cancels from this conjugation, while it remains
in the quadratic form through RPC9. Put $B_*=\beta+\eta>0$.
For $|u|\ge1$, $\epsilon\le B_*|u|^{-2/5}$.

If $c_3\ne0$, define the completely specified radius
\[
 \begin{split}
 T_*&=\max\left\{1,\frac{B_*}{\epsilon_0},
 \frac{2\bigl(|c_3|C_{K,3}B_*+|c_1|(\kappa+C_K)\bigr)}
 {|c_3g_4/g_1|}\right\},\qquad R_*=T_*^{5/2}.
 \end{split}
 \tag{RPC21}
\]
For every $u$ on its original branch with $|u|\ge R_*$,
RPC19--20 imply
\[
 \left|u^{-3/5}\mathcal H_{41}-c_3g_4/g_1\right|
 \le |u|^{-2/5}\bigl(|c_3|C_{K,3}B_*
                           +|c_1|(\kappa+C_K)\bigr)
 \le\tfrac12|c_3g_4/g_1|.
 \tag{RPC22}
\]
In particular $\mathcal H_{41}\ne0$. In the quadratic matrix
RPC14 this is the nonzero $(1,1)$ entry
$4i\delta\gamma\,\mathcal H_{41}\Omega_{41}$, whose cyclic
weight is two.

If $c_3=0$, then $c_1\ne0$. In this case define
\[
 T_* =\max\left\{1,\frac{B_*}{\epsilon_0},
                   \frac{2C_KB_*}{|g_2/g_1|}\right\},
 \qquad R_*=T_*^{5/2}.
 \tag{RPC23}
\]
For every $|u|\ge R_*$, the same estimates give
\[
 \left|u^{-1/5}\mathcal H_{21}-c_1g_2/g_1\right|
       \le |c_1|C_KB_*|u|^{-2/5}
       \le\tfrac12|c_1g_2/g_1|.
 \tag{RPC24}
\]
Thus $\mathcal H_{21}\ne0$. The $(1,3)$ entry of the quadratic
matrix is then $4i\delta\gamma\mathcal H_{21}\Omega_{23}$,
which has cyclic weight four; symmetry retains its $(3,1)$ partner.
Both cases prove the unconditional conclusion for the actual original
parameters and the explicit radius belonging to those parameters:
\[
 \boxed{\quad |u|\ge R_*\quad\Longrightarrow\quad
       \#\{[Q(D^jx)]:0\le j<5\}=5.\quad}
 \tag{RPC25}
\]
The bracket denotes the projective class of the quadratic polynomial.
RPC21 or RPC23 supplies its radius for every nonzero original unit;
neither case assumes a condition on an uncomputed period coefficient.

Finally, the original periods are holomorphic on the connected
logarithm cover of $\mathbb C^*$, and on each branch domain: locally the contours may
be held in common decay sectors and differentiated under the dominated
integral, and contour deformation identifies these local definitions.
Their determinant never vanishes. Hence all entries of $\mathcal H$
are holomorphic there. RPC22 or RPC24 proves that its respective
decisive entry is not identically zero on the logarithm cover and,
by the identity theorem on that connected cover, on any open branch domain. The zeros
of that one entry are isolated, by its convergent local power series.
The exact exceptional set
\[
 \mathcal E=\{u:\mathcal H_{ab}(u)=0\text{ for all }a\ne b\}
 \tag{RPC26}
\]
is a subset of those isolated zeros and satisfies $|u|<R_*$.
Thus it has no accumulation inside the branch domain; accumulation
at $u=0$ or a branch boundary has not been excluded. RPC15 is the
full calculation specifying this set. No assertion that it is empty
at arbitrary finite complex $u$ is made or used.

% END COMPLETE RPC

\clearpage
% BEGIN COMPLETE OKV
% Complete finite estimate on the original m=1 packet and its actual kernel.
% Original source and quotient providers: IFB1--11 and GMB1--16.
\section{A finite bound for the original kernel volume from its actual dimension}
\label{sec:OKV}

\paragraph{The original packet, sources, and exact coordinate map.}
Throughout this calculation the original multiplicity is $m=1$ and
\[
 k=4l+1,\quad l\ge1,\quad q=(k+1)^2\ge36,\quad c=k/2,
 \qquad 0<\delta<1/2,\quad\gamma>2.
 \tag{OKV1}
\]
Retain every original root and its two coordinates:
\[
\begin{aligned}
 \chi(S)&=\prod_{a,b=0}^k
  [S-c-(2a-k)\delta-i(2b-k)\gamma],\\
 \omega_{ab}&=(2b-k)\gamma-i(2a-k)\delta,\\
 P(y)&=i^{-q}\chi(c+iy)=\prod_{a,b=0}^k(y-\omega_{ab}).
\end{aligned}
 \tag{OKV2}
\]
The identity follows factor by factor from
$iy-(2a-k)\delta-i(2b-k)\gamma=i(y-\omega_{ab})$;
it retains the full phase $i^q$. The algebra map
\[
 \Phi:\mathbb C[y]/(P)\longrightarrow E_\chi:=\mathbb C[S]/(\chi),
 \qquad [r(y)]\longmapsto[r((S-c)/i)]
 \tag{OKV3}
\]
is well defined because $P((S-c)/i)=i^{-q}\chi(S)$.
Its inverse is $[f(S)]\mapsto[f(c+iy)]$, as verified by substitution.
Both maps preserve degrees and are inverse on each polynomial ring
before taking the quotients. In particular division by $P$ is the
original remainder operation through this exact coordinate map.

Use only the two original source orders $s\in\{1,k\}$ and write
\[
 b_s=s/2,\qquad M_s=(2\pi)^{s/2},\qquad
 dm_{0,s}(y)=\frac{M_s2^{b_s-2}}{\pi\Gamma(b_s)}
     \left|\Gamma\left(\frac{b_s}{2}+\frac{iy}{2}\right)\right|^2dy.
 \tag{OKV4}
\]
The exact source calculation IFB2--5 supplies mass $M_s$, the transform
$\int e^{ty}dm_{0,s}=M_s(\cos t)^{-b_s}$ for $|t|<\pi/2$, and
the original monic orthogonal polynomials and norms
\[
\begin{gathered}
 p_0=1,\quad p_1=y,\qquad
 p_{d+1}=yp_d-d(b_s+d-1)p_{d-1},\\
 h_d^{(s)}=\int|p_d|^2dm_{0,s}=M_s d!(b_s)_d,\qquad
 \varphi_d(y)=p_d(y)/\sqrt{h_d^{(s)}},\qquad
 u_d(S)=\varphi_d((S-c)/i).
\end{gathered}
 \tag{OKV5}
\]
All these integrals are on the same real line, with the original
source mass. For example the stated orthogonality and norms also
follow directly from its transform and the exact generating function
\[
 \sum_{d\ge0}p_d(y)\frac{z^d}{d!}
       =(1+z^2)^{-b_s/2}e^{y\arctan z}.
\]
Indeed the source integral of the product at sufficiently small real
$z,w$ is $M_s(1-zw)^{-b_s}$. The cosine addition identity and the
displayed transform give this equality; the source exponential bound
justifies each finite coefficient derivative under the integral.
The coefficients of $z^dw^e$ then give zero for $d\ne e$ and
$M_s d!(b_s)_d$ for $d=e$. The initial coefficients are $1,y$,
and differentiating the generating function gives precisely the
recurrence in OKV5. Thus the source norms used below retain their
original constants and all finite-degree factors.

Let $C_s:\mathbb C^q\to E_\chi$ have columns $[u_0],\ldots,[u_{q-1}]$
in the original frame $(1,S,\ldots,S^{q-1})$. Its exact leading
coefficients give
\[
 \det C_s=i^{-q(q-1)/2}\prod_{d=0}^{q-1}(h_d^{(s)})^{-1/2}\ne0.
\]
Retain the original GMB vectors and matrices, including the last row:
\[
 v_j=C_s^{-1}[u_{q+j}],\qquad
 K_j=I_q+\sum_{a=0}^{j-1}v_av_a^*,\qquad 0\le j\le q+1.
 \tag{OKV6}
\]
The vector $v_q$ has original polynomial degree $2q$.

\paragraph{Exact remainder coefficients in all original roots.}
Index the full list in OKV2 by $\omega_1,\ldots,\omega_q$, with
each pair $(a,b)$ occurring exactly once. Define the elementary and
complete homogeneous polynomials by their finite sums
\[
 e_h(\omega)=\sum_{1\le a_1<\cdots<a_h\le q}
                      \omega_{a_1}\cdots\omega_{a_h},\qquad
 h_t(\omega)=\sum_{n_1+\cdots+n_q=t}\prod_{a=1}^q\omega_a^{n_a},
 \tag{OKV7}
\]
with $e_0=h_0=1$, $e_h=0$ outside $0\le h\le q$, and
nonnegative integer $n_a$ in the second sum. The notation $h_t(\omega)$
in this paragraph is a polynomial in the roots; the original squared
norm continues to be $h_d^{(s)}$ of OKV5. Multiplication of formal
power series gives the exact identity
\[
 \left(\sum_{h=0}^q(-1)^he_h(\omega)z^h\right)
 \left(\sum_{t\ge0}h_t(\omega)z^t\right)=1,
\]
because each geometric series $(1-\omega_a z)^{-1}$ multiplies its
own factor $1-\omega_a z$. Consequently, for $D=q+j$ with $0\le j\le q$,
the quotient and remainder of $y^D$ upon division by $P$ are exactly
\[
\begin{aligned}
 Q_D(y)&=\sum_{t=0}^jh_t(\omega)y^{j-t},\qquad
 \operatorname{rem}_P(y^D)=y^D-P(y)Q_D(y),\\
 [y^a]\operatorname{rem}_P(y^D)
   &=-\sum_{t=0}^j h_t(\omega)(-1)^{D-t-a}e_{D-t-a}(\omega),
       \qquad 0\le a<q.
\end{aligned}
 \tag{OKV8}
\]
To verify the quotient assertion, the power-series identity makes
the coefficients of $y^D,y^{D-1},\ldots,y^q$ in $P Q_D$ equal
to those of $y^D$. The remaining coefficients are exactly the last
line of OKV8. Their degree is less than $q$, so uniqueness of monic
polynomial division proves both assertions. For $0\le D<q$ the
remainder is $y^D$ itself.

The actual maximal root modulus is
\[
 R=k\sqrt{\delta^2+\gamma^2},\qquad
 D_*:=\max\{3,\sqrt{\delta^2+\gamma^2}\},\qquad T_*:=D_*q.
 \tag{OKV9}
\]
Every $|\omega_{ab}|\le R$, and equality holds at each corner,
so this expression for $R$ is exact. In particular $R\le T_*$.
Counting the summands of OKV7 proves
\[
 |e_h(\omega)|\le\binom qh R^h,\qquad
 |h_t(\omega)|\le\binom{q+t-1}{t}R^t.
\]
The second count follows by placing $q-1$ separators among $t+q-1$
positions to encode $(n_1,\ldots,n_q)$. For fixed $j$, the binomial
$\binom{q+t-1}{t}$ is increasing in $0\le t\le j$, because its
successive ratio is $(q+t)/(t+1)\ge1$. In each nonzero summand
of OKV8 the total root exponent is $(D-t-a)+t=D-a$.
As $t$ varies the indices $D-t-a$ are distinct, so the sum of
their $\binom q{D-t-a}$ is at most $2^q$. It follows that
\[
 \left|[y^a]\operatorname{rem}_P(y^{q+j})\right|
 \le2^q\binom{q+j-1}{j}R^{q+j-a}
 \le2^{3q-1}R^{q+j-a},\qquad 0\le a<q,\quad0\le j\le q.
 \tag{OKV10}
\]
For the last inequality use
$\binom{q+j-1}{j}\le2^{q+j-1}\le2^{2q-1}$, by the binomial
expansion of $(1+1)^{q+j-1}$. These estimates retain the exact
coefficient formula OKV8, including every signed cancellation.

\paragraph{Original Gamma norms of the low monomials.}
The same recurrence and norms in OKV5 give, without changing the measure,
\[
 y\varphi_a=\sqrt{(a+1)(b_s+a)}\,\varphi_{a+1}
          +\sqrt{a(b_s+a-1)}\,\varphi_{a-1},
 \quad \varphi_{-1}=0.
 \tag{OKV11}
\]
For a polynomial of degree at most $q-2$, its orthonormal coefficient
vector is therefore sent by multiplication by $y$ to the sum of
one upward and one downward weighted shift, with output degree at most
$q-1$. Each shift has operator norm at most
$\sqrt{q(b_s+q)}$: the squared norm of its coefficient vector is
bounded by the largest squared weight times the input squared norm.
Since $s\le k<q$ gives $b_s\le q/2$, the triangle inequality yields
\[
 \|yf\|_{L^2(m_{0,s})}
       \le2\sqrt{q(b_s+q)}\,\|f\|_{L^2(m_{0,s})}
       \le3q\,\|f\|_{L^2(m_{0,s})}
       \le T_*\|f\|_{L^2(m_{0,s})}.
\]
Here $2\sqrt{3/2}=\sqrt6<3$. Starting from
$\|1\|_{L^2(m_{0,s})}=\sqrt{M_s}$ and applying this inequality
to $1,y,\ldots,y^{a-1}$ proves
\[
 \boxed{\|y^a\|_{L^2(m_{0,s})}\le\sqrt{M_s}\,T_*^a,
                      \qquad0\le a<q.}
 \tag{OKV12}
\]
All intermediate polynomial degrees satisfy the stated operator domain.

For any polynomial $f=\sum_a f_a y^a$, write
$\mathfrak c_{T_*}(f)=\sum_a|f_a|T_*^a$. This is a finite
coefficient sum, not a change of source or quotient norm.
The exact original recurrence gives
\[
 \mathfrak c_{T_*}(p_{d+1})
 \le T_*\mathfrak c_{T_*}(p_d)
       +d(b_s+d-1)\mathfrak c_{T_*}(p_{d-1}).
\]
For $1\le d<2q$ its actual positive coefficient is bounded by
$d(b_s+d-1)\le2q(b_s+2q)\le5q^2\le T_*^2$.
The initial values are $1,T_*$. Induction now proves
\[
 \boxed{\mathfrak c_{T_*}(p_d)\le(2T_*)^d,
                            \qquad0\le d\le2q.}
 \tag{OKV13}
\]
Indeed substituting the two preceding bounds in the recurrence gives
$T_*(2T_*)^d+T_*^2(2T_*)^{d-1}
 =\tfrac34(2T_*)^{d+1}\le(2T_*)^{d+1}$.

Combining all $q$ coefficients of OKV10 with OKV12, by the triangle
inequality in the original $L^2$ space, gives
\[
 \|\operatorname{rem}_P(y^D)\|_{L^2(m_{0,s})}
           \le\sqrt{M_s}\,q2^{3q}T_*^D\quad(0\le D\le2q).
\]
For $D\ge q$, each summand is bounded by
$\sqrt{M_s}2^{3q-1}R^{D-a}T_*^a
 \le\sqrt{M_s}2^{3q-1}T_*^D$ and there are $q$ summands.
For $D<q$, the assertion follows from OKV12 and $q2^{3q}\ge1$.
Linearity of the exact remainder map, followed by OKV13, proves
\[
 \boxed{\|\operatorname{rem}_P(p_d)\|_{L^2(m_{0,s})}
         \le\sqrt{M_s}\,q2^{3q}(2T_*)^d,
                           \qquad0\le d\le2q.}
 \tag{OKV14}
\]

\paragraph{The original remainder vectors and the full matrix norm.}
The polynomial $\operatorname{rem}_P(p_d)$ has degree less than $q$.
Its expansion in the orthonormal basis
$\varphi_0,\ldots,\varphi_{q-1}$ has squared coefficient norm equal
to its squared original $L^2(m_{0,s})$ norm, by expanding the finite
inner product and using OKV5. The exact quotient map OKV3 and
definition OKV6 therefore give, for $d=q+j$, $0\le j\le q$,
\[
 \|v_j\|_2^2
   =\frac{\|\operatorname{rem}_P(p_{q+j})\|_{L^2(m_{0,s})}^2}
                 {M_s(q+j)!(b_s)_{q+j}}
   \le\frac{q^22^{6q}(2T_*)^{2(q+j)}}{(q+j)!(b_s)_{q+j}}.
 \tag{OKV15}
\]
This equality is the actual map of the original quotient coordinates
to the remainder polynomial. The full source factor $M_s$ occurs in
both the numerator bound and its original norm denominator before
the displayed cancellation. No estimate on the conditioning of
$C_s$ or a companion matrix is used.

Since $b_s\ge1/2$, for each integer $d\ge1$ the first factor of
$(b_s)_d$ is at least $1/2$ and every subsequent factor is at least
one. Thus $d!(b_s)_d\ge1/2$. Also $T_*\ge1$ and $q+j\le2q$.
Consequently
\[
 \|v_j\|_2^2\le2q^22^{6q}(2D_*q)^{4q},\qquad0\le j\le q,
\]
and the original complete matrix obeys
\[
 \boxed{1\le\|K_{q+1}\|_2\le
   1+2(q+1)q^22^{6q}(2D_*q)^{4q}.}
 \tag{OKV16}
\]
To check its norm bound, $\|v_jv_j^*\|_2=\|v_j\|_2^2$
follows directly from Cauchy--Schwarz, with equality on a nonzero
$v_j$; it is zero for a zero vector. Add all $q+1$ summands of
OKV6 and the norm of $I_q$. Positivity also proves its lower bound.

Here is an explicit bound of the requested scale, with constants
depending only on the original two root coordinates:
\[
 \boxed{\log\|K_{q+1}\|_2
       \le C(\delta,\gamma)\,q\log(q+1),\qquad
 C(\delta,\gamma)=12+4\log(2D_*).}
 \tag{OKV17}
\]
For a direct verification, put
$U=2(q+1)q^22^{6q}(2D_*q)^{4q}\ge1$ and use
$\log(1+U)\le\log(2U)$. This upper bound is exactly
\[
 \log4+\log(q+1)+2\log q+6q\log2
       +4q\log(2D_*)+4q\log q.
\]
Write $L=\log(q+1)$. Since $q\ge36$, one has
$\log4\le L$, $\log q\le L$, $6\log2\le2L$, and $L\ge1$.
The last inequality follows already from
$\log4=\int_1^4dt/t>1$, by bounding the integrand below by
$1/2$ on $[1,2]$ and $1/4$ on $[2,4]$.
The preceding display is therefore at most
$4L+6qL+4\log(2D_*)qL$, which is bounded by the right side
of OKV17. The completely finite and sharper bound OKV16 remains
available before this logarithmic estimate.

\paragraph{The four original kernel Grams with their full coefficient frame.}
Retain the original surjection $\Lambda:E_\chi\to B_{\rm obs}$,
the actual subspace $\mathcal K_\Lambda=\ker\Lambda$, and its original
inclusion $I:\mathcal K_\Lambda\hookrightarrow E_\chi$. Put
$r_K=\dim\mathcal K_\Lambda$; no separate formula or upper bound
for this actual dimension is needed here. Use its fixed original
coefficient frame and write
\[
 T_s=C_s^{-1}I,\qquad
 G_{q+j-1}^{(s)}=C_s^{-*}K_j^{-1}C_s^{-1},\qquad
 H_j^{(s)}=I^*G_{q+j-1}^{(s)}I=T_s^*K_j^{-1}T_s,
             \qquad0\le j\le q+1.
 \tag{OKV18}
\]
For completeness, the quotient metric in this identity follows from
the original orthonormal source-coordinate map
$A_j=[I_q,v_0,\ldots,v_{j-1}]$.
It obeys $A_jA_j^*=K_j$. The minimum norm solution of $A_jz=x$
is $z=A_j^*K_j^{-1}x$: it has the required image and is orthogonal
to $\ker A_j$. Every other solution differs by such a kernel vector,
so the Pythagorean identity gives minimum squared norm $x^*K_j^{-1}x$.
The actual remainder is $C_sx$. Substitution proves the original
quotient Gram and then its full restriction in OKV18.

Let $\lambda_s=\|K_{q+1}\|_2$. The original update chain gives
$I_q=K_0\preceq K_j\preceq K_{q+1}\preceq\lambda_s I_q$.
For positive definite matrices, inversion reverses the order:
if $A\preceq B$, conjugate by $A^{-1/2}$ to obtain a matrix
with eigenvalues at least one, invert those eigenvalues and conjugate
back. Thus the full original kernel forms satisfy
\[
 \lambda_s^{-1}H_0^{(s)}\preceq H_j^{(s)}\preceq H_0^{(s)},
 \qquad H_0^{(s)}=T_s^*T_s,
 \qquad H_{j+1}^{(s)}\preceq H_j^{(s)}.
 \tag{OKV19}
\]
For $r_K>0$, the matrix $T_s$ is injective, so $H_0^{(s)}$ is
positive definite. Conjugation in OKV19 by its inverse square root
gives eigenvalues in $[\lambda_s^{-1},1]$. The determinant identity
for that exact congruence is
\[
 \det\bigl[(H_0^{(s)})^{-1/2}H_j^{(s)}(H_0^{(s)})^{-1/2}\bigr]
       =\frac{\det H_j^{(s)}}{\det(T_s^*T_s)}.
\]
Therefore
\[
 -r_K\log\lambda_s
 \le\log\det H_j^{(s)}-\log\det(T_s^*T_s)\le0.
 \tag{OKV20}
\]
In particular the determinant and conditioning of the actual
$T_s^*T_s$ are retained in this identity; no new kernel basis is
substituted for its original frame.

Define the original four-endpoint kernel volume, with its literal
degrees $q-1,q,2q-1,2q$, by
\[
 F_K^{(s)}=
 \log\det H_0^{(s)}+\log\det H_1^{(s)}
       -\log\det H_q^{(s)}-\log\det H_{q+1}^{(s)}.
 \tag{OKV21}
\]
Monotonicity in OKV19 makes each of
$\log\det H_0^{(s)}-\log\det H_q^{(s)}$ and
$\log\det H_1^{(s)}-\log\det H_{q+1}^{(s)}$ nonnegative.
Each is at most $r_K\log\lambda_s$ by OKV20. Consequently the
actual kernel volume satisfies the complete finite estimates
\[
 \boxed{\begin{aligned}
 0\le F_K^{(s)}
 &\le2r_K\log\|K_{q+1}\|_2\\
 &\le2r_K\log\left[1+2(q+1)q^22^{6q}(2D_*q)^{4q}\right]\\
 &\le2r_K C(\delta,\gamma)q\log(q+1),
                    \qquad s\in\{1,k\}.
 \end{aligned}}
 \tag{OKV22}
\]
For $r_K=0$, every kernel determinant has its original empty value
one, and $F_K^{(s)}=0$. Hence OKV22 includes this case.
The four copies of $\log\det(T_s^*T_s)$ in OKV21 have coefficients
$+1,+1,-1,-1$ and cancel exactly only after their presence in
OKV20 has been recorded.

\paragraph{The original boundary and arithmetic transitions.}
At each original degree set
\[
 Q_j^{(s)}=
   \left(\Lambda (G_{q+j-1}^{(s)})^{-1}\Lambda^*\right)^{-1},
 \qquad
 S_j^{(s)}=(G_{q+j-1}^{(s)})^{-1}\Lambda^*Q_j^{(s)}.
 \tag{OKV23}
\]
Their zero dimensional cases have their empty meaning. Surjectivity
of $\Lambda$ and positivity give the displayed inverses. Direct
multiplication proves $\Lambda S_j^{(s)}=I$ and
$I^*G_{q+j-1}^{(s)}S_j^{(s)}=0$.
For any fixed original section $S_*$ of $\Lambda$, the map $[I,S_*]$
is bijective: its inverse records $\Lambda z$ and the unique kernel
coordinate of $z-S_*\Lambda z$. The minimum section differs from
$S_*$ by a kernel column. The two frames therefore differ by a
block triangular matrix with identity diagonal and determinant one.
Taking the determinant of the full metric in the orthogonal frame
$[I,S_j^{(s)}]$ gives
\[
 \det H_j^{(s)}\det Q_j^{(s)}
       =|\det[I,S_*]|^2\det G_{q+j-1}^{(s)}.
 \tag{OKV24}
\]
This proves the exact relation with all kernel and boundary directions
retained. The same frame factor occurs at all four endpoints.
Their signed sum is consequently
\[
 F_K^{(s)}+F_B^{(s)}=\mathcal B_k^{(s)},\qquad
 F_B^{(s)}=\log\det Q_0^{(s)}+\log\det Q_1^{(s)}
               -\log\det Q_q^{(s)}-\log\det Q_{q+1}^{(s)}.
 \tag{OKV25}
\]
The original matrices $G_{q+j-1}^{(s)}$ decrease with $j$, so their
inverse forms increase. Compression by $\Lambda$ and inversion then
show that $Q_j^{(s)}$ decreases as well. Thus $F_B^{(s)}\ge0$,
and the finite companion bound is
\[
 \max\{0,\mathcal B_k^{(s)}-2r_K C(\delta,\gamma)q\log(q+1)\}
       \le F_B^{(s)}\le\mathcal B_k^{(s)}.
 \tag{OKV26}
\]
Here the finer middle bound in OKV22 can be used in place of its
last line, with the same proof.

Finally the original source identifications are
$\mathcal B_k^{(1)}=\mathcal B_k^\sigma$ and
$\mathcal B_k^{(k)}=\mathcal B_k^0$. Their already proved full
source transition is $\mathcal H_k=\mathcal B_k^\sigma-\mathcal B_k^0$,
and the original arithmetic transition is
$\mathcal B_k^{\rm ar}=\mathcal B_k^\sigma+\delta_k^\sigma$,
with its original finite signed allowance for $\delta_k^\sigma$.
Substitution of the proved exact splitting OKV25 gives, with every
source and arithmetic term still present,
\[
 \boxed{\mathcal B_k^{\rm ar}
   =F_K^{(k)}+F_B^{(k)}+\mathcal H_k+\delta_k^\sigma
   =F_K^{(1)}+F_B^{(1)}+\delta_k^\sigma.}
 \tag{OKV27}
\]
Thus OKV22 controls the original Gamma kernel summands in this
same arithmetic identity. It keeps the separate source transition
and arithmetic scalar with their established signs and bounds.
The result is finite for the actual $r_K=\dim\ker\Lambda$ and
uses no additional asserted dimension estimate or replacement source.

% END COMPLETE OKV

\clearpage
% BEGIN COMPLETE OMV
\section{The actual simple-quartet kernel and boundary at the original volume scale}
\label{sec:OMV}

This section combines the complete original period calculation,
the exact cyclic-orbit rank calculation, and the finite original
Gamma kernel estimate. The hypotheses are the actual simple-quartet
data already used in the programme, and the period range below has
an explicit bound proved from those data. No orbit condition or
kernel-rank premise is added to the statement.

\subsection{The original domain and the proved period range}
Fix the original simple quartet $m=1$, its coordinates
$0<\delta<1/2$, $\gamma>2$, and its actual nonzero amplitude value
$a=v_h(1/2+\delta+i\gamma)$, where $v_h=2\xi/h$ and every
zero order in this quartet is complete. Reflection and conjugation
give the other three unit values exactly as in PQO3 and RPC11.
For every integer $l\geq1$ retain
\[
 k=4l+1,\qquad q=(k+1)^2,\qquad c=k/2,
 \qquad s\in\{1,k\}.
 \tag{OMV1}
\]
The polynomial $\chi$, the source measures $m_{0,s}$, their masses
$(2\pi)^{s/2}$, the original observation $\Lambda_k:E_\chi\to B_k$,
its kernel, and all four endpoint degrees are unchanged.

Let $R_*$ be the finite positive real number defined by RPC17 and
RPC21 when the actual $c_3$ of RPC12 is nonzero, and by RPC17 and
RPC23 when it is zero. Those two explicitly evaluated alternatives
exhaust the original nonzero unit, because $c_1,c_3$ cannot both
vanish. In particular $R_*$ depends on the fixed quartet and its
actual unit and is independent of $k$. Its constants are finite
positive Gamma values and the convergent positive integrals RPC17.
The full period proof RPC18--25 establishes
\[
 |u|\geq R_*\quad\Longrightarrow\quad
 \#\{[Q_u(D^j x)]:0\leq j<5\}=5
 \tag{OMV2}
\]
on every original logarithm branch. To recall the actual nonvanishing
calculation, RPC22 bounds the error in
$u^{-3/5}\mathcal H_{41}$ by half the nonzero value
$|c_3g_4/g_1|$ in the first alternative. RPC24 bounds the error
in $u^{-1/5}\mathcal H_{21}$ by half $|c_1g_2/g_1|$ in the
second. Each is a nonzero off-diagonal entry of the exact
$\mathcal H=P(c_1M+c_3M^3)P^{-1}$; the flat residue identity
RPC14--15 makes it a nonzero noninvariant coefficient of the actual
quadric. The degree-five cyclic action and its nondegenerate
quadratic matrix then give OMV2. All powers of $u$ use its original
branch, and the absolute error estimates hold for every such branch.

The exact invariant-ideal proof OQX1--3 now applies at every point
of this proved domain. It gives, for all the original degrees OMV1,
\[
 \boxed{\quad \dim B_k=k^2-6k+17,\qquad
                 r_K:=\dim\ker\Lambda_k=8k-16.\quad}
 \tag{OMV3}
\]
Specifically, an invariant polynomial vanishing on the original
quadric is divisible by all five distinct original orbit factors;
its quotient by their invariant product is invariant. Thus the
actual observed rank is $d_{5,k,0}-d_{5,k-10,0}$, with negative
degrees zero. The full character calculation OQX3 evaluates this
difference to $k^2-6k+17$, including the two original low degrees
$5,9$. Subtraction from $q=(k+1)^2$ proves the second equality.
This is the dimension of the kernel of the literal period observation,
not an assigned ambient dimension or a presumed maximal rank.

\subsection{A uniform finite bound on the actual kernel volume}
Define, from the original root coordinates only,
\[
 D_*=\max\{3,\sqrt{\delta^2+\gamma^2}\},\quad
 C_*=12+4\log(2D_*),\qquad
 \mathfrak M_k=(16k-32)C_*q\log(q+1).
 \tag{OMV4}
\]
Retain the exact original source-coordinate vectors
$v_j=C_s^{-1}[u_{q+j}]$, $K_j=I+\sum_{a<j}v_av_a^*$,
and the full original kernel frame $T_s=C_s^{-1}I$ of OKV6,18.
The original kernel metrics at degrees $q+j-1$ are
$H_j^{(s)}=T_s^*K_j^{-1}T_s$. Their actual four-endpoint volume
satisfies
\[
 \boxed{\begin{aligned}
 0\leq F_K^{(s)}(u)
 &\leq(16k-32)\log\!\left[
       1+2(q+1)q^22^{6q}(2D_*q)^{4q}\right]\\
 &\leq\mathfrak M_k,
 \qquad s\in\{1,k\},\quad |u|\geq R_*.
 \end{aligned}}\tag{OMV5}
\]
Here the source index and the actual period dependence of the kernel
are retained. The proof is direct substitution of the now proved
$r_K=8k-16$ into OKV22. Its finite coefficient proof retains
every original root in the signed remainder formula, all Gamma
recurrence coefficients through degree $2q$, and the full source
mass before its exact norm-ratio cancellation. Its metric proof
retains $\det(T_s^*T_s)$ at each of the four degrees, and cancels
only the four copies with coefficients $+1,+1,-1,-1$.
In particular neither the conditioning nor the changing position
of the original period kernel has been suppressed. The bound is
uniform in the entire period domain OMV2 because its right side
depends only on $k,\delta,\gamma$, and the dimension in OMV3 is
the same at every point of that domain.

Write the four original total quotient baselines as
$\mathcal B_k^{(1)}=\mathcal B_k^\sigma$ and
$\mathcal B_k^{(k)}=\mathcal B_k^0$. The full determinant identity
OKV24--25, with the original fixed coefficient section, gives
\[
 F_B^{(s)}(u)=\mathcal B_k^{(s)}-F_K^{(s)}(u),\qquad
 \max\{0,\mathcal B_k^{(s)}-\mathfrak M_k\}
           \leq F_B^{(s)}(u)\leq\mathcal B_k^{(s)}.
 \tag{OMV6}
\]
Both factors in the exact splitting are nonnegative by monotonicity
of their original minimum metrics. The literal boundary term is
the actual period-vector expression
\[
 F_B^{(s)}(u)=\log\frac{
       \det\mathcal M_{2q}^{(s)}(u)\det\mathcal M_{2q+1}^{(s)}(u)}{
       \det\mathcal M_q^{(s)}(u)\det\mathcal M_{q+1}^{(s)}(u)},
 \qquad
 \mathcal M_D^{(s)}(u)=\sum_{d=0}^{D-1}
       \lambda_d^{(s)}(u)(\lambda_d^{(s)}(u))^*.
 \tag{OMV7}
\]
The vectors are the complete original primary-period coefficients
OPG14 and OCS21 on the actual image of $\Lambda_k$. Formula OPG17
proves that each consecutive determinant ratio is $1+\xi_j$;
the unchanged weights $1,2,\ldots,2,1$ give exactly OMV7.
The companion kernel vectors retain the minimum-section correction
OPR2 and the action defect OPG22/OQC22. No invariance of that kernel
under the original sum action is needed for OMV5--7.

\subsection{The two individual leading coefficients}
The complete original baseline proofs BSL16--19 and HBL establish,
for fixed original $\delta,\gamma$ and $m=1$, both source limits
\[
 \frac{\mathcal B_k^{(s)}}{q^2}\longrightarrow
 C_B=9-8\log2+F(2,\pi)>\frac{769}{17010},
 \qquad s=1\text{ and }s=k.
 \tag{OMV8}
\]
Here $F(2,\pi)$ is exactly the energy functional evaluated in those
complete proofs; its original endpoint convention and coefficients
are retained. No value of it is redefined by the present comparison.
Since $q=(k+1)^2$,
\[
 \frac{\mathfrak M_k}{q^2}
   =C_*\frac{(16k-32)\log((k+1)^2+1)}{(k+1)^2}
               \longrightarrow0.
\]
For an elementary verification, $((k+1)^2+1)<(k+2)^2$, and
the last fraction is at most $32\log(k+2)/k$. The inequality
$\log x=\int_1^x dt/t\leq2(\sqrt x-1)$ for $x\geq1$
proves that this upper bound tends to zero. The finite estimate
OMV5 and the exact identity OMV6 therefore prove
\[
 \boxed{\quad
 \frac{F_K^{(1)}(u)}{q^2},\ \frac{F_K^{(k)}(u)}{q^2}
       \longrightarrow0,\qquad
 \frac{F_B^{(1)}(u)}{q^2},\ \frac{F_B^{(k)}(u)}{q^2}
       \longrightarrow C_B.
 \quad}\tag{OMV9}
\]
These limits are uniform for all $|u|\geq R_*$ and all original
branches. Indeed the kernel error is bounded by
$\mathfrak M_k/q^2$ uniformly, while the boundary error is at
most
$|\mathcal B_k^{(s)}/q^2-C_B|+\mathfrak M_k/q^2$.
The baseline is independent of $u$. This proves the asserted
uniformity, rather than only a conclusion for each separately
fixed observation. The original combined coefficient has now been
resolved into its two individual contributions on the proved
simple-quartet period domain.

\subsection{Exact propagation to the arithmetic identity and the finer scale}
Keep the original signed transitions
$\mathcal H_k=\mathcal B_k^\sigma-\mathcal B_k^0$ and
$\delta_k^\sigma=\mathcal B_k^{\rm ar}-\mathcal B_k^\sigma$.
Their complete finite allowances, including the QGQ simple-
multiplicity refinement, remain those of the full current receiver.
The exact arithmetic receiving equality becomes
\[
 \begin{aligned}
 \mathcal B_k^{\rm ar}-F_B^{(1)}(u)
       &=F_K^{(1)}(u)+\delta_k^\sigma,\\
 \mathcal B_k^{\rm ar}-F_B^{(k)}(u)
       &=F_K^{(k)}(u)+\mathcal H_k+\delta_k^\sigma.
 \end{aligned}\tag{OMV10}
\]
This is substitution in OKV27, so every source and sign is the
original one. In particular the finite upper and lower allowances
are obtained by adding $[0,\mathfrak M_k]$ to the actual interval
for $\delta_k^\sigma$, or to the actual interval for
$\mathcal H_k+\delta_k^\sigma$, respectively. The original proved
relations $\delta_k^\sigma/(kq)\to0$ and
$\mathcal H_k/(kq)\to-\mathcal C_\Gamma/4$, together with
$k/q\to0$ and OMV5, prove that both differences in OMV10,
divided by $q^2$, tend to zero uniformly on the stated period domain.
This is an exact comparison with the full arithmetic baseline.
The literal arithmetic kernel and boundary metrics retain their
own source until their corresponding metric comparison is applied.

The finer original Gamma return is still carried by an exact mixed
identity:
\[
 \begin{aligned}
 F_B^{(1)}(u)-F_B^{(k)}(u)
       &=\mathcal H_k-
                   \bigl(F_K^{(1)}(u)-F_K^{(k)}(u)\bigr),\\
 -\mathfrak M_k\leq F_K^{(1)}(u)-F_K^{(k)}(u)
       &\leq\mathfrak M_k,\\
 \mathcal H_k-\mathfrak M_k
       \leq F_B^{(1)}(u)-F_B^{(k)}(u)
       &\leq\mathcal H_k+\mathfrak M_k.
 \end{aligned}\tag{OMV11}
\]
Subtract OMV6 for the two actual source orders to obtain the first
line. Each kernel term lies in $[0,\mathfrak M_k]$, so their
difference lies in the stated interval, proving the other two lines.
The error bound here is of order $kq\log(q+1)$, whereas the
already proved signed return $\mathcal H_k$ has order $kq$.
No sign or limiting coefficient for either finer mixed difference
is assigned by OMV11. Its explicit remaining term is the difference
of the two original kernel metrics, on the exactly calculated
$8k-16$ dimensional period kernel. The full source multiplier and
its action on this actual kernel remain the input to that next
calculation, rather than a freely selected source order or kernel.

% END COMPLETE OMV

\clearpage
% BEGIN COMPLETE AMT
% Exact original arithmetic mixed transfer. Historical sources remain unchanged.
\section{The original arithmetic kernel and boundary receive the full source comparison}
\label{sec:AMT}

This calculation uses the proved full-polynomial comparison TW12/TWA3,
the original observation OPG1--5, and the complete original Gamma
kernel estimate OKV1--27. It gives the exact maps between their source,
quotient, kernel and boundary metrics, before taking any determinant.
The subsequent simple-quartet conclusion uses the evaluated original
period test RPC1--26 and its exact rank calculation OQX1--3. Every
source order below is one of the already present orders $1,k$.

\subsection{The complete original source and the same observation}
Retain the original full-order quartet, with $0<\delta<1/2$,
$\gamma>2$, and positive integer complete zero order $m$. Thus
\[
\begin{gathered}
 \mathcal R=\{\tfrac12+\delta+i\gamma,\tfrac12+\delta-i\gamma,
              \tfrac12-\delta+i\gamma,\tfrac12-\delta-i\gamma\},\quad
 h(z)=\prod_{\alpha\in\mathcal R}(z-\alpha)^m,\quad v_h=2\xi/h,\\
 k\ge3,\quad c=k/2,\quad e=1+k(m-1),\quad q=e(k+1)^2,\\
 \lambda_{a,b}=c+(2a-k)\delta+i(2b-k)\gamma,\quad
 \chi(S)=\prod_{a,b=0}^k(S-\lambda_{a,b})^e,\quad E=\mathbb C[S]/(\chi),\\
 w_h(y)=|v_h(1/2+iy)|^2/(2\pi),\quad m_{h,k}=w_h^{*k},\quad
 d\sigma(y)=|\Gamma(1/4+iy/2)|^2dy/(2\pi).
\end{gathered}\tag{AMT1}
\]
The original source inner products, conjugate-linear in their first
argument, are
\[
 \langle P,Q\rangle_{\rm ar}
  =\int_{\mathbb R}\overline{P(c+iy)}Q(c+iy)m_{h,k}(y)\,dy,
 \qquad
 \langle P,Q\rangle_\sigma
  =\int_{\mathbb R}\overline{P(c+iy)}Q(c+iy)\,d\sigma(y).
 \tag{AMT2}
\]
No measure mass or complex-coordinate phase is changed. In particular
$\int d\sigma=\sqrt{2\pi}$. Write $\mathcal P_N$ for polynomials
in the original $S$ of degree at most $N$. In its coefficient frame
$(1,S,\ldots,S^N)$ let the two Grams be $H_N^{\rm ar},H_N^\sigma$.
For $N\ge q-1$, the literal remainder map
$J_N:\mathcal P_N\to E$ is onto and has kernel
$\chi\mathcal P_{N-q}$; at $N=q-1$ this is the zero vector space.
The frame on $E$ is $(1,S,\ldots,S^{q-1})$.

Here is the original observation being retained. Put $n=4m$,
$E_h=\mathbb C[z]/(h)$, $U=M_{j_hv_h}$, and
$\iota[P]=[P(z_1+\cdots+z_k)]\in E_h^{\otimes k}$.
The map $j_h$ records every divided Taylor coefficient through order
$m-1$ at every original root. Its unit $j_hv_h$ is invertible because
the stated order is complete. Fix the original $u\ne0$ and its
logarithm branch. With $\Phi_h'=h$ and $\Phi_h(0)=0$, let
$\ell_j$ be the outward ray of angle
$(\arg u+\pi+2\pi j)/(n+1)$ and $\Gamma_j=-\ell_0+\ell_j$.
The period map and projector are exactly
\[
\begin{gathered}
 (\Pi[P])_j=\int_{\Gamma_j}e^{\Phi_h(z)/u}
                           \operatorname{rem}_hP(z)\,dz,\quad1\le j\le n,\\
 C:\Gamma_j\mapsto\Gamma_{j+1}-\Gamma_1\ (j<n),\quad
        \Gamma_n\mapsto-\Gamma_1,\qquad
 P_k=\frac1{n+1}\sum_{a=0}^{n}(C^{-a})^{\otimes k},\\
 L=P_k\Pi^{\otimes k}U^{\otimes k}\iota,\quad
 B=\operatorname{im}L,\quad\jmath:B\hookrightarrow(\mathbb C^n)^{\otimes k},
 \quad\jmath\Lambda=L,\quad\Lambda:E\twoheadrightarrow B,\quad K=\ker\Lambda.
\end{gathered}\tag{AMT3}
\]
These are the convergent original representative integrals; their
existence and full determinant are proved in PDE and OPG. Fix the
original coefficient inclusion $I:K\hookrightarrow E$ and section
$S_*:B\to E$ with $\Lambda S_*=I_B$. Let
$r=\dim K$ and $b=\dim B=q-r$. All matrices below use these same
frames and this same full unit and period observation. Empty Grams
have determinant one, and maps to or from a zero space have their
unique empty values.

\subsection{All constants of the actual full-polynomial inequality}
Retain exactly the constants of TW and TWA:
\[
\begin{gathered}
 \alpha=\pi/2,\quad p=8m-\tfrac92,\quad B_h=42+8m,\quad M=21+4m,\\
 \vartheta_h=\int_{-1}^1w_h(y)\,dy>0,\quad
 M_\alpha=\int_{\mathbb R}e^{\alpha y}w_h(y)\,dy\in(0,\infty),\\
 c_\Gamma=\sqrt2e^{-7/6},\quad C_\Gamma=\sqrt{2\sqrt5}e^{2/3},\\
 a_k=\frac{c_h\vartheta_h^{k-3}e^{-\alpha(k-3)}(k-2)^{-B_h}}
                   {C_\Gamma2^{B_h/2}},\qquad
 A_k=\frac{C_h^{\rm tilt}}{c_\Gamma}k^{p+1}M_\alpha^{k-1},\qquad
 D_N=[1+4(N+M)^2]^M.
\end{gathered}\tag{AMT4}
\]
Here $c_h$ is the proved original convolution lower-envelope constant
TW7 and $C_h^{\rm tilt}$ is the original compact/tail maximum BT12.
Their full constructions occur in the preserved arithmetic endpoint
and BT providers; they depend on this fixed $h$ and not on $k,N,u$.
Their defining inequalities, proved there with the original numerator,
are
\[
\begin{aligned}
 m_{h,k}(y)&\ge c_h\vartheta_h^{k-3}
    e^{-\alpha(|y|+k-3)}(k-2+|y|)^{-B_h},\\
 m_{h,k}(y)&\le k C_h^{\rm tilt}M_\alpha^{k-1}
              e^{-\alpha|y|}(1+|y|/k)^{-p},\\
 c_\Gamma e^{-\alpha|y|}(1+|y|)^{-1/2}
 &\le d\sigma/dy\le
 C_\Gamma e^{-\alpha|y|}(1+|y|)^{-1/2}.
\end{aligned}\tag{AMT5}
\]
For clarity the full polynomial step is given here. The first and
third bounds, $k-2+|y|\le(k-2)(1+|y|)$, and
$(1+|y|)^{B_h}\le2^{B_h/2}(1+y^2)^M$, give
$m_{h,k}\ge a_k(1+y^2)^{-M}d\sigma/dy$. The second and third,
using $1+|y|/k\ge(1+|y|)/k$ and $p>1/2$, give
$m_{h,k}\le A_kd\sigma/dy$.
The fixed Gamma orthonormal polynomials obey
\[
 ye_j=\sqrt{(j+1)(j+1/2)}e_{j+1}
                         +\sqrt{j(j-1/2)}e_{j-1}.
\]
The two shifts on degree at most $N$ have norms at most $N+1$ and
$N$, respectively. Thus $\|yP\|_\sigma\le2(N+1)\|P\|_\sigma$.
Iterating at the successively increased degrees and expanding
$(1+y^2)^M$ proves
$\int|P|^2(1+y^2)^M d\sigma\le D_N\int|P|^2d\sigma$.
Cauchy--Schwarz applied to $|P|(1+y^2)^{-M/2}$ and
$|P|(1+y^2)^{M/2}$ then gives the reciprocal lower bound. The
coordinate map $P(S)\mapsto P(c+iy)$ has inverse
$f(y)\mapsto f((S-c)/i)$, preserves the degree and retains every
coefficient through this explicitly invertible substitution. It
therefore proves on the complete original source
\[
 \boxed{\quad \frac{a_k}{D_N}H_N^\sigma
                  \preceq H_N^{\rm ar}\preceq A_kH_N^\sigma.\quad}
 \tag{AMT6}
\]
All coefficients can be complex in this proof. In particular AMT6
holds on every affine remainder fibre, not only on monic vectors.
The two Grams are positive: the Gamma density is positive and the
lower bound in AMT5 is positive everywhere. Their finite moments
exist by the exponential upper bound.

\subsection{Exact remainder, minimum-section, kernel and boundary maps}
For $a\in\{\sigma,\mathrm{ar}\}$ put
\[
\begin{gathered}
 G_N^a=(J_N(H_N^a)^{-1}J_N^*)^{-1},\qquad
 V_N^a=(H_N^a)^{-1}J_N^*G_N^a:E\longrightarrow\mathcal P_N,\\
 H_{K,N}^a=I^*G_N^aI,\qquad
 Q_N^a=(\Lambda(G_N^a)^{-1}\Lambda^*)^{-1},\qquad
 S_N^a=(G_N^a)^{-1}\Lambda^*Q_N^a:B\longrightarrow E.
\end{gathered}\tag{AMT7}
\]
Inverses are on the indicated nonzero spaces. For an onto matrix
$J$ and positive $H$, direct multiplication gives $JV=I$ for
$V=H^{-1}J^*(JH^{-1}J^*)^{-1}$. Moreover
$Z^*HV=Z^*J^*G=0$ for $JZ=0$. Therefore every lift of $z$ is
$Vz+Z$, with squared norm
$z^*Gz+Z^*HZ$, proving its unique minimum and its Gram. Applying
this proof first to $J_N$, then to $\Lambda$, proves every formula
AMT7 and $\Lambda S_N^a=I_B$, $I^*G_N^aS_N^a=0$.
It also proves that $V_N^aS_N^a$ is the minimum lift for the
composite $\Lambda J_N$.

The exact comparison map on $\mathcal P_N$ is the identity on
original coefficients; it commutes with $J_N$ and $\Lambda J_N$.
It induces the identity on $E,K,B$. Between the two spaces of
minimum lifts the exact bijection and its inverse are
\[
 \operatorname{im}V_N^\sigma\longrightarrow\operatorname{im}V_N^{\rm ar},
 \quad V_N^\sigma z\longmapsto V_N^{\rm ar}z,
 \qquad T_N=V_N^{\rm ar}J_N\big|_{\operatorname{im}V_N^\sigma},
 \quad T_N^{-1}=V_N^\sigma J_N\big|_{\operatorname{im}V_N^{\rm ar}}.
 \tag{AMT8}
\]
Their difference $V_N^{\rm ar}z-V_N^\sigma z$ lies in
$\chi\mathcal P_{N-q}$, with its unique coefficient obtained by
division by the unchanged monic $\chi$. This also gives the
minimum-lift bijection on the preimage of $K$ and the analogous
bijection $V_N^\sigma S_N^\sigma y\mapsto
V_N^{\rm ar}S_N^{\rm ar}y$ for boundary lifts. The latter
differences lie in $\ker(\Lambda J_N)$, exactly as multiplication
by $\Lambda J_N$ verifies.

Minimize both sides of AMT6 over $J_NP=z$. Taking infima preserves
each inequality because these are the same nonempty fibres. Restrict
the resulting inequality to $z=Ix$, and independently minimize it
over $\Lambda z=y$. This proves the three full-vector comparisons
\[
\boxed{\begin{gathered}
 (a_k/D_N)G_N^\sigma\preceq G_N^{\rm ar}\preceq A_kG_N^\sigma,\\
 (a_k/D_N)H_{K,N}^\sigma\preceq H_{K,N}^{\rm ar}
                                      \preceq A_kH_{K,N}^\sigma,\\
 (a_k/D_N)Q_N^\sigma\preceq Q_N^{\rm ar}\preceq A_kQ_N^\sigma.
\end{gathered}}\tag{AMT9}
\]
The source identity is an exact bounded bijection with the two
constants in AMT6. No invariance of $K$ under an additional
multiplier was used to obtain its induced identities.

The matrix $[I,S_*]$ is invertible: its inverse sends $z$ to the
unique kernel coordinate of $z-S_*\Lambda z$ and the boundary
coordinate $\Lambda z$. Each $S_N^a-S_*$ lies in $K$, so the
change from this frame to $[I,S_N^a]$ is triangular with identity
diagonal and determinant one. The latter frame is orthogonal in
$G_N^a$ by AMT7. Consequently
\[
 \boxed{\quad\det H_{K,N}^a\det Q_N^a
       =|\det[I,S_*]|^2\det G_N^a,\qquad a\in\{\sigma,\mathrm{ar}\}.\quad}
 \tag{AMT10}
\]
This includes $r=0$ and $b=0$ and retains the original frame factor.

\subsection{The exact nested deficits and their rank factors}
Define the actual finite source determinant deficit and its width
\[
 \Xi_{k,N}=(N+1)\log A_k-
                  \log\frac{\det H_N^{\rm ar}}{\det H_N^\sigma},
 \qquad L_N=\log(A_kD_N/a_k)\ge0.
 \tag{AMT11}
\]
The last sign follows by applying AMT6 to any nonzero vector.
Put
\[
\begin{aligned}
 d_{E,N}&=q\log A_k-\log\frac{\det G_N^{\rm ar}}{\det G_N^\sigma},\\
 d_{K,N}&=r\log A_k-\log\frac{\det H_{K,N}^{\rm ar}}{\det H_{K,N}^\sigma},\\
 d_{B,N}&=b\log A_k-\log\frac{\det Q_N^{\rm ar}}{\det Q_N^\sigma}.
\end{aligned}\tag{AMT12}
\]
AMT9, by congruence with the inverse square root of each comparison
Gram, places each eigenvalue of its relative Gram in
$[a_k/D_N,A_k]$. Multiplication of all eigenvalues gives
$0\le d_{K,N}\le rL_N$, $0\le d_{B,N}\le bL_N$ and
$0\le d_{E,N}\le qL_N$.

Here is the full source-to-quotient determinant calculation, including
all relation directions. Choose an $H_N^\sigma$-orthonormal
coefficient frame $R_N$ of $\ker J_N$ solely for this computation,
and put $W_N=[R_N,V_N^\sigma(G_N^\sigma)^{-1/2}]$.
AMT7 gives $W_N^*H_N^\sigma W_N=I_{N+1}$. In this frame the
positive contraction
\[
 \mathscr D_N=A_k^{-1}W_N^*H_N^{\rm ar}W_N
       =\begin{pmatrix}F_N&C_N'\\(C_N')^*&A_N'\end{pmatrix}
       \preceq I_{N+1}
\]
has Schur complement
\[
 T_N'=A_N'-(C_N')^*F_N^{-1}C_N'
       =A_k^{-1}(G_N^\sigma)^{-1/2}G_N^{\rm ar}(G_N^\sigma)^{-1/2}.
\]
Indeed minimization in its first block is exactly the original
remainder minimum at $z=(G_N^\sigma)^{-1/2}w$.
Both $F_N$ and $T_N'$ are positive contractions: positivity follows
from positive block elimination, and the upper bound for the Schur
complement follows by testing its minimum at zero first coordinate.
Since $|\det W_N|^2=(\det H_N^\sigma)^{-1}$, elimination gives
\[
 \Xi_{k,N}=d_{R,N}+d_{E,N},\quad d_{R,N}:=-\log\det F_N\ge0,
 \qquad d_{E,N}=-\log\det T_N'\ge0.
 \tag{AMT13}
\]
When $N=q-1$, $F_N$ is empty, $d_{R,N}=0$, and the same equalities
hold without a relation inverse. The determinant identities are
independent of the auxiliary orthonormal frame.
Dividing the two instances of AMT10 proves
$d_{K,N}+d_{B,N}=d_{E,N}$. Thus the complete nested statement is
\[
\boxed{\begin{gathered}
 \Xi_{k,N}=d_{R,N}+d_{K,N}+d_{B,N},\qquad d_{R,N},d_{K,N},d_{B,N}\ge0,\\
 0\le d_{K,N}\le\min\{rL_N,d_{E,N}\}\le\min\{rL_N,\Xi_{k,N}\},\\
 0\le d_{B,N}\le\min\{bL_N,d_{E,N}\}\le\min\{bL_N,\Xi_{k,N}\}.
\end{gathered}}\tag{AMT14}
\]
In particular each metric keeps its own rank factor, and the two
deficits retain their exact common sum. Any proved finite upper
bound for the original $\Xi_{k,N}$, including the degree-specific
HMR44/TWA15 bound, can be substituted in these displayed minima.
No new density or source determinant is introduced by this substitution.

\subsection{All four original endpoints and the correlated mixed return}
Let
\[
 (N_0,N_1,N_2,N_3)=(q-1,q,2q-1,2q),\qquad
 (\epsilon_0,\epsilon_1,\epsilon_2,\epsilon_3)=(1,1,-1,-1),
\]
and define, in their original signs,
\[
\begin{aligned}
 F_K^a&=\sum_{j=0}^3\epsilon_j\log\det H_{K,N_j}^a,\qquad
 F_B^a=\sum_{j=0}^3\epsilon_j\log\det Q_{N_j}^a,\\
 \mathcal B_k^a&=\sum_{j=0}^3\epsilon_j\log\det G_{N_j}^a,
 \quad\Delta_K^\sigma=F_K^{\rm ar}-F_K^\sigma,
 \quad\Delta_B^\sigma=F_B^{\rm ar}-F_B^\sigma,
 \quad\delta_k^\sigma=\mathcal B_k^{\rm ar}-\mathcal B_k^\sigma.
\end{aligned}\tag{AMT15}
\]
AMT10 cancels its four identical frame factors only after these
signs are applied. It gives
$F_K^a+F_B^a=\mathcal B_k^a$ and
$\Delta_K^\sigma+\Delta_B^\sigma=\delta_k^\sigma$ exactly.
The spaces $\mathcal P_N$ increase with the same source norm and
remainder map, so their quotient minima decrease. Restriction to
$K$ and minimization onto $B$ preserve that order. Pairing $N_0$
with $N_2$ and $N_1$ with $N_3$ consequently proves
$F_K^a,F_B^a\ge0$.

For $X=K,B$, put $r_K=r$, $r_B=b$, and retain the exact caps
\[
 c_{X,N}=\min\{r_XL_N,d_{E,N}\},\qquad
 C_X^-=c_{X,q-1}+c_{X,q},\qquad
 C_X^+=c_{X,2q-1}+c_{X,2q}.
 \tag{AMT16}
\]
The entirely source-evaluated caps
$\widehat c_{X,N}=\min\{r_XL_N,\Xi_{k,N}\}$ and the corresponding
$\widehat C_X^\pm$ are at least these values; all conclusions below
also hold with the hats. Substituting AMT12 in AMT15 records each
endpoint and its sign:
\[
\boxed{\begin{aligned}
 \Delta_K^\sigma&=-d_{K,q-1}-d_{K,q}+d_{K,2q-1}+d_{K,2q},\\
 \Delta_B^\sigma&=-d_{B,q-1}-d_{B,q}+d_{B,2q-1}+d_{B,2q},\\
 \delta_k^\sigma&=-d_{E,q-1}-d_{E,q}+d_{E,2q-1}+d_{E,2q},\\
 -C_X^-&\le\Delta_X^\sigma\le C_X^+,\qquad X=K,B.
\end{aligned}}\tag{AMT17}
\]
The four copies of $r_X\log A_k$ cancel here because their signs
sum to zero. Since the original total correction is retained, the
joint interval is stronger than two independent intervals:
\[
\boxed{\begin{aligned}
 \max\{-C_K^-,\delta_k^\sigma-C_B^+\}
   &\le\Delta_K^\sigma\le
          \min\{C_K^+,\delta_k^\sigma+C_B^-\},\\
 \max\{-C_B^-,\delta_k^\sigma-C_K^+\}
   &\le\Delta_B^\sigma\le
          \min\{C_B^+,\delta_k^\sigma+C_K^-\}.
\end{aligned}}\tag{AMT18}
\]
These follow by subtracting the interval for either component from
the exact common sum; they require no independence assertion.

Here are the explicit original four endpoint widths:
\[
\begin{aligned}
 E_k^+&=L_{q-1}+L_q
   =2\log(A_k/a_k)+M\log[1+4(q-1+M)^2]
                         +M\log[1+4(q+M)^2],\\
 E_k^-&=L_{2q-1}+L_{2q}
   =2\log(A_k/a_k)+M\log[1+4(2q-1+M)^2]
                         +M\log[1+4(2q+M)^2].
\end{aligned}\tag{AMT19}
\]
They are the unchanged TWA5 constants, both nonnegative. AMT17
therefore gives in particular
\[
 \boxed{-rE_k^+\le F_K^{\rm ar}-F_K^\sigma\le rE_k^-,\qquad
        -bE_k^+\le F_B^{\rm ar}-F_B^\sigma\le bE_k^-.}
 \tag{AMT20}
\]
The total correction retains TWA15 with all four individual source
deficits:
\[
\begin{aligned}
 \underline\delta_k^\sigma
   &=\max\{-qE_k^+,-\Xi_{k,q-1}-\Xi_{k,q}\},\\
 \overline\delta_k^\sigma
   &=\min\{qE_k^-,\Xi_{k,2q-1}+\Xi_{k,2q}\},\qquad
 \underline\delta_k^\sigma\le\delta_k^\sigma\le\overline\delta_k^\sigma.
\end{aligned}\tag{AMT21}
\]
This also follows immediately from AMT14 and AMT17. It can be
inserted into AMT18 when an interval, rather than the actual finite
total determinant, is wanted. Thus the lower and upper endpoints
are kept in their distinct places throughout the mixed transfer.

For later use, subtraction of the exact logarithms in AMT4 gives
\[
\begin{gathered}
 \mathfrak s_h=\alpha+\log(M_\alpha/\vartheta_h),\qquad
 \mathfrak c_h=-\log M_\alpha+3\log\vartheta_h-3\alpha
       +\log\frac{C_h^{\rm tilt}C_\Gamma2^{B_h/2}}{c_hc_\Gamma},\\
 \log(A_k/a_k)=k\mathfrak s_h+(p+1)\log k
                         +B_h\log(k-2)+\mathfrak c_h,\qquad
 E_k^\pm=O_h(k+\log q).
\end{gathered}\tag{AMT22}
\]
The last statement follows from the displayed exact expressions,
since $M,p,B_h$ and the two integral constants are fixed and
$\log D_N=O_h(1+\log(N+1))$. It is a consequence of the original
constants and contains no change of measure.

\subsection{The full original multiplier and its transported observation}
On the existing tensor class $k=4l+1$, retain
\[
\begin{gathered}
 t_j=2j+\tfrac12,\quad\xi_j=c+t_j,\quad
 f_l(S)=\prod_{j=0}^{l-1}(S-c-t_j),\\
 \beta_k=\frac{(2\pi)^{k/2}}{\sqrt2\Gamma(k/2)},\quad
 dm_{0,k}(y)=\beta_k|f_l(c+iy)|^2d\sigma(y).
\end{gathered}\tag{AMT23}
\]
This is the exact Gamma recurrence identity TWA26 with masses
$(2\pi)^{k/2}$ and $\sqrt{2\pi}$. For $l=0$ both $f_l$ and
$\beta_k$ are one. Write $\mathsf E_\chi:E\to\mathbb C^q$ for
the complete divided Taylor map
$[P]\mapsto(P^{(r)}(\lambda_{a,b})/r!)_{a,b,\,0\le r<e}$.
It is invertible: a vector in its kernel is divisible by the
coprime factors of $\chi$, and dimensions agree. The full
multiplication matrix $U_f$ has, in each primary block, entries
\[
\begin{gathered}
 (U_f)_{r,a}=\begin{cases}
 f_l^{(r-a)}(\lambda)/(r-a)!,&0\le a\le r<e,\\
 0,&a>r.
 \end{cases}\\
 \det U_f=\prod_{a,b=0}^k\prod_{j=0}^{l-1}
       [(2a-k)\delta-t_j+i(2b-k)\gamma]^e\ne0.
\end{gathered}\tag{AMT24}
\]
Indeed $k$ is odd, so every $\lambda_{a,b}$ has nonzero imaginary
part, while all $\xi_j$ are real and distinct. The inverse in each
block is the finite Taylor inverse of its nonzero constant term.
The exact polynomial map and its inverse on its stated range are
\[
 \mathcal P_N\longrightarrow f_l\mathcal P_N\subseteq\mathcal P_{N+l},
 \quad P\longmapsto f_lP,
 \qquad Q\longmapsto Q/f_l.
 \tag{AMT25}
\]
It sends the original fibre $J_NP=z$ bijectively to the product
jet fibre $(\mathsf E_fQ,\mathsf E_\chi Q)=(0,U_f\mathsf E_\chi z)$.
Vanishing of the $l$ distinct $f$-values proves divisibility by
$f_l$, and the inverse of $U_f$ proves the converse fibre statement.
The original relation space maps to $f_l\chi\mathcal P_{N-q}$.
On the free $\chi$ jet coordinate $Y$ of this constrained product
fibre, the exact observed map, its kernel, and a section are
\[
\boxed{\begin{gathered}
 \widetilde\Lambda=\Lambda\mathsf E_\chi^{-1}U_f^{-1},\qquad
 \widetilde\Lambda U_f\mathsf E_\chi=\Lambda,\\
 \ker\widetilde\Lambda=U_f\mathsf E_\chi K,\qquad
 \widetilde I=U_f\mathsf E_\chi I,\qquad
 \widetilde S_*=U_f\mathsf E_\chi S_*,\quad
 \widetilde\Lambda\widetilde S_*=I_B.
\end{gathered}}\tag{AMT26}
\]
Both inclusions in the kernel equality follow by applying the
displayed invertible map and its inverse. This retains the full
original amplitude in $\Lambda$ and every derivative in $U_f$.
The resulting kernel is proved by transport; it has not been
assigned the coordinates of $\mathsf E_\chi K$.

For an explicit metric on this transported space, let $b_n(y)$ be
the fixed monic Gamma polynomials and put
$\pi_n(S)=i^n b_n((S-c)/i)$, whose leading coefficient is one
and squared norm is $\gamma_n=\sqrt{2\pi}(2n)!/4^n$.
The evaluation kernels are the full finite sums
\[
 \mathscr K_{\chi,L}=\sum_{n=0}^L
   \frac{\mathsf E_\chi\pi_n(\mathsf E_\chi\pi_n)^*}{\gamma_n},\quad
 \mathscr K_{f,L}=\sum_{n=0}^L
   \frac{\mathsf E_f\pi_n(\mathsf E_f\pi_n)^*}{\gamma_n},\quad
 \mathscr K_{\chi,f;L}=\sum_{n=0}^L
   \frac{\mathsf E_\chi\pi_n(\mathsf E_f\pi_n)^*}{\gamma_n}.
\]
The product evaluation is onto for $L=N+l\ge q+l-1$, because
a polynomial of degree below $q+l$ vanishing at all its divided
jets is divisible by the degree-$q+l$ product $f_l\chi$ and is
zero. Thus its positive kernel has positive Schur complement
\[
\begin{gathered}
 C_N=\mathscr K_{\chi,N+l}-\mathscr K_{\chi,f;N+l}
            \mathscr K_{f,N+l}^{-1}\mathscr K_{f,\chi;N+l}>0,\\
 U_N=(\mathsf E_\chi\pi_{N+j}/\sqrt{\gamma_{N+j}})_{j=1}^l,
 \quad V_N=\mathscr K_{\chi,f;N+l}\mathscr K_{f,N+l}^{-1/2},\quad
 C_N=\mathscr K_{\chi,N}+U_NU_N^*-V_NV_N^*.
\end{gathered}\tag{AMT27}
\]
Empty $f$ blocks have the identity convention. The inverse of the
joint positive kernel, evaluated on $(0,Y)$, gives the minimum
$Y^*C_N^{-1}Y$: block elimination of the $f$-block proves this
identity with every mixed entry present. Combining it with AMT25
and the coefficient $\beta_k$ proves the exact tensor-source metrics
\[
\boxed{\begin{gathered}
 G_N^0=\beta_k\mathsf E_\chi^*U_f^*C_N^{-1}U_f\mathsf E_\chi,\qquad
 H_{K,N}^0=\beta_k\widetilde I^*C_N^{-1}\widetilde I,\\
 Q_N^0=\beta_k(\widetilde\Lambda C_N\widetilde\Lambda^*)^{-1},\qquad
 S_N^0=\mathsf E_\chi^{-1}U_f^{-1}C_N\widetilde\Lambda^*
                   (\widetilde\Lambda C_N\widetilde\Lambda^*)^{-1}.
\end{gathered}}\tag{AMT28}
\]
Direct multiplication verifies the section and its orthogonality.
The formulas hold on the original $E,K,B$, with the transported
observation exactly as in AMT26. No block or mixed product was
dropped from $C_N$.

The same coefficient identity also supplies a quantitative comparison
directly with this original tensor source. Define the positive exact
finite factors
\[
 \Pi_l=\prod_{j=0}^{l-1}t_j^2,\qquad
 U_{l,N}^{\rm poly}=\prod_{j=0}^{l-1}[4(N+j+1)^2+t_j^2],\qquad
 \widetilde L_N=L_N+\log(U_{l,N}^{\rm poly}/\Pi_l).
 \tag{AMT29}
\]
Pointwise $|f_l(c+iy)|^2\ge\Pi_l$. Also
$\|(iy-t)P\|_\sigma^2=\|yP\|_\sigma^2+t^2\|P\|_\sigma^2$,
so the degree-specific multiplication estimate in AMT6, applied
successively, gives
$\|f_lP\|_\sigma^2\le U_{l,N}^{\rm poly}\|P\|_\sigma^2$.
Writing $H_N^0$ for the unchanged $m_{0,k}$ source Gram yields
\[
 \beta_k\Pi_lH_N^\sigma\preceq H_N^0
                  \preceq\beta_k U_{l,N}^{\rm poly}H_N^\sigma,
 \qquad
 \frac{a_k}{D_N\beta_k U_{l,N}^{\rm poly}}H_N^0
       \preceq H_N^{\rm ar}\preceq\frac{A_k}{\beta_k\Pi_l}H_N^0.
 \tag{AMT30}
\]
The second pair follows by using the first upper bound for its lower
inequality and the first lower bound for its upper inequality.
Minimizing and restricting with the same $J_N,I,\Lambda$ gives
AMT30 for all three metrics in AMT28. Its common upper scalar is
$A_k/(\beta_k\Pi_l)$; it is retained before its four signed copies
cancel. The same determinant proof as AMT12--17 therefore gives
\[
\begin{gathered}
 \widetilde E_k^+=E_k^+
   +\log(U_{l,q-1}^{\rm poly}U_{l,q}^{\rm poly}/\Pi_l^2),\qquad
 \widetilde E_k^-=E_k^-
   +\log(U_{l,2q-1}^{\rm poly}U_{l,2q}^{\rm poly}/\Pi_l^2),\\
 \boxed{-r\widetilde E_k^+\le F_K^{\rm ar}-F_K^0
                 \le r\widetilde E_k^-,\qquad
        -b\widetilde E_k^+\le F_B^{\rm ar}-F_B^0
                 \le b\widetilde E_k^-.}
\end{gathered}\tag{AMT31}
\]
For an explicit bound on the extra widths, $t_j\ge1/2$ gives
\[
 0\le\log(U_{l,N}^{\rm poly}/\Pi_l)
   =\sum_{j=0}^{l-1}\log\left(1+\frac{4(N+j+1)^2}{t_j^2}\right)
   \le l\log[1+16(N+l)^2].
 \tag{AMT32}
\]
Thus $\widetilde E_k^\pm=O_h(k\log(q+k))$ on the stated endpoints.
This is a proved polynomial-source estimate for the exact original
$f_l$, with no inference that its multiplication preserves $K$.

\subsection{The evaluated original simple-quartet arithmetic kernel}
Now take the original simple-quartet class $m=1$, $k=4l+1$,
$l\ge1$, $q=(k+1)^2$. The parameter $u$ is precisely the period
parameter in AMT3, not the real integration coordinate $y$.
RPC21 and RPC23 give an explicit finite radius $R_*$ for the
unchanged quartet and arithmetic unit. For reference its exact
constants and cases are recorded here. Set
\[
\begin{gathered}
 \beta=2(\gamma^2-\delta^2),\quad\eta=(\delta^2+\gamma^2)^2,
 \quad a=v_h(\tfrac12+\delta+i\gamma)\ne0,\\
 c_3=\frac{a^{-2}-\overline a^{-2}}{4i\delta\gamma},\qquad
 c_1=\frac{a^{-2}+\overline a^{-2}}2-(\delta^2-\gamma^2)c_3,\\
 g_j=5^{j/5-1}e^{\pi ij/5}\Gamma(j/5)\ (1\le j\le4),\quad
 g_+=\max_j|g_j|,\quad g_-^{-1}=\max_j|g_j|^{-1},\quad
 \kappa=\max_{1\le j<4}|g_{j+1}/g_j|,\\
 E_j^{\rm ray}=\frac85\int_0^\infty r^j(r^3/3+r)
                  e^{-r^5/5+r^3/3+r}\,dr\ (0\le j\le3),\quad
 C_T=2\left(\sum_{j=0}^3(E_j^{\rm ray})^2\right)^{1/2},\\
 C_K=4C_Tg_-^{-1}+2g_+g_-^{-1}+2g_+g_-^{-2}C_T,\quad
 C_{K,3}=C_K(3\kappa^2+3\kappa C_K+C_K^2),\\
 \epsilon_*=\min\{1,(2C_Tg_-^{-1})^{-1}\},\qquad B_*=\beta+\eta.
\end{gathered}\tag{AMT33}
\]
Here $g_-^{-2}=(g_-^{-1})^2$. All integrals converge by their
negative fifth-degree exponent, and $(c_1,c_3)\ne(0,0)$ by the
original unit values. The exact radius is
\[
 R_*=
 \begin{cases}
 \displaystyle\left[\max\left\{1,\frac{B_*}{\epsilon_*},
 \frac{2(|c_3|C_{K,3}B_*+|c_1|(\kappa+C_K))}{|c_3g_4/g_1|}\right\}\right]^{5/2},
                  &c_3\ne0,\\[6pt]
 \displaystyle\left[\max\left\{1,\frac{B_*}{\epsilon_*},
                     \frac{2C_KB_*}{|g_2/g_1|}\right\}\right]^{5/2},
                  &c_3=0.
 \end{cases}\tag{AMT34}
\]
The complete period calculation RPC16--25 proves, for every original
branch and every $|u|\ge R_*$, a nonzero original commutator entry:
its $(4,1)$ entry in the first case and its $(2,1)$ entry in the
second. RPC14--15 identifies those entries with the exact test for
the original rank-four quadric's cyclic invariance. Thus its actual
projective orbit has five elements throughout this domain. OQX1--3,
using the product of precisely these five original quadrics, proves
\[
 \boxed{r=8k-16,\qquad b=k^2-6k+17,
        \qquad m=1,\ k=4l+1,\ l\ge1,\ |u|\ge R_*.}
 \tag{AMT35}
\]
This is the evaluated original period outcome; no value of the
amplitude or of a period coefficient is chosen. The radius depends
on the fixed one-factor quartet and unit and is independent of $k$.

For the same original packet, OKV22 proves for both existing source
orders $s=1,k$, and for the actual $K$ of any rank,
\[
\begin{gathered}
 D_*=\max\{3,\sqrt{\delta^2+\gamma^2}\},\quad
 C(\delta,\gamma)=12+4\log(2D_*),\\
 Z_q=1+2(q+1)q^2 2^{6q}(2D_*q)^{4q},\\
 0\le F_K^{(s)}\le2r\log Z_q
                \le2r C(\delta,\gamma)q\log(q+1),
 \\ F_K^{(1)}=F_K^\sigma,\quad F_K^{(k)}=F_K^0.
\end{gathered}\tag{AMT36}
\]
Its proof includes the exact source coefficient map and the unchanged
kernel inclusion. Combining it with the already proved AMT17--20
gives the finite arithmetic conclusion
\[
\boxed{\begin{aligned}
 \max\{0,F_K^\sigma-C_K^-\}
       &\le F_K^{\rm ar}\le F_K^\sigma+C_K^+,\\
 0\le F_K^{\rm ar}
       &\le U_K:=2r\log Z_q+\widehat C_K^+\\
       &\le2r C(\delta,\gamma)q\log(q+1)+rE_k^-.
\end{aligned}}\tag{AMT37}
\]
The cap symbols $C_K^\pm,\widehat C_K^\pm$ in this formula mean
the endpoint sums AMT16; the constant $C_K$ without a superscript
in AMT33 belongs only to the explicit period radius. They have
different stated indices and are not identified.
The alternative finite tensor-source bound
$F_K^{\rm ar}\le2r\log Z_q+r\widetilde E_k^-$ also follows from
AMT31. Formula AMT37 retains the sharper fixed-source widths.

On AMT35 the exact ranks give
\[
 \boxed{\quad F_K^{\rm ar}-F_K^\sigma=O_h(k^2),\qquad
 F_K^{\rm ar}=O_h(kq\log(q+1)),\qquad
 \frac{F_K^{\rm ar}}{q^2}\longrightarrow0.\quad}
 \tag{AMT38}
\]
For the first estimate use $r=8k-16$ in AMT20 and AMT22,
with $\log q=2\log(k+1)$. The second follows from AMT37.
For the limit, divide its last line by $q^2$: its first term is
at most $16C(\delta,\gamma)k\log(q+1)/q\to0$, and its second
is $O_h(k^2/q^2)\to0$. Nonnegativity gives the limit by the two
bounds. These estimates are uniform over the entire original
domain $|u|\ge R_*$, since none of their right sides depends on $u$.
The tensor comparison likewise gives
$F_K^{\rm ar}-F_K^0=O_h(k^2\log q)=o(q^2)$ from AMT31--32.

\subsection{The original boundary receives the full baseline coefficient}
Before using a limit, the exact finite arithmetic identity and its
consequence are
\[
\boxed{\begin{gathered}
 F_B^{\rm ar}=\mathcal B_k^\sigma+\delta_k^\sigma-F_K^{\rm ar},\\
 \max\{0,\mathcal B_k^\sigma+\delta_k^\sigma-U_K\}
    \le F_B^{\rm ar}\le\mathcal B_k^\sigma+\delta_k^\sigma,\\
 \max\{0,\mathcal B_k^\sigma+\underline\delta_k^\sigma-U_K\}
    \le F_B^{\rm ar}\le\mathcal B_k^\sigma+\overline\delta_k^\sigma.
\end{gathered}}\tag{AMT39}
\]
The first identity follows from AMT10 and AMT15; AMT37 and
$F_K^{\rm ar},F_B^{\rm ar}\ge0$ prove the first interval, and
AMT21 proves the second. The correlated boundary return also has
the sharper exact interval AMT18, with every original signed
total correction present.

BSL13--19 supplies the complete original finite baseline and its
proved coefficient
\[
 \mathcal B_k^\sigma=V_q+\eta_q+\eta_{q+1}-\eta_1,\qquad
 \frac{\mathcal B_k^{\rm ar}}{q^2}\longrightarrow
 C_B^{\rm base}:=9-8\log2+\mathcal F(2,\pi)
                      >\frac{769}{17010}>0.
 \tag{AMT40}
\]
Here $V_q$, the three distinct original relation errors $\eta_r$,
and the energy $\mathcal F$ are exactly the BSL/HBL/GEL quantities,
whose complete definitions and proofs are in the retained closure.
The superscript on $C_B^{\rm base}$ only distinguishes that existing
baseline coefficient from the endpoint caps $C_B^\pm$ in AMT16.
Its finite BSL14 enclosure can be substituted directly for
$\mathcal B_k^\sigma$ in AMT39, keeping the signs of all three
$\eta$ terms; no new evaluation is needed for that substitution.
Subtract AMT38 from the exact original identity
$F_B^{\rm ar}=\mathcal B_k^{\rm ar}-F_K^{\rm ar}$. It proves
\[
 \boxed{\quad
 \frac{F_B^{\rm ar}}{q^2}\longrightarrow C_B^{\rm base}>0,
 \qquad
 F_K^{\rm ar}+F_B^{\rm ar}=\mathcal B_k^{\rm ar},
 \quad m=1,\ k=4l+1,\ l\to\infty,\ |u|\ge R_*.
 \quad}\tag{AMT41}
\]
The convergence is uniform over that original period domain by
AMT38 and because the full-source baseline is independent of $u$.
Thus the proved original kernel estimate now passes to the actual
arithmetic kernel, and the original arithmetic boundary carries the
entire positive $q^2$ coefficient. All earlier full-sum identities
remain exact, with their individual mixed contributions now receiving
AMT37--41 and their signed finite corrections AMT17--21.

\subsection{The same comparison evaluates the mixed Gamma return on the \texorpdfstring{$kq$}{kq} scale}
The first inequality pair AMT30 directly compares the two existing
Gamma sources on the same original coefficient space. Apply its
proved minimum and restriction maps AMT7--9 to $X=E,K,B$, writing
$H_{E,N}^a=G_N^a$, $H_{B,N}^a=Q_N^a$, and
$r_E=q,r_K=r,r_B=b$. At every original endpoint it gives
\[
 \beta_k\Pi_l H_{X,N}^\sigma\preceq H_{X,N}^0
       \preceq\beta_k U_{l,N}^{\rm poly}H_{X,N}^\sigma.
 \tag{AMT42}
\]
Define the exact relative determinant and the positive degree width
\[
\begin{gathered}
 z_{X,N}=\log\frac{\det H_{X,N}^0}{\det H_{X,N}^\sigma}
                              -r_X\log(\beta_k\Pi_l),\qquad
 \omega_N^\Gamma=\log(U_{l,N}^{\rm poly}/\Pi_l),\\
 0\le z_{X,N}\le r_X\omega_N^\Gamma,\qquad
 z_{K,N}+z_{B,N}=z_{E,N},\\
 \Omega_k^{\rm low}=\omega_{q-1}^\Gamma+\omega_q^\Gamma,\qquad
 \Omega_k^{\rm high}=\omega_{2q-1}^\Gamma+\omega_{2q}^\Gamma.
\end{gathered}\tag{AMT43}
\]
Congruence and multiplication of the eigenvalues in AMT42 prove
the two inequalities. The determinant identity AMT10, applied to
the tensor metric AMT28 as well, proves the common sum. These
facts preserve the complete original kernel and boundary rather
than assigning separate arbitrary errors to their determinants.

Use the precise return orientation
\[
 \Delta_X^\Gamma=F_X^\sigma-F_X^0\quad(X=K,B),\qquad
 \mathcal H_k=\mathcal B_k^\sigma-\mathcal B_k^0=-\mathfrak T_k.
\]
The last equality is the original TWA30/BSL18 sign. Substituting
all four $z$ terms, and canceling $\log(\beta_k\Pi_l)$ only
after the signed sum, proves
\[
\boxed{\begin{aligned}
 \Delta_X^\Gamma
   &=-z_{X,q-1}-z_{X,q}+z_{X,2q-1}+z_{X,2q},\\
 -r_X\Omega_k^{\rm low}
   &\le\Delta_X^\Gamma\le r_X\Omega_k^{\rm high},\\
 \Delta_K^\Gamma+\Delta_B^\Gamma&=\mathcal H_k,\\
 \mathcal H_k-r\Omega_k^{\rm high}
   &\le\Delta_B^\Gamma\le\mathcal H_k+r\Omega_k^{\rm low}.
\end{aligned}}\tag{AMT44}
\]
AMT29 displays every factor of both $\Omega$ values explicitly;
AMT32 gives the finite bounds
\[
\begin{aligned}
 0\le\Omega_k^{\rm low}
 &\le l\log[1+16(q-1+l)^2]+l\log[1+16(q+l)^2],\\
 0\le\Omega_k^{\rm high}
 &\le l\log[1+16(2q-1+l)^2]+l\log[1+16(2q+l)^2].
\end{aligned}\tag{AMT45}
\]
In particular both are $O(k\log(q+1))$ on the original
$m=1$ class. Its actual rank $r=8k-16$ therefore proves directly
\[
 \boxed{\Delta_K^\Gamma=O_h(k^2\log(q+1))=o(kq),\qquad
 \frac{F_K^{\rm ar}-F_K^\sigma}{kq}\longrightarrow0,\qquad
 \frac{F_K^{\rm ar}-F_K^0}{kq}\longrightarrow0.}
 \tag{AMT46}
\]
For the first equality divide the AMT44 bounds by $kq$ and use
$q=(k+1)^2$, so $k\log(q+1)/q\to0$. The second limit uses
AMT38, since $k^2/(kq)=k/q\to0$. The last uses the exact
identity $F_K^{\rm ar}-F_K^0=\Delta_K^\sigma+\Delta_K^\Gamma$.
All bounds are uniform for $|u|\ge R_*$.

For completeness the entire finite arithmetic boundary return is
also evaluated by those same quantities:
\[
\boxed{\begin{aligned}
 F_B^{\rm ar}-F_B^0
   &=\mathcal H_k+\delta_k^\sigma
                         -(F_K^{\rm ar}-F_K^0),\\
 \mathcal H_k+\delta_k^\sigma-r\widetilde E_k^-
   &\le F_B^{\rm ar}-F_B^0
             \le\mathcal H_k+\delta_k^\sigma+r\widetilde E_k^+.
\end{aligned}}\tag{AMT47}
\]
This uses AMT31 with its exact endpoint constants
$\widetilde E_k^+=E_k^++\Omega_k^{\rm low}$ and
$\widetilde E_k^-=E_k^-+\Omega_k^{\rm high}$, rather than a
new source comparison. Both tensor and arithmetic source masses
and the full original resultant have already been retained in
AMT23--30 before these signed cancellations.

The proved GEL18--19 constant is
\[
 \mathcal C_\Gamma=2\int_{a_*^2}^{b_*^2}\log x\,
                       \rho_{2,\pi}(x)\,dx-4(2\log2-1)>0,
 \qquad \frac{\mathfrak T_k}{lq}\longrightarrow\mathcal C_\Gamma.
 \tag{AMT48}
\]
Here $a_*,b_*$ and $\rho_{2,\pi}$ are exactly the proved GEL8/EIQ
equilibrium endpoints and density; their complete definitions and
integral proof remain in that provider. The endpoint letters here
avoid the already fixed original period parameter $u$. GEL18a
gives $4\log(27/(8\pi))\le\mathcal C_\Gamma\le6$ and proves
the strict positivity. TWA15 also proves the actual source return
$\delta_k^\sigma/(kq)\to0$. As $l/k\to1/4$, AMT44, AMT46
and AMT47 now prove
\[
\boxed{\begin{gathered}
 \frac{F_B^\sigma-F_B^0}{kq}\longrightarrow-\frac{\mathcal C_\Gamma}{4},
 \qquad
 \frac{F_B^{\rm ar}-F_B^0}{kq}\longrightarrow-\frac{\mathcal C_\Gamma}{4},
 \qquad
 \frac{F_B^{\rm ar}-F_B^\sigma}{kq}\longrightarrow0,\\
 \frac{F_B^\sigma-F_B^0}{\mathcal H_k}\longrightarrow1,
 \qquad
 \frac{F_B^{\rm ar}-F_B^0}{\mathcal H_k}\longrightarrow1,
 \quad m=1,\ k=4l+1,\ l\to\infty,\ |u|\ge R_*.
\end{gathered}}\tag{AMT49}
\]
Indeed the first two limits subtract an $o(kq)$ kernel return
from $\mathcal H_k=-\mathfrak T_k$, whose quotient by $kq$
tends to $-\mathcal C_\Gamma/4$. The third follows from
$\Delta_B^\sigma=\delta_k^\sigma-\Delta_K^\sigma$.
The nonzero limit of $\mathcal H_k/(kq)$ makes its denominator
nonzero for all sufficiently large original $k$ and proves the
last two ratios by division. The uniformity in $u$ follows from
the uniform kernel bounds and the period-independent total returns.
Thus the full existing source comparison evaluates the mixed
Gamma return also at its $kq$ scale, with its finite signed
remainders AMT44 and AMT47 still present.

% END COMPLETE AMT

\clearpage
% BEGIN COMPLETE KAF
% Complete original-source estimate; preceding editions are unchanged.
\section{Factorial control of the original common kernel at the \texorpdfstring{$kq$}{kq} scale}
\label{sec:KAF}

The exact denominator in OKV15 contains the information needed to improve
its absolute volume estimate. This proof keeps that denominator and also
allows every original primary multiplicity. The original Gamma sources,
root coordinates, coefficient maps and observation are retained. The
simple-quartet conclusion uses the actual period and rank calculations
PQO/RPC and OQC/OQX, with their complete proofs supplied as dependencies.

\subsection{Original coordinates, all primary multiplicities, and source norms}
Fix the original full-order quartet with $0<\delta<1/2$, $\gamma>2$,
and a positive integer primary order $m$. On the original tensor class set
\[
 k=4l+1,\quad l\ge1,\quad e_k=1+k(m-1),\quad
 q=e_k(k+1)^2,\quad c=k/2,\quad s\in\{1,k\},\quad b_s=s/2.
 \tag{KAF1}
\]
The original polynomial and its exact real-coordinate representative are
\[
\begin{split}
 \chi(S)&=\prod_{a,b=0}^k
 [S-c-(2a-k)\delta-i(2b-k)\gamma]^{e_k},\\
 \omega_{ab}&=(2b-k)\gamma-i(2a-k)\delta,\\
 P(y)&=i^{-q}\chi(c+iy)
       =\prod_{a,b=0}^k(y-\omega_{ab})^{e_k}.
\end{split}\tag{KAF2}
\]
In the full list $\omega_1,\ldots,\omega_q$, each root $\omega_{ab}$
occurs $e_k$ times. The exact maximal modulus remains
$R=k\sqrt{\delta^2+\gamma^2}$. Substitution gives inverse algebra maps
\[
 \Phi:\mathbb C[y]/(P)\longrightarrow E=\mathbb C[S]/(\chi),
 \quad[r]\longmapsto[r((S-c)/i)],\qquad
 \Phi^{-1}[f]=[f(c+iy)].\tag{KAF3}
\]
Indeed $P((S-c)/i)=i^{-q}\chi(S)$, and the two substitutions are
inverse before quotienting. Thus every root, multiplicity and factor of
$i$ in the original remainder operation is present.

Retain the complete measures and their monic orthogonal polynomials:
\[
\begin{gathered}
 M_s=(2\pi)^{s/2},\qquad
 dm_{0,s}(y)=\frac{M_s2^{b_s-2}}{\pi\Gamma(b_s)}
       \left|\Gamma\left(\frac{b_s}{2}+\frac{iy}{2}\right)\right|^2dy,\\
 p_0^{(s)}=1,\quad p_1^{(s)}=y,\qquad
 p_{d+1}^{(s)}=yp_d^{(s)}-d(b_s+d-1)p_{d-1}^{(s)},\\
 h_d^{(s)}=\int|p_d^{(s)}|^2dm_{0,s}=M_s d!(b_s)_d,\qquad
 \varphi_d^{(s)}=p_d^{(s)}/\sqrt{h_d^{(s)}},\qquad
 u_d^{(s)}(S)=\varphi_d^{(s)}((S-c)/i).
\end{gathered}\tag{KAF4}
\]
These are precisely the existing orders $1,k$. Their complete source
transform and norm calculation are IFB2--5 and OKV4--5. In particular
$\int e^{ty}dm_{0,s}=M_s(\cos t)^{-b_s}$ for $|t|<\pi/2$.
The generating series
$(1+z^2)^{-b_s/2}\exp(y\arctan z)$ has coefficients
$p_d^{(s)}z^d/d!$. Integrating its product at two sufficiently small
real variables gives $M_s(1-zw)^{-b_s}$, proving the displayed norms
and orthogonality by coefficient comparison. The original exponential
bound justifies those finite coefficient differentiations under the
integral. None of these source statements depends on $m$.

In the original polynomial frame of $E$ let $C_s$ have columns
$[u_0^{(s)}],\ldots,[u_{q-1}^{(s)}]$. Its determinant is exactly
\[
\begin{gathered}
 \det C_s=i^{-q(q-1)/2}\prod_{d=0}^{q-1}(h_d^{(s)})^{-1/2}\ne0,
 \\ v_j^{(s)}=C_s^{-1}[u_{q+j}^{(s)}],\quad
 K_j^{(s)}=I_q+\sum_{a=0}^{j-1}v_a^{(s)}(v_a^{(s)})^*,
 \quad0\le j\le q+1.
\end{gathered}\tag{KAF5}
\]
The last included vector has original degree $2q$.

\subsection{A degree-dependent constant in the unchanged remainder estimate}
Define the actual finite constants
\[
 A_{k,m}=\max\left\{2\sqrt{1+\frac{k}{2q}},\frac{R}{q}\right\},
 \qquad T=qA_{k,m},\qquad
 \mathcal A_{k,m}=6\log2+4+4\log A_{k,m}.
 \tag{KAF6}
\]
The letter $A_{k,m}$ in this section is a coefficient-estimate constant;
it is not the arithmetic upper source constant $A_k$ of TWA/AMT.
In particular $A_{k,m}\ge2$ and $R\le T$. For fixed original quartet
and fixed $m$, $q\ge(k+1)^2$ implies
\[
 A_{k,m}\longrightarrow2,\qquad
 \mathcal A_{k,m}\longrightarrow4+10\log2.
 \tag{KAF7}
\]
Indeed $k/q\to0$ and $R/q=\sqrt{\delta^2+\gamma^2}\,k/q\to0$.
The root-dependent term in the maximum is retained at every finite $k$.

Here is the complete remainder estimate needed below. For the full
root list define elementary and complete homogeneous polynomials
$e_h(\omega)$ and $h_t(\omega)$ by their finite sums, with
$e_0=h_0=1$ and $e_h=0$ outside $0\le h\le q$. The formal identity
\[
 \left(\sum_{h=0}^q(-1)^he_h(\omega)z^h\right)
 \left(\sum_{t\ge0}h_t(\omega)z^t\right)=1
 \tag{KAF8}
\]
follows by multiplying each of the $q$ factors $1-\omega_a z$
by its geometric series. Repeated roots do not affect this identity.
It proves, for $D=q+j$ and $0\le j\le q$, the exact signed formula
\[
 [y^a]\operatorname{rem}_P(y^{q+j})
 =-\sum_{t=0}^j h_t(\omega)(-1)^{q+j-t-a}e_{q+j-t-a}(\omega),
 \qquad0\le a<q.\tag{KAF9}
\]
In detail the quotient is $\sum_{t=0}^jh_t(\omega)y^{j-t}$;
KAF8 cancels all coefficients of degree at least $q$ in its product
with $P$, and the remaining coefficients are KAF9. This proves the
formula by uniqueness of monic division, including all multiplicities.

The defining sums give
$|e_h|\le\binom qh R^h$ and
$|h_t|\le\binom{q+t-1}{t}R^t$. In every nonzero summand of KAF9
the total root exponent is $q+j-a$. The binomial coefficients of
$h_t$ increase with $t$, and the distinct indices in the $e_h$ sum
contribute at most $2^q$. Hence
\[
 |[y^a]\operatorname{rem}_P(y^{q+j})|
 \le2^q\binom{q+j-1}{j}R^{q+j-a}
 \le2^{3q-1}R^{q+j-a}.\tag{KAF10}
\]
The last bound is the binomial expansion at one, since $j\le q$.

The exact orthonormal recurrence from KAF4 has upward and downward
weights $\sqrt{(a+1)(b_s+a)}$ and $\sqrt{a(b_s+a-1)}$.
On degree at most $q-2$ each shift has norm at most
$\sqrt{q(b_s+q)}$. Because $b_s\le k/2$, their sum has norm at
most $2\sqrt{q(k/2+q)}\le T$. Iterating from the constant, with all
intermediate degrees in that range, proves
\[
 \|y^a\|_{L^2(m_{0,s})}\le\sqrt{M_s}\,T^a,
 \qquad0\le a<q.\tag{KAF11}
\]
For the finite coefficient sum $c_T(f)=\sum_a|[y^a]f|T^a$ the
original recurrence gives
$c_T(p_{d+1}^{(s)})\le T c_T(p_d^{(s)})+
d(b_s+d-1)c_T(p_{d-1}^{(s)})$.
For $1\le d<2q$,
\[
 d(b_s+d-1)\le2q(k/2+2q)=4q^2+kq
       \le4q^2+2kq\le T^2.
\]
Starting with $c_T(p_0)=1,c_T(p_1)=T$, induction therefore proves
$c_T(p_d)\le(2T)^d$ for $0\le d\le2q$: the recurrence bound is
$(3/4)(2T)^{d+1}\le(2T)^{d+1}$.
Combining KAF10 with all $q$ monomial norms KAF11, then using the
linearity of the original remainder map, gives
\[
 \|\operatorname{rem}_P(p_d^{(s)})\|_{L^2(m_{0,s})}
       \le\sqrt{M_s}\,q2^{3q}(2T)^d,
       \qquad0\le d\le2q.\tag{KAF12}
\]
For a monomial of degree $D\ge q$, each coefficient term is at
most $\sqrt{M_s}2^{3q-1}R^{D-a}T^a\le
\sqrt{M_s}2^{3q-1}T^D$. For $D<q$ use KAF11 directly.
These two cases and the proved coefficient sum establish KAF12
without deleting a remainder coefficient or changing the source norm.

\subsection{Retaining the factorial denominator}
Expansion of the degree-below-$q$ remainder in the orthonormal
polynomials of KAF4, followed by the exact map KAF3, proves
\[
 \|v_j^{(s)}\|_2^2
 =\frac{\|\operatorname{rem}_P(p_{q+j}^{(s)})\|_{L^2(m_{0,s})}^2}
          {M_s(q+j)!(b_s)_{q+j}}
 \le\frac{q^22^{6q}(2T)^{2(q+j)}}{(q+j)!(b_s)_{q+j}}.
 \tag{KAF13}
\]
The factor $M_s$ occurs in the original numerator and denominator
before this cancellation. For every $d\ge0$,
\[
 d!(b_s)_d\ge d!(1/2)_d=\frac{(2d)!}{4^d},
 \qquad
 \|v_j^{(s)}\|_2^2\le q^22^{6q}\frac{(4T)^{2(q+j)}}{(2q+2j)!}.
 \tag{KAF14}
\]
The inequality is factor by factor because $b_s\ge1/2$.
The equality follows by separating even and odd factors of $(2d)!$.
This is precisely the denominator whose further retention sharpens
OKV16--17.

Define the positive exact quantities
\[
 \eta_{k,m}=\frac{1-1/(4q)}{A_{k,m}^2}<\frac14,\qquad
 Z_{k,m}=1+q^22^{6q}\frac{(4qA_{k,m})^{4q}}{(4q)!}
              \frac{1-\eta_{k,m}^{q+1}}{1-\eta_{k,m}}.
 \tag{KAF15}
\]
For $t_d=(4T)^{2d}/(2d)!$, $q+1\le d\le2q$ gives
$t_{d-1}/t_d=(2d)(2d-1)/(16T^2)\le\eta_{k,m}$.
Consequently $\sum_{d=q}^{2q}t_d\le
t_{2q}\sum_{a=0}^q\eta_{k,m}^a$. Adding the $q+1$ actual rank-one
updates in KAF5 and using KAF14 proves the finite operator bound
\[
 \boxed{1\le\|K_{q+1}^{(s)}\|_2\le Z_{k,m},\qquad s\in\{1,k\}.}
 \tag{KAF16}
\]
Here $\|vv^*\|_2=\|v\|_2^2$ follows by Cauchy--Schwarz with
equality on $v\ne0$; a zero vector contributes zero.

For an explicit linear bound, monotonicity of $\log x$ on each
interval $[j-1,j]$ gives
\[
 \log n!=\sum_{j=1}^n\log j
       \ge\int_1^n\log x\,dx=n\log n-n+1.
\]
Thus $(4qA_{k,m})^{4q}/(4q)!\le
\mathrm e^{-1}(\mathrm e A_{k,m})^{4q}$, where $\mathrm e$ is
the exponential constant and differs from the primary integer $e_k$.
The geometric sum in KAF15 is at most $4/3$. Therefore
$Z_{k,m}\le1+[4/(3\mathrm e)]X$, with
$X=q^2\exp(q\mathcal A_{k,m})$. The positive exponential series
gives $\mathrm e>1+1+1/2+1/6=8/3$, so $4/(3\mathrm e)<1/2$.
Since $A_{k,m}\ge2$ and $q\ge36$, $X>2$, and hence
$1+[4/(3\mathrm e)]X\le X$. We have proved
\[
 \boxed{\log\|K_{q+1}^{(s)}\|_2\le\log Z_{k,m}
       \le q\mathcal A_{k,m}+2\log q.}
 \tag{KAF17}
\]
The full factorial and finite geometric sum in KAF15 remain available
as the sharper finite bound.

\subsection{All four actual kernel and boundary Grams}
Let $\Lambda:E\twoheadrightarrow B$ be the original period and full-unit
observation, let $I:\mathcal K\hookrightarrow E$ be its actual kernel
inclusion, and put $r=\dim\mathcal K$. No rank formula is needed in
this subsection, so it applies to every original $m$. The original
quotient and restricted Grams are
\[
 G_{q+j-1}^{(s)}=C_s^{-*}(K_j^{(s)})^{-1}C_s^{-1},\qquad
 T_s=C_s^{-1}I,\qquad H_j^{(s)}=T_s^*(K_j^{(s)})^{-1}T_s.
 \tag{KAF18}
\]
Indeed the original orthonormal source-coordinate remainder matrix is
$[I_q,v_0,\ldots,v_{j-1}]$. Multiplication by its adjoint gives $K_j$;
its unique minimum lift of $x$ is its adjoint applied to $K_j^{-1}x$,
with squared norm $x^*K_j^{-1}x$. This proves both matrices KAF18
in the original coefficient frame. The positive update chain implies
\[
 Z_{k,m}^{-1}T_s^*T_s\preceq H_j^{(s)}\preceq T_s^*T_s,
 \qquad H_{j+1}^{(s)}\preceq H_j^{(s)}.\tag{KAF19}
\]
Inversion reverses positive order by congruence and inversion of
positive eigenvalues. Restricting this order with the same injective
$T_s$ proves KAF19. After congruence with $(T_s^*T_s)^{-1/2}$,
every eigenvalue lies between $Z_{k,m}^{-1}$ and one. Thus
$-r\log Z_{k,m}\le\log\det H_j^{(s)}-
\log\det(T_s^*T_s)\le0$.

The exact four original degrees are $q-1,q,2q-1,2q$. With their
unchanged signs define
\[
 F_K^{(s)}=\log\det H_0^{(s)}+\log\det H_1^{(s)}
              -\log\det H_q^{(s)}-\log\det H_{q+1}^{(s)}.
\]
Pairing degrees $q-1$ with $2q-1$, and $q$ with $2q$, proves
nonnegativity by KAF19. Each difference is at most $r\log Z_{k,m}$.
The four original frame factors cancel only after these four signs
have been recorded. It follows that
\[
 \boxed{0\le F_K^{(s)}\le2r\log Z_{k,m}
       \le2r\,[q\mathcal A_{k,m}+2\log q],\qquad s\in\{1,k\}.}
 \tag{KAF20}
\]
If $r=0$, every kernel determinant is one and the statement has its
same empty-space meaning.

At each original degree the boundary Gram is
$Q_j^{(s)}=(\Lambda(G_{q+j-1}^{(s)})^{-1}\Lambda^*)^{-1}$.
For any fixed section $S_*$ of $\Lambda$, its minimum section
$S_j^{(s)}=(G_{q+j-1}^{(s)})^{-1}\Lambda^*Q_j^{(s)}$ differs from
$S_*$ by a kernel column and is orthogonal to $I$ in that metric.
The change of frame has determinant one. Taking determinants in this
orthogonal frame proves, with its original coefficient factor,
\[
 \det H_j^{(s)}\det Q_j^{(s)}
       =|\det[I,S_*]|^2\det G_{q+j-1}^{(s)}.
\]
The same four signs now prove the exact splitting
$F_K^{(s)}+F_B^{(s)}=\mathcal B_k^{(s)}$. Boundary minimum norms
decrease with degree, so $F_B^{(s)}\ge0$. In particular
\[
 \boxed{\max\{0,\mathcal B_k^{(s)}-2r\log Z_{k,m}\}
       \le F_B^{(s)}\le\mathcal B_k^{(s)}.}\tag{KAF21}
\]
These are statements about the original kernel and boundary, with
all multiplicities and the same observation throughout.

\subsection{The actual simple quartet and the arithmetic source}
For $m=1$, the complete RPC calculation supplies an explicit finite
$R_*$ from the original unit and quartet. On every logarithm branch
and for $|u|\ge R_*$ its original period quadric has five distinct
cyclic orbit members. The full invariant restriction ideal calculation
OQX then gives $r=8k-16$. Substitution in the finite KAF20 gives
\[
 \boxed{0\le\frac{F_K^{(s)}}{kq}
 \le\left(16-\frac{32}{k}\right)
       \left(\mathcal A_{k,1}+\frac{2\log q}{q}\right),
 \qquad s\in\{1,k\},\quad |u|\ge R_*.}\tag{KAF22}
\]
The original arithmetic-source comparison AMT17--22 proves
$-rE_k^+\le F_K^{\rm ar}-F_K^{(1)}\le rE_k^-$, where the
unchanged full constants are
\[
\begin{gathered}
 E_k^+=2\log(A_k/a_k)+\log D_{q-1}+\log D_q,\\
 E_k^-=2\log(A_k/a_k)+\log D_{2q-1}+\log D_{2q},
 \\ D_N=[1+4(N+21+4m)^2]^{21+4m}.
\end{gathered}
\]
Here $a_k,A_k$ are exactly the original arithmetic constants of AMT4,
whose full source inequality and minimum-section proof are included
there. They retain the actual numerator $2\xi/h$ and all source
masses; AMT22 proves $E_k^\pm=O_h(k+\log q)$.
Thus the actual arithmetic metric satisfies the finite bound
\[
 \boxed{0\le F_K^{\rm ar}
       \le(16k-32)\log Z_{k,1}+(8k-16)E_k^-.}\tag{KAF23}
\]
Combining KAF7, KAF22--23, and $q=(k+1)^2$ proves uniformly on
the entire stated original period domain, for
$a\in\{\sigma,0,\mathrm{ar}\}$,
\[
 \boxed{0\le\liminf_{k\to\infty}\frac{F_K^a}{kq}
       \le\limsup_{k\to\infty}\frac{F_K^a}{kq}
       \le16(4+10\log2).}\tag{KAF24}
\]
For precision, uniformity means the same finite upper bound holds for
every such $u$ and branch at each original $k$; therefore it also
bounds every sequence of those parameters. The arithmetic extra term
divided by $kq$ is at most $(8-16/k)E_k^-/q\to0$.
AMT46 further proves that the three ratios differ by quantities
tending uniformly to zero. Thus their limit points coincide for any
fixed allowed sequence of period parameters; KAF24 does not assert
that this bounded sequence already has a unique limit.

Finally the exact arithmetic identity gives
\[
 \max\{0,\mathcal B_k^{\rm ar}-(16k-32)\log Z_{k,1}
                              -(8k-16)E_k^-\}
       \le F_B^{\rm ar}\le\mathcal B_k^{\rm ar}.\tag{KAF25}
\]
Its full-source value remains
$\mathcal B_k^{\rm ar}=\mathcal B_k^\sigma+\delta_k^\sigma
=\mathcal B_k^0+\mathcal H_k+\delta_k^\sigma$.
All exact Gamma and arithmetic signed returns of AMT44--49 are
retained. The new calculation improves the absolute common kernel
bound to the $kq$ scale, while its exact four-endpoint determinant
and actual period dependence remain the objects for the next leading
term calculation.

% END COMPLETE KAF

\clearpage
% BEGIN COMPLETE OCF
% Original coefficient reduction. All maps are complex-linear until a metric is supplied.
\section{An exact conductor frame and a fixed invariant-generator recursion}
\label{sec:OCF}

Retain the original simple-quartet data, $m=1$, $k=4l+1$,
$l\geq1$, the original root order $(++),(+-),(-+),(--)$, and
the complete original period and amplitude matrix of OQC1:
\[
                  \mathsf H=F^{-1}\Pi U\mathcal E_{\rm prim}.
\]
Thus $\det\mathsf H\ne0$ and
$Q(x)=Q_0(\mathsf H^{-1}x)$, where
$Q_0(t)=t_{++}t_{--}-t_{+-}t_{-+}$.
Use $g f(x)=f(D^{-1}x)$,
$D=\operatorname{diag}(\zeta,\zeta^2,\zeta^3,\zeta^4)$,
$\zeta=e^{2\pi i/5}$, and the original invariant algebra
$\mathscr I=\mathbb C[x_1,x_2,x_3,x_4]^{\langle g\rangle}$.
This calculation applies on the actual five-orbit period locus,
which includes every original logarithm branch with $|u|\geq R_*$
by PQO and RPC. On that locus the five prime polynomials
$Q_a=g^aQ$, $0\leq a<5$, are pairwise nonassociate. Keep their
actual products
\[
 \mathcal A=\prod_{a=1}^4Q_a,\qquad
 \mathcal H_Q=Q\mathcal A=\prod_{a=0}^4Q_a,
 \qquad \deg\mathcal A=8,\quad\deg\mathcal H_Q=10.
 \tag{OCF1}
\]
No coefficient in these polynomials is replaced. The exact orbit
decision and the period domain are those of the cited original
calculations. The construction below treats every possible nonzero
leading coefficient of their actual restriction.

\subsection{The original Segre coefficient isomorphism}

For $d\geq0$ let $\mathcal B_d$ be the space of biforms of
bidegree $(d,d)$ in $(z_0,z_1)$ and $(w_0,w_1)$. Its ordered basis is
\[
 e_{ab}^{(d)}=z_0^a z_1^{d-a}w_0^b w_1^{d-b},
 \qquad 0\leq a,b\leq d.
 \tag{OCF2}
\]
Pairs are ordered lexicographically, first by $a$, then by $b$.
The order is used only for the exact coefficient elimination.
Set $\mathcal B_d=0$ for $d<0$. Multiplication satisfies
$e_{ab}^{(d)}e_{ij}^{(e)}=e_{a+i,b+j}^{(d+e)}$.
The original substitution is
\[
 \begin{gathered}
 t=(z_0w_0,z_0w_1,z_1w_0,z_1w_1)^T,\qquad x=\mathsf Ht,\\
 \mathscr P_d=\mathbb C[x_1,x_2,x_3,x_4]_d,\qquad
 \psi_d:\mathscr P_d\longrightarrow\mathcal B_d,\quad
 f\longmapsto f(\mathsf Ht).
 \end{gathered}\tag{OCF3}
\]
The induced map
$R_d=\mathscr P_d/(Q\mathscr P_{d-2})\longrightarrow\mathcal B_d$
is an isomorphism. Here is a coefficient proof with the given
orientation. For each $(a,b)$, choose an integer
$h$ between $\max(0,a+b-d)$ and $\min(a,b)$. The monomial
$t_{++}^{h}t_{+-}^{a-h}t_{-+}^{b-h}t_{--}^{d-a-b+h}$ maps
to $e_{ab}^{(d)}$, proving surjectivity. Two monomials with
the same pair differ in their exponent vectors by an integer
multiple of $(1,-1,-1,1)$. After removing their common monomial,
their difference is divisible by $t_{++}t_{--}-t_{+-}t_{-+}$,
using $A^n-B^n=(A-B)\sum_{j=0}^{n-1}A^{n-1-j}B^j$.
A polynomial in the kernel has coefficient sum zero within every
pair, so it is a sum of such differences. The reverse inclusion
in the kernel follows by direct substitution. The invertible
substitution $x=\mathsf Ht$ then gives exactly the asserted kernel.
Thus these maps preserve multiplication as well as every grading.

Write the full conductor restriction as
\[
 A=\psi_8(\mathcal A)
     =\sum_{r=0}^8\sum_{s=0}^8 a_{rs}e_{rs}^{(8)}\ne0.
 \tag{OCF4}
\]
To verify its nonvanishing, $Q$ is a rank-four irreducible quadric,
hence prime in the polynomial ring, and none of $Q_1,\ldots,Q_4$
is its associate. Therefore $Q$ does not divide their product.
The kernel calculation OCF3 proves OCF4. The ring
$\mathcal B=\bigoplus_{d\geq0}\mathcal B_d$, with the indicated
diagonal grading, embeds in the four-variable polynomial ring;
it is an integral domain. Multiplication by $A$ is consequently
injective on every $\mathcal B_d$.

\subsection{Every leading-coefficient case gives an exact strip frame}

Let $(r_*,s_*)$ be the lexicographically largest pair for which
$a_{r_*s_*}\ne0$, and put $a_*=a_{r_*s_*}$. This is an
actual nonzero coefficient of OCF4. For $d\geq8$ define
\[
 \begin{split}
 P_d&=\{(r_*+i,s_*+j):0\leq i,j\leq d-8\},\\
 E_d^{\rm strip}&=\{0,\ldots,d\}^2\setminus P_d,\qquad
 W_d=\mathbb C^{E_d^{\rm strip}}.
 \end{split}\tag{OCF5}
\]
For $0\leq d<8$ put $P_d=\varnothing$,
$E_d^{\rm strip}=\{0,\ldots,d\}^2$ and use the same definition
of $W_d$; for $d<0$ put $W_d=0$. Let
$S_d:W_d\hookrightarrow\mathbb C^{(d+1)^2}$ insert zeros in
the positions $P_d$. All coefficient columns below use OCF2.

The following finite elimination defines
$N_d:\mathbb C^{(d+1)^2}\to W_d$ and, for $d\geq8$,
$V_d:\mathbb C^{(d+1)^2}\to\mathbb C^{(d-7)^2}$.
Starting with the complete input biform $f$, traverse all
$(i,j)\in\{0,\ldots,d-8\}^2$ in decreasing lexicographic order.
At the step $(i,j)$, let $c$ be the current coefficient at
$(r_*+i,s_*+j)$, record $b_{ij}=c/a_*$, and replace the current
biform by
\[
       f_{\rm current}-b_{ij}A e_{ij}^{(d-8)}.
 \tag{OCF6}
\]
At completion take the coefficients on $E_d^{\rm strip}$ as
$N_df$ and take $(b_{ij})$ as $V_df$. If $d<8$, $N_d$ is the
identity and $V_d$ is the zero map into the zero space.
The exact multiplication matrix $M_{A,d}$ gives
\[
 S_dN_d+M_{A,d}V_d=I,\qquad
 N_dS_d=I_{W_d},\qquad
 \ker N_d=A\mathcal B_{d-8}.
 \tag{OCF7}
\]
To prove these identities, the largest term of
$A e_{ij}^{(d-8)}$ is $a_*e_{r_*+i,s_*+j}^{(d)}$;
every other term is strictly smaller in the same lexicographic
order. Therefore each step zeros its designated position without
altering an earlier zero. Summing all subtractions gives the
first identity. Input supported on the strip produces zero
coefficients at every elimination step, giving the second.
Finally, the leading term of $Ab$, for a nonzero homogeneous
biform $b$, is the product of the two leading terms and belongs
to $P_d$. Indeed every other product is smaller in the additive
lexicographic order. Thus
$A\mathcal B_{d-8}\cap S_dW_d=0$. The first identity now proves
the kernel formula, the uniqueness of the decomposition, and all
three identities in OCF7. All operations are linear in the input;
their coefficients are explicit rational expressions in the
original $a_{rs}$ with powers of the displayed $a_*$ as denominators.
The polynomial $A$ itself remains the full polynomial OCF4.

In particular the quotient map induced by $N_d$ is the exact
isomorphism
\[
 \mathcal B_d/(A\mathcal B_{d-8})\xrightarrow{\ \sim\ }W_d,
 \qquad
 n_d:=\dim W_d=
 \begin{cases}(d+1)^2,&0\leq d<8,\\16d-48,&d\geq8.
 \end{cases}
 \tag{OCF8}
\]
The second formula is
$(d+1)^2-(d-7)^2=16d-48$. The possible choices
$(r_*,s_*)\in\{0,\ldots,8\}^2$ form an exhaustive finite
partition: the chart for a pair specifies that its coefficient
is nonzero and every larger coefficient is zero. Hence all actual
coefficient cases have been included. At the first original
degree $k=5$ this construction retains all $36$ coordinates;
at every original $k\geq9$ it retains $16k-48$ coordinates.

\subsection{Twenty-five fixed invariant monomials generate all degrees}

Define the following fixed finite sets of exponent vectors:
\[
 G_e=\{\nu\in\mathbb N^4:|\nu|=e,
                   \nu_1+2\nu_2+3\nu_3+4\nu_4\equiv0\pmod5\},
 \quad 1\leq e\leq5,\qquad G=\bigcup_{e=1}^5G_e.
 \tag{OCF9}
\]
Here $G_1$ is empty. The complete remaining lists, in the original
$x_1,x_2,x_3,x_4$ order, are
\[
\begin{array}{c|l}
e&G_e\\\hline
2&(0,1,1,0),(1,0,0,1)\\
3&(0,0,2,1),(0,1,0,2),(1,2,0,0),(2,0,1,0)\\
4&(0,0,1,3),(0,2,2,0),(0,3,0,1),(1,0,3,0),\\
 & (1,1,1,1),(2,0,0,2),(3,1,0,0)\\
5&(0,0,0,5),(0,0,5,0),(0,1,3,1),(0,2,1,2),\\
 & (0,5,0,0),(1,0,2,2),(1,1,0,3),(1,3,1,0),\\
 & (2,1,2,0),(2,2,0,1),(3,0,1,1),(5,0,0,0).
\end{array}\tag{OCF10}
\]
Thus $(|G_1|,\ldots,|G_5|)=(0,2,4,7,12)$ and $|G|=25$.
These lists can also be obtained directly by choosing
$0\leq\nu_1,\nu_2,\nu_3\leq e$ with their sum at most $e$,
putting $\nu_4=e-\nu_1-\nu_2-\nu_3$, and applying the single
displayed congruence, which verifies completeness of each list.

Every positive-degree invariant monomial is a product of members
$x^\nu$, $\nu\in G$. To prove this, list the weights of its
individual factors in any order. If its degree is at most five,
the entire list gives a member of $G$. If its degree exceeds five,
consider the six partial sums of the first five weights, including
the empty partial sum, in $\mathbb Z/5\mathbb Z$. Two coincide;
the nonempty block between them has length at most five and total
weight zero. Its factors give a divisor $x^\nu$ with $\nu\in G$.
The complementary monomial again has total weight zero and smaller
degree. Induction completes the factorization. Because the action
is diagonal, a polynomial is invariant exactly when all its
nonzero monomials have weight zero: comparison of each monomial
coefficient in $gf=f$ proves this statement. Thus
\[
                \mathscr I=\mathbb C[x^\nu:\nu\in G].
 \tag{OCF11}
\]
All intermediate degrees here are degrees in the original
coefficient algebra. This construction introduces no new source
distribution and leaves the original source orders $1$ and $k$ intact.

For each $\nu\in G_e$ the complete generator biform is
\[
 f_\nu=\psi_e(x^\nu)
  =\prod_{j=1}^4
   (H_{j1}z_0w_0+H_{j2}z_0w_1+
                  H_{j3}z_1w_0+H_{j4}z_1w_1)^{\nu_j}
  =\sum_{a,b=0}^e c_{\nu,ab}e_{ab}^{(e)}.
 \tag{OCF12}
\]
For an entirely explicit coefficient formula, sum over all
nonnegative $4$ by $4$ integer arrays $(n_{jt})$ satisfying
$\sum_t n_{jt}=\nu_j$,
$\sum_j(n_{j1}+n_{j2})=a$ and
$\sum_j(n_{j1}+n_{j3})=b$. Then
\[
 c_{\nu,ab}=
 \sum_{(n_{jt})}
     \prod_{j=1}^4\left(
        \frac{\nu_j!}{\prod_{t=1}^4n_{jt}!}
        \prod_{t=1}^4H_{jt}^{n_{jt}}\right).
 \tag{OCF13}
\]
The multinomial theorem proves the formula, including every
original period and amplitude entry. Every factor has degree at
most five and at most $36$ possible biform coefficients.

\subsection{Exact multiplication on the strip and the invariant recursion}

Let $\mathcal L=\mathcal B/(A\mathcal B)$, a graded algebra.
For $\nu\in G_e$ and $d\geq e$ define
\[
 m_{\nu,d}:W_{d-e}\longrightarrow W_d,
 \qquad m_{\nu,d}=N_d M_{f_\nu}S_{d-e}.
 \tag{OCF14}
\]
Here $M_{f_\nu}$ is the exact coefficient convolution:
\[
       (M_{f_\nu}S_{d-e}v)_{ab}
          =\sum_{i,j}c_{\nu,a-i,b-j}(S_{d-e}v)_{ij},
\]
with out-of-range indices giving zero.
Therefore every entry of OCF14 is supplied
by the finite formulas OCF6 and OCF13.
For a representative $h\in\mathcal B_{d-e}$, OCF7 writes
$h-S_{d-e}N_{d-e}h=A b$. Multiplication by $f_\nu$ preserves
this ideal. Applying $N_d$ yields the exact identity
\[
       N_d(f_\nu h)=m_{\nu,d}N_{d-e}h.
 \tag{OCF15}
\]
This proves that OCF14 is multiplication by the specified element
of $\mathcal L$, including its independence of representatives.

The original invariant image is
$J_d=\psi_d(\mathscr I_d)\subseteq\mathcal B_d$.
Set $\overline J_d=N_dJ_d\subseteq W_d$, and set
$\overline J_d=0$ for $d<0$. The exact recursion is
\[
 \overline J_0=\mathbb C\cdot1,\qquad
 \overline J_d=
   \sum_{\substack{\nu\in G_e\\2\leq e\leq5,\ e\leq d}}
                 m_{\nu,d}\overline J_{d-e}
                 \quad(d\geq1).
 \tag{OCF16}
\]
Indeed OCF11 factors every invariant monomial of degree $d>0$
as $x^\nu$ times an invariant monomial of degree $d-e$.
Applying $\psi$, then OCF15, gives the inclusion from left to
right. Conversely every summand is the image of a product of two
invariant polynomials, so it belongs to the left side. Linearity
proves equality for the complete invariant space. In particular
the recursion gives $\overline J_1=0$ and includes all small degrees.

The conductor ideal is contained in this invariant image with
its original trace coefficient:
\[
 A\mathcal B_{d-8}\subseteq J_d,\qquad
 [5\mathsf R(\mathcal A h)]_{(Q)}=[\mathcal A h]_{(Q)},
 \quad\mathsf R=\frac15\sum_{a=0}^4g^a.
 \tag{OCF17}
\]
To prove the identity, $g^a\mathcal A$ is the product of all
five orbit equations except $Q_a$; for $a\ne0$ it contains $Q$.
Thus all four corresponding terms vanish modulo $Q$, and the
$a=0$ term is the displayed product. Every biform in
$\mathcal B_{d-8}$ has a polynomial lift by OCF3, so the identity
proves the inclusion. It also shows exactly why the quotient
used in OCF16 retains the original invariant image.

\subsection{Exact ranks, a finite pivot atlas, and the original quotient map}

For $d\geq0$ put
\[
 d_0(d)=\frac{\binom{d+3}{3}+4\epsilon_d}{5},\qquad
 \epsilon_d=\mathbf1_{d\equiv0\ (5)}-\mathbf1_{d\equiv1\ (5)},
 \qquad d_0(d)=0\quad(d<0).
 \tag{OCF18}
\]
This is the number of degree-$d$ weight-zero monomials: the
identity character has generating function $(1-z)^{-4}$, each
of the four nonidentity characters has generating function
$\prod_{j=1}^4(1-z\zeta^j)^{-1}=(1-z)/(1-z^5)$, and averaging
their coefficients gives OCF18.
Every invariant polynomial divisible by $Q$ is divisible by all
five $Q_a$, hence by $\mathcal H_Q$, since they are pairwise
nonassociate primes. Its quotient by this product is invariant:
apply $g$ and cancel the nonzero invariant product in the integral
domain. Conversely each $\mathcal H_Q$ times an invariant
polynomial is invariant and vanishes modulo $Q$. Consequently
$\ker(\mathscr I_d\to J_d)=\mathcal H_Q\mathscr I_{d-10}$,
and injection of multiplication by the nonzero product gives
$\dim J_d=d_0(d)-d_0(d-10)$.
Injection of multiplication by $A$, OCF7 and OCF17 then prove
\[
 j_d:=\dim\overline J_d
     =d_0(d)-d_0(d-10)-\max(d-7,0)^2.
 \tag{OCF19}
\]
In particular $j_0=1$, $j_1=0$, $j_5=12$, and
$j_d=8d-32$ for $d\geq8$. For the last formula at $d\geq10$,
$\epsilon_d=\epsilon_{d-10}$ and direct expansion gives
$\binom{d+3}{3}-\binom{d-7}{3}=5(d^2-6d+17)$;
subtracting $(d-7)^2$ gives $8d-32$. At $d=8,9$ the ranks
before this subtraction are $33,44$, respectively, which gives
$32,40$ and verifies both remaining cases.

Here is a complete finite construction of all the columns used
by OCF16. Start with $Z_0=(1)$ and the invariant polynomial
certificate $F_{0,1}=1$. Suppose all earlier matrices $Z_b$ have
independent columns spanning $\overline J_b$, and their columns
are $N_b\psi_b(F_{b,h})$ for invariant polynomials $F_{b,h}$.
Form the candidate matrix
\[
 C_d=\bigl[\,m_{\nu,d}Z_{d-e}\,\bigr]_{
                    \nu\in G_e,\ 2\leq e\leq5,\ e\leq d}.
 \tag{OCF20}
\]
Each candidate column has the explicit invariant certificate
$x^\nu F_{d-e,h}$. OCF16 proves that $C_d$ has column image
$\overline J_d$ and rank $j_d$. Select $j_d$ columns and $j_d$
rows giving a nonzero determinant. Let $Z_d$ be those columns,
let $I_d$ be those row positions, and let $T_d$ be the complementary
row positions in $E_d^{\rm strip}$. Retain the corresponding
invariant certificates. This completes the induction. When
$j_d=0$, use the empty matrix with zero columns, empty $I_d$,
and $T_d=E_d^{\rm strip}$.

Every selected determinant and every matrix entry is an explicit
expression obtained from the original entries of $\mathsf H$ and
the full $a_{rs}$. To cover all cases, order the finite choices
of column and row subsets and select the first nonzero determinant;
the corresponding chart asserts its nonvanishing and the vanishing
of all preceding candidates. A nonzero determinant exists because
the rank was proved in OCF19. At rank zero no determinant is needed.
Together with the exhaustive leading-coefficient charts after
OCF8, this gives a finite atlas for each specified $k$. It requires
no generic coefficient premise and supplies no numerical zero test
for an unevaluated analytic coefficient: each chart is specified
by its exact minors.

Define $\rho_d:W_d\to\mathbb C^{T_d}$ and the full original
coefficient map $\pi_d$ by
\[
 \rho_d=\operatorname{row}_{T_d}
            -(Z_d)_{T_d}(Z_d)_{I_d}^{-1}
                                  \operatorname{row}_{I_d},
 \qquad \pi_d=\rho_dN_d.
 \tag{OCF21}
\]
At rank zero the second term is the zero map. If $E_{T_d}$
inserts a column into its $T_d$ positions in $W_d$, then
\[
 \begin{split}
 w&=Z_d(Z_d)_{I_d}^{-1}w_{I_d}+E_{T_d}\rho_dw,\\
 \rho_dE_{T_d}&=I,\qquad\ker\rho_d=\operatorname{im}Z_d
                                  =\overline J_d,\\
 \ker\pi_d&=J_d,\qquad
       \pi_d S_d E_{T_d}=I_{\mathbb C^{T_d}}.
 \end{split}\tag{OCF22}
\]
For the first line, the subtraction on its right leaves zero
$I_d$ rows and exactly $\rho_dw$ on the complementary rows.
This proves the first three claims. If $\pi_df=0$, then
$N_df\in\overline J_d$; choose $j\in J_d$ with $N_dj=N_df$.
Now $f-j\in\ker N_d=A\mathcal B_{d-8}\subseteq J_d$ by
OCF17, so $f\in J_d$. The reverse inclusion follows immediately
from $N_dJ_d=\overline J_d$. The last equality uses $N_dS_d=I$.
This proves every map and kernel in OCF22.

At the original degrees, these are the exact dimensions
\[
 \begin{array}{c|c|c|c}
 k&n_k&j_k&|T_k|=n_k-j_k\\\hline
 5&36&12&24\\
 k\geq9,\ k=4l+1&16k-48&8k-32&8k-16.
 \end{array}\tag{OCF23}
\]
The original map $\Theta_k:R_k\xrightarrow{\sim}E_k^\vee$
in OQC7 sends the coefficient basis OCF2 to the original dual
primary basis in the same order. Indeed its evaluation on
$e^{wY_k}[1]$ is
$\exp(((2b-k)\gamma-i(2a-k)\delta)w)$. The corresponding
eigenvalue of $A_k$ is
$\beta_{ab}=k/2+(2a-k)\delta+i(2b-k)\gamma$; the eigenvalue
of $Y_k=(A_k-kI/2)/i$ is
$\omega_{ab}=(\beta_{ab}-k/2)/i=(2b-k)\gamma-i(2a-k)\delta$.
These eigenvalues are distinct because
$\delta>0$ and $\gamma>0$. Their exponentials are independent
by the Vandermonde determinant, so the evaluation identifies the
covectors exactly. By the original observation identity OQC8,
$\Theta_kJ_k=\operatorname{im}\Lambda_k^\vee$. Therefore OCF22
supplies the specified isomorphism
\[
 E_k^\vee/\operatorname{im}\Lambda_k^\vee
       \xrightarrow{\ \pi_k\ }\mathbb C^{T_k}
       \cong (\ker\Lambda_k)^\vee.
 \tag{OCF24}
\]
The displayed $\pi_k$ uses exactly the original primary coefficient
coordinates under $\Theta_k$. Every dualization in this statement
is complex-linear. No conjugation is inserted in this coefficient
map; the original metric maps can subsequently be composed with it.

Finally, this construction is finite with explicitly bounded
compressed matrices. OCF20 has $n_d$ rows and
$\sum_{e=2}^5 |G_e|j_{d-e}$ columns, with negative subscripts zero.
Both dimensions grow at most linearly in $d$: for example
$n_d\leq16d+1$ and the column number is at most $25(16d+1)$.
For $d\geq13$ the exact column number is $200d-1632$, obtained
from $\sum_e|G_e|=25$, $\sum_e e|G_e|=104$, and
$j_{d-e}=8(d-e)-32$. Only the preceding five graded bases are
needed to form the next matrix. Each column is obtained from
the at-most-$36$-coefficient formula OCF12, one exact convolution,
and OCF6. A full original biform coefficient column, if used as
temporary storage, has $(d+1)^2$ entries; every invariant relation
matrix in OCF20 has the stated linear row and column counts.
The invariant certificates can be retained as their products and
parent-column references, so their expanded high-degree monomials
need never be enumerated. The proof of their completeness is
OCF11--16, and the exact original quotient and its right inverse
are OCF21--24.

\subsection{The same recursion with fourteen proved generators}

The complete fixed set OCF9 can be reduced to the following
fourteen specific members while preserving every map and image
just proved:
\[
\begin{array}{c|l}
e&\mathcal G_e\\\hline
2&x_1x_4,\ x_2x_3\\
3&x_1^2x_3,\ x_1x_2^2,\ x_2x_4^2,\ x_3^2x_4\\
4&x_1^3x_2,\ x_1x_3^3,\ x_2^3x_4,\ x_3x_4^3\\
5&x_1^5,\ x_2^5,\ x_3^5,\ x_4^5 .
\end{array}\tag{OCF25}
\]
Here is a complete proof. Call a positive-degree invariant
monomial indecomposable when it is not the product of two
positive-degree invariant monomials. The divisor proof preceding
OCF11 shows that each indecomposable has degree at most five.
If it contains both $x_1,x_4$, it is their product, since otherwise
division by this invariant quadratic leaves a nonconstant
invariant monomial. The same assertion holds for $x_2,x_3$.
After these cases are removed, at most two variables occur:
every three indices contain one of the two opposite pairs.
One variable gives precisely its fifth power. The remaining
two-element supports are $\{1,2\},\{1,3\},\{2,4\},\{3,4\}$.
For positive exponents $(a,b)$ with $a+b\leq5$, the equation
$a+2b\equiv0\pmod5$ has exactly $(a,b)=(3,1),(1,2)$, and
$a+3b\equiv0\pmod5$ has exactly $(a,b)=(2,1),(1,3)$.
For $\{2,4\}$ divide the weights by $2$ in the field
$\mathbb Z/5\mathbb Z$, giving the first list; for $\{3,4\}$
divide by $3$, giving the second. These are exactly the eight
mixed monomials of degrees three and four in OCF25. Every
listed monomial has no proper invariant divisor, as is also
seen from these lists and the excluded opposite pairs.
This proves the classification. Factoring any invariant
monomial repeatedly into indecomposables terminates because
positive total degree strictly decreases. Hence OCF25 generates
the entire invariant algebra.

Use $\mathcal G=\bigcup_{e=2}^5\mathcal G_e$ in place of
$\{x^\nu:\nu\in G\}$ in OCF12--16 and OCF20. The same
factorization argument proves the same exact invariant image;
all coefficient formulas, invariant certificates, pivot charts,
quotient maps and ranks remain valid. The new recursion has
$2,4,4,4$ generators in degrees $2,3,4,5$, respectively. Thus
for $d\geq13$ its precise matrix dimensions and rank are
\[
 C_d^{(14)}\in
 \operatorname{Mat}_{(16d-48)\,\times\,(112d-864)}(\mathbb C),
 \qquad \operatorname{rank}C_d^{(14)}=8d-32.
 \tag{OCF26}
\]
Indeed the number of generators is $14$ and the sum of their
degrees is $2\cdot2+4\cdot3+4\cdot4+4\cdot5=52$; substituting
the already proved $j_{d-e}=8(d-e)-32$ gives
$14(8d-32)-8\cdot52=112d-864$. At smaller degrees the exact
number of columns is the same sum of the appropriate $j_{d-e}$
from OCF19, with negative subscripts zero. All coefficients
in this reduced recursion are still the original coefficients
of OCF12--14.

% END COMPLETE OCF

\clearpage
% BEGIN COMPLETE OKM
% Complete independent metric/CD provider. Original source orders only.
\section{The exact common kernel frame and its original Gamma metrics}
\label{sec:OKM}

Retain the original simple quartet $m=1$, $k=4l+1$, $l\geq1$,
$q=(k+1)^2$, $c=k/2$, $0<\delta<1/2$ and $\gamma>2$.
The original period parameter and its branch are fixed, on the proved
original OCF rank domain $|u|\geq R_*$, with $R_*$ the retained
explicit period-orbit radius. The kernel
below is the kernel of that one fixed original observation
$\Lambda:E\to B$. It is common to the two source metrics and to
their four degree endpoints. No intersection over different period
parameters or different observations is made.

Order all pairs $(a,b)$, $0\leq a,b\leq k$, lexicographically once
and retain that order. Define
\[
 \begin{gathered}
 \beta_{ab}=c+(2a-k)\delta+i(2b-k)\gamma,\qquad
 \omega_{ab}=\frac{\beta_{ab}-c}{i}
                =(2b-k)\gamma-i(2a-k)\delta,\\
 \chi(S)=\prod_{a,b=0}^k(S-\beta_{ab}),\qquad
 E=\mathbb C[S]/(\chi),\qquad
 \varepsilon_{ab}=
 \left[\prod_{(a',b')\ne(a,b)}
       \frac{S-\beta_{a'b'}}{\beta_{ab}-\beta_{a'b'}}\right],\\
 V_\beta=((\beta_{ab})^d)_{(a,b),\,0\leq d<q}.
 \end{gathered}\tag{OKM1}
\]
The points are distinct: equality of their real and imaginary parts
forces both indices equal. Hence every denominator displayed above
is nonzero. Evaluation at all points gives
$\varepsilon_{ab}(\beta_{a'b'})=\mathbf1_{(a,b)=(a',b')}$.
Polynomial interpolation proves that these are a basis of $E$, and
the evaluation map from increasing polynomial coordinates to this
primary basis is exactly $V_\beta$. Its determinant is the ordered
Vandermonde product, so it is invertible.

\subsection{The exact dual and primal maps of the strip quotient}
Use the coefficient presentation of OQC7--8 in Segre coordinates.
A biform of bidegree $(k,k)$ has the full coefficient expression
\[
 F=\sum_{a,b=0}^k f_{ab}
       X_1^aX_0^{k-a}Y_1^bY_0^{k-b}.
 \tag{OKM2}
\]
The literal variable dictionary to the OCF biform coefficients is
$X_1=z_0$, $X_0=z_1$, $Y_1=w_0$, $Y_0=w_1$.
The Segre substitution has
$t_{++}=X_1Y_1$, $t_{+-}=X_1Y_0$,
$t_{-+}=X_0Y_1$, $t_{--}=X_0Y_0$.
At the original exponential vector take
$X_1=e^{-i\delta w}$, $X_0=e^{i\delta w}$,
$Y_1=e^{\gamma w}$, $Y_0=e^{-\gamma w}$.
Then the monomial with coefficient $f_{ab}$ is $e^{\omega_{ab}w}$.
The original coefficient identification includes the entire invertible
matrix $\mathsf H$ of OQC1 before this Segre presentation; its unit
and period entries remain in the matrices produced by the strip
calculation. OQC7 is characterized by evaluation on $e^{wY}[1]$.
Since
$e^{wY}[1]=\sum_{a,b}e^{\omega_{ab}w}\varepsilon_{ab}$,
the independent exponentials in OQC4 therefore give the precise
complex-linear dual pairing
\[
 \Theta(F)(\varepsilon_{ab})=f_{ab},\qquad
 \Theta(F)(x)=f^{\mathsf T}x_{\rm prim}.
 \tag{OKM3}
\]
There is no binomial coefficient or complex conjugation in this
complex-linear dual pairing. The statement follows by equating all
coefficients of the distinct exponentials, whose independence is
proved by their Vandermonde derivative matrix in OQC4.

Let $J$ be exactly the coefficient subspace of the original observed
dual image, as identified in OQC8. The original strip and invariant
elimination provider OCF constructs its literal surjection
\[
 \pi:\mathbb C^q\twoheadrightarrow\mathbb C^{r_K},
 \qquad\ker\pi=J,\qquad r_K=8k-16.
 \tag{OKM4}
\]
This input is the actual OCF matrix, including its stated period
coefficients and its valid pivot choice. The coefficient pairing
OKM3 and OQC8 identify
$K:=\ker\Lambda=J^\perp$ in primary coordinates, where
$J^\perp=\{x:f^{\mathsf T}x=0\text{ for every }f\in J\}$.
Consequently its exact primal frame and polynomial frame are
\[
 I_{\rm prim}=\pi^{\mathsf T}:\mathbb C^{r_K}\hookrightarrow E_{\rm prim},
 \qquad I_{\rm poly}=V_\beta^{-1}\pi^{\mathsf T}:
                        \mathbb C^{r_K}\hookrightarrow E.
 \tag{OKM5}
\]
Here is a direct proof, including inverse coordinates. Choose any of
the explicit right sections $L_\pi$ of the surjective matrix $\pi$,
so $\pi L_\pi=I_{r_K}$. For $x=\pi^{\mathsf T}v$ and $f\in J$,
$f^{\mathsf T}x=(\pi f)^{\mathsf T}v=0$. The transpose is
injective because $L_\pi^{\mathsf T}\pi^{\mathsf T}=I_{r_K}$.
Conversely, if $x\in J^\perp$, then
$(I-L_\pi\pi)f\in J$ for every $f$, so
$x^{\mathsf T}(I-L_\pi\pi)f=0$ for every $f$. Thus
$x=\pi^{\mathsf T}L_\pi^{\mathsf T}x$.
This proves equality of the image and $K$, and its inverse on that
image is $L_\pi^{\mathsf T}$. In polynomial coordinates the inverse
is $L_\pi^{\mathsf T}V_\beta$. This proves the full typed primal
arrow, rather than selecting a Hermitian transpose in place of it.

\subsection{The complete original source and quotient metric}
For either original source $s\in\{1,k\}$ write $b_s=s/2$ and
$M_s=(2\pi)^{s/2}$. The unchanged real-coordinate measure is
\[
 dm_{0,s}(y)=M_s\frac{2^{b_s}}{4\pi\Gamma(b_s)}
      |\Gamma(b_s/2+iy/2)|^2\,dy,
 \qquad S=c+iy.
 \tag{OKM6}
\]
The complete canonical OGN1--3 proof gives the monic real
orthogonal polynomials and all their squared norms:
\[
 \begin{gathered}
 p_0(y)=1,\quad p_1(y)=y,\quad
 p_{d+1}(y)=yp_d(y)-d(b_s+d-1)p_{d-1}(y),\\
 h_d=M_s d!(b_s)_d,\qquad
 u_d(S)=\frac{p_d((S-c)/i)}{\sqrt{h_d}}.
 \end{gathered}\tag{OKM7}
\]
The square root is positive. Thus $u_d$ has the original leading
coefficient $i^{-d}/\sqrt{h_d}$ in $S$ and is orthonormal on the
line in OKM6. No source mass or phase has been removed.

For a polynomial cutoff $D\geq q-1$ let
\[
 \begin{gathered}
 T_D=\left(\frac{p_d(\omega_{ab})}{\sqrt{h_d}}\right)_
                        {(a,b),\,0\leq d\leq D},
 \qquad R_D=T_DT_D^*,\\
 G_D^{\rm prim}=R_D^{-1},\qquad
 G_D^{\rm poly}=V_\beta^*R_D^{-1}V_\beta,\qquad
 H_D^K=\overline\pi\,R_D^{-1}\pi^{\mathsf T}.
 \end{gathered}\tag{OKM8}
\]
The star denotes conjugate transpose, while the superscript
$\mathsf T$ is ordinary transpose. In particular
$(\pi^{\mathsf T})^*=\overline\pi$, with precisely the indicated
matrix dimensions.

These are the actual minimum quotient metrics, as follows directly
from the original remainder map. In the orthonormal coefficient
frame $u_0,\ldots,u_D$, the remainder map followed by primary
evaluation is exactly $T_D$, column by column. Its first $q$
columns are invertible: the monic polynomials $p_0,\ldots,p_{q-1}$
evaluated at the distinct $\omega_{ab}$ have a nonzero Vandermonde
determinant, and all $\sqrt{h_d}$ are nonzero. Thus $R_D>0$.
For a prescribed primary vector $x$, the vector
$T_D^*R_D^{-1}x$ maps to $x$. It is orthogonal to $\ker T_D$,
because its inner product with a kernel vector contains $T_D$
applied to that vector. Every other lift differs by such a vector,
so its squared norm is
$x^*R_D^{-1}x$ plus the squared norm of that difference.
This proves the first quotient metric in OKM8.
The identity $x_{\rm prim}=V_\beta x_{\rm poly}$ proves the
second by direct pairing. Restriction to the exact frame OKM5
then gives
\[
 I_{\rm poly}^*G_D^{\rm poly}I_{\rm poly}
 = (\pi^{\mathsf T})^*R_D^{-1}\pi^{\mathsf T}
 =H_D^K>0.
 \tag{OKM9}
\]
The last positivity follows from injectivity of $\pi^{\mathsf T}$.
The map $K\hookrightarrow E$ is the fixed original coefficient
map. No invariance of that kernel under multiplication is used.

\subsection{Christoffel--Darboux with the exact source mass}
For arbitrary $z,w\in\mathbb C$ and $D\geq0$, the complete
polynomial identity is
\[
 \sum_{d=0}^D\frac{p_d(z)p_d(w)}{h_d}
 =\begin{cases}
 \displaystyle\frac{p_{D+1}(z)p_D(w)-p_D(z)p_{D+1}(w)}
                     {h_D(z-w)},&z\ne w,\\[6pt]
 \displaystyle\frac{p_{D+1}'(z)p_D(z)-p_D'(z)p_{D+1}(z)}{h_D},
                       &z=w.
 \end{cases}\tag{OKM10}
\]
To prove it, put $a_d=d(b_s+d-1)$, with $a_0=0$, and set
$p_{-1}=0$ for this recurrence expression. The recurrence
in OKM7 gives
\[
 \begin{split}
 (z-w)p_d(z)p_d(w)
 &=p_{d+1}(z)p_d(w)-p_d(z)p_{d+1}(w)\\
 &\quad-a_d\bigl(p_d(z)p_{d-1}(w)-p_{d-1}(z)p_d(w)\bigr).
 \end{split}\tag{OKM11}
\]
For $d\geq1$, the full norm identity is $a_d/h_d=1/h_{d-1}$.
Divide OKM11 by $h_d$ and sum. Every adjacent term cancels,
leaving just the degree-$D$ numerator divided by $h_D$; the
$d=0$ lower term is zero. This proves the first line of OKM10.
Both sides of the resulting equality after multiplication by
$z-w$ are polynomials. Differentiating with respect to $z$ and
then setting $z=w$ proves the second line with the displayed
order and sign. This also proves the value at every confluent
point without an assertion about a limiting singular quotient.
The denominator is the full $h_D=M_sD!(b_s)_D$, including $M_s$.

Since all coefficients of the recurrence are real, the matrix
entries in OKM8 are therefore exactly the values of OKM10 at
\[
 z=\omega_{ab},\qquad w=\overline{\omega_{a'b'}}.
 \tag{OKM12}
\]
Their indexed conjugation and confluent locus are
\[
 \begin{gathered}
 \overline{\omega_{a'b'}}=\omega_{k-a',b'},\qquad
 \omega_{ab}-\overline{\omega_{a'b'}}
       =2(b-b')\gamma-2i(a+a'-k)\delta,\\
 \omega_{ab}=\overline{\omega_{a'b'}}
       \quad\Longleftrightarrow\quad b=b'\text{ and }a+a'=k.
 \end{gathered}\tag{OKM13}
\]
These follow by substituting OKM1; positivity of $\delta,\gamma$
proves both implications in the last line. Thus the conjugation
map on the point indices is precisely $(a,b)\mapsto(k-a,b)$.
The confluent formula of OKM10 applies at these indexed entries.
For the retained odd $k$, its matrix diagonal has $a'=a,b'=b$,
so $2a=k$ cannot occur; the literal Hermitian diagonal is given
by the first line of OKM10. The conjugate-index entries use its
second line. All entries, including the positive diagonal, are
simultaneously defined by the original polynomial sum on the
left of OKM10.

\subsection{Every pivot change and the four original endpoints}
Let $\pi'$ be another valid strip/invariant elimination matrix for
the same original $J$, with the same surjective codomain dimension.
Choose right sections $L_\pi,L_{\pi'}$ and put
\[
 A_{\pi',\pi}=\pi'L_\pi,
 \qquad A_{\pi,\pi'}=\pi L_{\pi'}.
 \tag{OKM14}
\]
Because $\pi'$ vanishes on $J=\ker\pi$, one has
$\pi' L_\pi\pi=\pi'$. Thus
$\pi'=A_{\pi',\pi}\pi$. Reversing the roles gives the
opposite identity, and surjectivity of each matrix then proves
$A_{\pi',\pi}A_{\pi,\pi'}=I$ and
$A_{\pi,\pi'}A_{\pi',\pi}=I$. This proves invertibility and
also independence from the chosen right sections. Writing
$A=A_{\pi',\pi}$ for these identities gives the exact frame
and Gram changes
\[
 \begin{gathered}
 I'_{\rm prim}=I_{\rm prim}A^{\mathsf T},\qquad
 I'_{\rm poly}=I_{\rm poly}A^{\mathsf T},\\
 (H_D^K)'=\overline A H_D^K A^{\mathsf T},\qquad
 \det(H_D^K)'=|\det A|^2\det H_D^K.
 \end{gathered}\tag{OKM15}
\]
The first two formulas follow by transposing $\pi'=A\pi$.
Substitution into OKM9 proves the congruence and determinant.
This retains all actual pivot determinants. They have not been
assigned modulus one.

The kernel is the same fixed $K=\ker\Lambda$ used in MRI1--5.
If its original BRU inclusion is $I_0$ and its fixed-section kernel
coordinate map is $\kappa_0$, let
\[
 \mathsf C_K=\kappa_0I_{\rm poly},\qquad
 \mathsf C_K^{-1}=L_\pi^{\mathsf T}V_\beta I_0.
 \tag{OKM16}
\]
The equalities $I_0\kappa_0x=x$ for $x\in K$ and
$L_\pi^{\mathsf T}V_\beta I_{\rm poly}=I_{r_K}$ prove that
these are inverse matrices and $I_0\mathsf C_K=I_{\rm poly}$.
Consequently
$H_D^K=\mathsf C_K^*(I_0^*G_D^{\rm poly}I_0)\mathsf C_K$.
The matrix $\mathsf C_K$ is independent of $D$ and of the choice of source
order $s$. The source metrics change; the two specified coefficient
frames do not.

It follows that the complete original four-endpoint kernel scalar is
\[
 \boxed{\displaystyle
 F_K^{(s)}=
 \log\det H_{q-1}^K+\log\det H_q^K
       -\log\det H_{2q-1}^K-\log\det H_{2q}^K.}
 \tag{OKM17}
\]
Indeed MRI3 gives the actual consecutive kernel determinant ratio
\[
 \frac{\det(I_0^*G_{q+j-1}^{\rm poly}I_0)}
      {\det(I_0^*G_{q+j}^{\rm poly}I_0)}
 =\frac{1+t_j}{1+\xi_j}.
\]
Each factor $|\det\mathsf C_K|^2$ cancels in that ratio by OKM16.
The weighted sum with endpoint weights $1,2,\ldots,2,1$ therefore
has precisely the four signs in OKM17, as in MRI4. This verifies
the original degree cutoffs, rather than shifting the column count
in $T_D$. Under any valid pivot change OKM15, the added scalar
$2\log|\det A|$ has total coefficient $1+1-1-1=0$, proving
invariance of OKM17 with its original frames.

For emphasis, the full mass in OKM7 appears in every column of
$T_D$ and in the denominator of OKM10. If that common mass is
factored explicitly, write $T_D=M_s^{-1/2}\widehat T_D$, where
$\widehat T_D$ has exactly the same numerator polynomials and
the denominator $\sqrt{d!(b_s)_d}$. Then
$H_D^K=M_s\overline\pi(\widehat T_D\widehat T_D^*)^{-1}
\pi^{\mathsf T}$, so its determinant has factor $M_s^{r_K}$.
It cancels only in the displayed signed four-endpoint scalar,
whose four ranks are equal. The individual metrics in OKM8--9
retain it. No inserted Gamma order, discarded action defect, or
kernel-invariance assertion enters any of these calculations.

% END COMPLETE OKM

\clearpage
% BEGIN COMPLETE OKA
\section{The original kernel graph, boundary basis and complete action}
\label{sec:OKA}

This section uses the complete conductor coefficient construction OCF
and its exact original dual identification in OKM. Fix the original
simple quartet, its complete amplitude unit and periods, the original
parameter on any logarithm branch with $|u|\geq R_*$, and
$k=4l+1$, $l\geq1$. Put $q=(k+1)^2$, $r=8k-16$, and
$\mathcal I_k=\{(a,b):0\leq a,b\leq k\}$. The roots, centre and
centred eigenvalues remain
\[
 c=k/2,\qquad
 \beta_{ab}=c+(2a-k)\delta+i(2b-k)\gamma,\qquad
 \omega_{ab}=(\beta_{ab}-c)/i
             =(2b-k)\gamma-i(2a-k)\delta.
 \tag{OKA1}
\]
The original algebra is $E=\mathbb C[S]/(\chi)$, with
$\chi(S)=\prod_{(a,b)\in\mathcal I_k}(S-\beta_{ab})$.
Let $\varepsilon_{ab}$ be its original primary idempotents and
$\mathcal V:E\to\mathbb C^{\mathcal I_k}$ the evaluation
isomorphism $[p]\mapsto(p(\beta_{ab}))_{ab}$. Its inverse sends
$e_{ab}$ to
\[
 \varepsilon_{ab}(S)
 =\frac{\chi(S)}{(S-\beta_{ab})\chi'(\beta_{ab})}.
 \tag{OKA2}
\]
These denominators are nonzero: equality of two roots in OKA1,
on comparing real and imaginary parts and using $\delta>0$,
$\gamma>2$, implies equality of both indices. Evaluation proves
that the polynomials in OKA2 are the coordinate idempotents,
so their columns give the exact inverse, with no changed root
or polynomial coefficient convention.

\subsection{The coefficient quotient gives a graph in the original primary basis}

Let $\pi:\mathbb C^{\mathcal I_k}\to\mathbb C^r$ be the
surjective coefficient quotient of OCF, whose kernel is precisely
the image $J_k$ of the invariant polynomials under the original
restriction. The domain here consists of the coefficients of the
biforms
$z_0^a z_1^{k-a}w_0^b w_1^{k-b}$, and is the complex-linear
dual of the displayed primary coordinate space. Let
$T\subset\mathcal I_k$ be the final OCF complement of its
invariant-image pivot rows, viewed as a subset of the full
biform coefficient grid. Its ordered coordinate inclusion is
$E_T:\mathbb C^r\to\mathbb C^{\mathcal I_k}$. The explicit
OCF formula gives
\[
 \pi E_T=I_r,\qquad
 N=\pi^T:\mathbb C^r\longrightarrow\mathbb C^{\mathcal I_k},
 \qquad E_T^TN=I_r.
 \tag{OKA3}
\]
Indeed a retained strip monomial in $T$ already has zero
conductor remainder outside its own coordinate, zero coordinates
on the invariant pivot rows, and its indicated coordinate equal
to one. Both reductions in OCF therefore leave that quotient
coordinate equal to its standard unit vector. Transposing gives
the last identity. This uses ordinary transpose, because the
restriction quotient and its dual are complex-linear maps.

Write $J=\mathcal I_k\setminus T$ only for the complementary
index set in this section; the invariant image retains the name
$J_k$. Denote coordinate restrictions by $x_T,x_J$, and put
$N_J=E_J^TN$. The original kernel inclusion is consequently
\[
 I_K=\mathcal V^{-1}N:\mathbb C^r\hookrightarrow E,
 \qquad
 \operatorname{im}I_K=K:=\ker\Lambda,\qquad
 N=\begin{pmatrix}I_r\\N_J\end{pmatrix}_{T,J}.
 \tag{OKA4}
\]
For completeness, a primary vector $v$ pairs with a biform
coefficient vector $f$ by $f^Tv$. The annihilator of
$\ker\pi=J_k$ is $\operatorname{im}\pi^T$: containment
follows by multiplication, while surjectivity of $\pi$ gives
both subspaces dimension $r$. OQC8 identifies precisely this
annihilator with the primary coordinates of the fixed original
$\ker\Lambda$. This proves both the original image assertion
and its direction in OKA4. In polynomial coordinates the same
inclusion is $V_\beta^{-1}N$, where
$(V_\beta)_{(a,b),d}=\beta_{ab}^d$ for $0\leq d<q$;
this is the full coefficient matrix of OKA2.

\subsection{An explicit basis of the actual boundary}

For $j\in J$ define the actual boundary vector
$b_j=\Lambda\varepsilon_j$. These vectors form a basis of
the original image $B=\operatorname{im}\Lambda$. To prove
independence, if $\sum_{j\in J}c_j\varepsilon_j$ lies in
$K$, its primary coordinates equal $N\zeta$. Its $T$
coordinates vanish, so OKA3 gives $\zeta=0$ and all $c_j=0$.
To prove spanning, use the literal decomposition
\[
 x=N x_T+E_J(x_J-N_Jx_T),\qquad
 \Lambda\mathcal V^{-1}x
       =\sum_{j\in J}b_j(x_J-N_Jx_T)_j.
 \tag{OKA5}
\]
The first equality follows separately on its $T$ and $J$
coordinates. Apply $\Lambda$, which kills the first term by
OKA4, to obtain the second. Thus no observation has been
substituted: in this basis the specified original observation
has the exact coefficient matrix $(-N_J,I_{q-r})$.

Here are the complete original period and tensor coordinates of
these boundary vectors. Write the quartet roots as
$\alpha_{\epsilon,\eta}=1/2+\epsilon\delta+i\eta\gamma$,
where $\epsilon,\eta\in\{+1,-1\}$ and the signs $+,-$ in
subscripts denote these same two numbers. In a subscript
$\alpha$ below means that complex root, indexed by its sign
pair. Keep the original one-factor primary idempotents
$\varepsilon^h_\alpha\in E_h$. Set
\[
 v_\alpha=\Pi U\varepsilon^h_\alpha
       =v_h(\alpha)\Pi\varepsilon^h_\alpha,
 \qquad
 Z_+=z_0,\ Z_-=z_1,\ W_+=w_0,\ W_-=w_1,
 \quad
 \mathfrak v(z,w)=\sum_{\epsilon,\eta}
                         v_{\epsilon\eta}Z_\epsilon W_\eta.
 \tag{OKA6}
\]
The vectors belong to the original one-factor period space;
all entries of $\Pi$ and all four actual nonzero amplitude
values remain present. For every $(a,b)\in\mathcal I_k$,
\[
 \boxed{\quad
 \jmath\Lambda\varepsilon_{ab}
  =[z_0^a z_1^{k-a}w_0^b w_1^{k-b}]
       P_k\bigl(\mathfrak v(z,w)^{\otimes k}\bigr),\qquad
 P_k=\frac15\sum_{h=0}^4(C^{-h})^{\otimes k}.
 \quad}\tag{OKA7}
\]
The original sum inclusion $\iota:E\hookrightarrow E_h^{\otimes k}$
sends $\varepsilon_{ab}$ to the sum of the primary tensor
idempotents whose word has $a$ positive real signs and $b$
positive imaginary signs. This assertion follows by evaluation:
on a primary word, the sum variable takes exactly the value
$\beta_{ab}$ of its two counts, and OKA2 evaluates to one
for those counts and zero otherwise. Applying
$(\Pi U)^{\otimes k}$ gives precisely the corresponding sum
of the tensors of the $v_\alpha$ in OKA6. Expanding the tensor
power and taking the indicated biform coefficient selects
exactly the same words, each once. The original definition
$\jmath\Lambda=P_k(\Pi U)^{\otimes k}\iota$ proves OKA7,
including the literal inverse powers and $1/5$. Its specialization
to $j\in J$ therefore supplies the actual basis $b_j$ of
OKA5. Neither a different target nor a numerical period matrix
has been introduced.

The fixed, source-independent kernel coordinate and section are
now explicitly
\[
 \begin{gathered}
 \kappa=E_T^T\mathcal V:E\to\mathbb C^r,
 \qquad S_*:B\to E,\quad S_*b_j=\varepsilon_j,
 \\
 I_K\kappa+S_*\Lambda=I_E,
 \quad\kappa I_K=I_r,\quad\kappa S_*=0,\quad\Lambda S_*=I_B.
 \end{gathered}
 \tag{OKA8}
\]
Every equality follows by applying OKA5 to an arbitrary primary
vector. In particular this is an actual splitting of the fixed
original short exact sequence, used unchanged for both source
orders and for all polynomial degrees.

\subsection{The whole original action and its retained defect}

The original centred sum action is $Y=(S-c)/i$ on $E$.
Under $\mathcal V$ it is the diagonal matrix of the
$\omega_{ab}$ in OKA1. Write its two diagonal submatrices
as $\Omega_T$ and $\Omega_J$. In the exact coordinates
$e=I_K\zeta+S_*b$ of OKA8, it has the complete matrix
\[
 \boxed{\quad
 Y\ \longleftrightarrow\
 \begin{pmatrix}
   \Omega_T&0\\
   \Omega_JN_J-N_J\Omega_T&\Omega_J
 \end{pmatrix},\qquad
 (\Omega_JN_J-N_J\Omega_T)_{jt}
                 =(\omega_j-\omega_t)(N_J)_{jt}.
 \quad}\tag{OKA9}
\]
Indeed the primary vector of $e$ is
$(\zeta,N_J\zeta+b)$. Applying the original diagonal matrix
gives $(\Omega_T\zeta,\Omega_JN_J\zeta+\Omega_Jb)$.
Applying the coordinates in OKA5 subtracts
$N_J\Omega_T\zeta$ from its second component and proves
every block in OKA9. This proves, with its precise codomain,
\[
 YI_K-I_K\Omega_T
     =S_*\bigl(\Omega_JN_J-N_J\Omega_T\bigr)
                :\mathbb C^r\to E,\qquad
 \Lambda YI_K=\Omega_JN_J-N_J\Omega_T:\mathbb C^r\to B,
 \tag{OKA10}
\]
where the matrix on the right uses the basis $b_j$.
The chosen complementary span of the original primary idempotents
is stable under $Y$; no stability of $K$ has been used. Its
entire failure is the displayed coefficient map.

For each original source $s\in\{1,k\}$ retain the monic
Gamma polynomials $p_d^{(s)}$, the complete norms
$h_d^{(s)}=(2\pi)^{s/2}d!(s/2)_d$, and the original
$e_d^{(s)}=[p_d^{(s)}(Y)]/\sqrt{h_d^{(s)}}$. Then the
fixed-section kernel and boundary coordinates of every actual
remainder, including all degrees $d\geq q$, are
\[
 \begin{split}
 (\kappa e_d^{(s)})_t
      &=\frac{p_d^{(s)}(\omega_t)}{\sqrt{h_d^{(s)}}},\\
 (\Lambda e_d^{(s)})_j
      &=\frac{p_d^{(s)}(\omega_j)
                -\sum_{t\in T}(N_J)_{jt}p_d^{(s)}(\omega_t)}
              {\sqrt{h_d^{(s)}}}.
 \end{split}\tag{OKA11}
\]
Polynomial evaluation in the original primary algebra gives the
full vector $p_d^{(s)}(\omega)/\sqrt{h_d^{(s)}}$ even when
its polynomial representative must be reduced modulo $\chi$.
Applying OKA5 and OKA8 proves OKA11. Thus the formula includes
the original finite remainder operation and every phase and source
mass. Composing its second line with the actual columns OKA7
gives exactly the period-vector increments of OPG. Its first line
is the fixed-section kernel vector used by OPG and OPR; the
degree-dependent minimum-section correction in OPR remains
required for the mixed norm. The metric computation in OKM
applies to the common inclusion OKA4 and does not replace that
correction by the fixed-section vector alone.

% END COMPLETE OKA

\clearpage
% BEGIN COMPLETE OMG
% Full original multiplicity geometry. No frozen predecessor is edited.
\section{The complete homogeneous geometry of the original multiplicity curve}
\label{sec:OMG}

\subsection{The actual primary, unit, period and Fourier coefficient map}
Retain the original full-order quartet and every original period
parameter. Write
\[
\begin{gathered}
 \alpha_{\epsilon,\eta}=\tfrac12+\epsilon\delta+i\eta\gamma,
 \quad\epsilon,\eta\in\{-1,1\},\quad0<\delta<\tfrac12,\quad\gamma>2,\\
 m\ge1,\quad d_m=m-1,\quad n=4m,\quad N=4m+1,\quad
 h(z)=\prod_{\epsilon,\eta}(z-\alpha_{\epsilon,\eta})^m,\\
 E_h=\mathbb C[z]/(h),\quad A_h=M_z,\quad
 v_h=2\xi/h,\quad U=M_{j_hv_h},\quad v_h(\alpha_{\epsilon,\eta})\ne0.
\end{gathered}\tag{OMG1}
\]
The last inequalities follow from the stipulated complete original
zero order $m$. The map $j_h$ records all divided Taylor coefficients
of orders $0,\ldots,m-1$ at all four roots. Let
$\mathcal E_{\rm prim}:\mathbb C^{4m}\to E_h$ send its coordinate
$(\alpha,d)$ to the original primary element $\varepsilon_\alpha
(z-\alpha)^d$, $0\le d<m$. Explicitly its representative is
\[
 E_{\alpha,d}(z)=\operatorname{rem}_h\left[
 h_\alpha(z)\sum_{j=0}^{m-1}
       \frac{(h_\alpha^{-1})^{(j)}(\alpha)}{j!}(z-\alpha)^{j+d}
                                      \right],\quad
 h_\alpha=h/(z-\alpha)^m.
\]
The Taylor inverse is defined since $h_\alpha(\alpha)\ne0$.
It gives the idempotent one on its own primary factor and zero
on the others. The four coprime factors therefore prove that
$\mathcal E_{\rm prim}$ is an isomorphism and include every
nilpotent degree. In this basis the full matrix of $U$ is
\[
 U_{\rm prim;\alpha,a;\alpha,d}
 =\begin{cases}v_h^{(a-d)}(\alpha)/(a-d)!,&a\ge d,\\0,&a<d,
   \end{cases}\qquad
 \det U=\prod_\alpha v_h(\alpha)^m\ne0.
 \tag{OMG2}
\]
Other primary blocks are zero. Inverting its constant coefficient
plus its nilpotent part by the finite geometric series proves
invertibility with all derivatives retained.

Keep the original $u\ne0$ and its logarithm branch, the primitive
$\Phi_h'=h$, $\Phi_h(0)=0$, the outward rays $\ell_j$ of angles
$(\arg u+\pi+2\pi j)/N$, and $\Gamma_j=-\ell_0+\ell_j$.
The exact period and Fourier maps are
\[
\begin{gathered}
 (\Pi[P])_j=\int_{\Gamma_j}e^{\Phi_h(z)/u}
                     \operatorname{rem}_hP(z)\,dz,\quad1\le j\le n,\\
 \zeta=e^{2\pi i/N},\quad F_{jt}=\zeta^{jt}-1,\quad
 (F^{-1})_{tj}=\zeta^{-tj}/N\quad(1\le j,t\le n),\quad
 C=F\operatorname{diag}(\zeta,\ldots,\zeta^n)F^{-1}.
\end{gathered}\tag{OMG3}
\]
The full original period theorem PDE/OPG gives the convergence and
invertibility of $\Pi$ on these representatives. The Fourier inverse
follows directly from
$\sum_{j=1}^{N-1}\zeta^{-tj}(\zeta^{sj}-1)=N\mathbf1_{t=s}$.
It is the same literal cyclic matrix as OCS4, including its inverse
powers in the original average. Its inherited Gram is
$F^*F=N(I_n+\mathbf1\mathbf1^*)$, as expansion of its entries gives.

Define the exact diagonal map and coefficient matrix
\[
\begin{gathered}
 D_m^{\rm jet}=\operatorname{diag}\left(i^{-d}/d!\right)_{\alpha,\,0\le d<m},
 \qquad \mathsf H_m=F^{-1}\Pi U\mathcal E_{\rm prim}D_m^{\rm jet},\\
 (\mathsf H_m)_{t;\alpha,d}
  =\frac{i^{-d}}{d!}\sum_{a=d}^{m-1}
     \left[\frac1N\sum_{j=1}^{n}\zeta^{-tj}
              \int_{\Gamma_j}e^{\Phi_h(z)/u}E_{\alpha,a}(z)\,dz\right]
                  \frac{v_h^{(a-d)}(\alpha)}{(a-d)!}.
\end{gathered}\tag{OMG4}
\]
This is exactly OCS13, now regarded as one square matrix. All
factors in its first line are invertible. In particular
\[
 \det\mathsf H_m=(\det F)^{-1}\det\Pi\,
      \left(\prod_\alpha v_h(\alpha)^m\right)
      \det\mathcal E_{\rm prim}
      \prod_\alpha\prod_{d=0}^{m-1}\frac{i^{-d}}{d!}\ne0,
\]
with the determinant of $\mathcal E_{\rm prim}$ taken in the fixed
original coefficient frame. None of these factors is prescribed.

Set $\omega_{\epsilon,\eta}=\eta\gamma-i\epsilon\delta$ and let
\[
 z_{\epsilon,\eta,d}(w)=w^d e^{\omega_{\epsilon,\eta}w},\qquad
 x(w)=F^{-1}\Pi Ue^{w(A_h-1/2)/i}[1].
 \tag{OMG5}
\]
The terminating exponential on each primary factor has coefficient
$i^{-d}w^d/d!$ on $\varepsilon_\alpha(z-\alpha)^d$.
Multiplication by OMG2 and application of the literal period maps
therefore prove the exact identity and inverse coordinate map
\[
 \boxed{x(w)=\mathsf H_m z(w),\qquad z(w)=\mathsf H_m^{-1}x(w).}
 \tag{OMG6}
\]
For the original Euclidean period metric, its Gram on these primary
curve coordinates is $\mathsf H_m^*N(I_n+\mathbf1\mathbf1^*)
\mathsf H_m$. Thus the coefficient change also retains its full
metric and all off-diagonal entries, without claiming orthogonality.

\subsection{The complete homogeneous ideal, with quadratic generators}
Use binary indices $a,b\in\{0,1\}$ for
$\epsilon=2a-1$, $\eta=2b-1$. Let
$\mathscr P_Z=\mathbb C[Z_{a,b,d}:a,b\in\{0,1\},0\le d\le d_m]$
with each variable of degree one. Define the algebra map
\[
 \Phi_m:\mathscr P_Z\longrightarrow
        \mathbb C[X_0,X_1,Y_0,Y_1,T_0,T_1],\qquad
 Z_{a,b,d}\longmapsto X_aY_bT_0^{d_m-d}T_1^d.
 \tag{OMG7}
\]
Its image is the complete diagonal graded algebra
\[
 R_m=\bigoplus_{k\ge0}
  \mathbb C[X_0,X_1]_k\otimes\mathbb C[Y_0,Y_1]_k
                              \otimes\mathbb C[T_0,T_1]_{kd_m}.
 \tag{OMG8}
\]
Indeed a basis monomial of its degree-$k$ component is
\[
 X_0^{k-a}X_1^aY_0^{k-b}Y_1^bT_0^{kd_m-D}T_1^D,
       \quad0\le a,b\le k,\quad0\le D\le kd_m.
 \tag{OMG9}
\]
It is a product of $k$ images in OMG7: choose $k$ binary first
indices with sum $a$, independently $k$ binary second indices with
sum $b$, and $k$ integers in $[0,d_m]$ with sum $D$; the last
choice follows by filling successive slots up to $d_m$. Combining
the three lists into rows gives the required preimage monomial.
Distinct displayed target monomials are linearly independent.

Let $\mathfrak t_m$ be generated by every quadratic binomial
\[
 Z_{a,b,d}Z_{a',b',d'}-Z_{c,e,f}Z_{c',e',f'}
 \quad\text{with}\quad
 a+a'=c+c',\quad b+b'=e+e',\quad d+d'=f+f',
 \tag{OMG10}
\]
where every index belongs to its original range. Repeated variables
are allowed. Then
\[
 \boxed{\ker\Phi_m=\mathfrak t_m,\qquad
           \mathscr P_Z/\mathfrak t_m\simeq R_m.}
 \tag{OMG11}
\]
Here is a complete ideal proof. Each displayed quadratic has zero
image. Two degree-$k$ monomials have the same image exactly when
the sums of their first indices, second indices and third indices
agree. Regard their $k$ factors temporarily as labelled rows. Align
the first indices with the target list by pairwise swaps of first
indices, keeping each row's other indices fixed. Each swap is an
OMG10 relation. Next align the second indices by the same operation,
keeping the aligned first indices fixed. Finally, if the third
indices differ from their target list, choose a row below its target
and a row above its target. Increase the former third index by one
and decrease the latter by one. These entries stay in $[0,d_m]$
because their target entries do; the total sum is unchanged and the
sum of absolute deviations decreases by two. This terminates at
the target list. Each transfer is another OMG10 relation. Thus the
difference between any two monomials with the same image is a sum
of multiples of the quadrics, by telescoping the successive monomial
differences. Labels only specify the sequence of polynomial maps;
the products are the original commutative monomials throughout.
For any polynomial in $\ker\Phi_m$, collect its coefficients by
their target monomial. Each such coefficient sum is zero by target
monomial independence. Subtracting one chosen monomial in each
group expresses that whole group as a linear combination of the
proved monomial differences. This proves the reverse inclusion for
all degrees and hence the entire ideal equality.

Every degree-$k$ original curve monomial has the function
\[
 w^D\exp\{[(2b-k)\gamma-i(2a-k)\delta]w\}
       \quad(0\le a,b\le k,\ 0\le D\le kd_m).
 \tag{OMG12}
\]
The frequencies are distinct: equality of real parts gives equality
of $b$ and equality of imaginary parts gives equality of $a$.
To prove independence with all polynomial powers, suppose
$\sum_{a,b}P_{a,b}(w)e^{\omega_{a,b}w}=0$, with
$\deg P_{a,b}\le kd_m$. Fix a pair and apply
$\prod_{(a',b')\ne(a,b)}(\partial_w-\omega_{a',b'})^{kd_m+1}$.
This kills every other summand. On its surviving polynomial each
factor is $\partial_w+(\omega_{a,b}-\omega_{a',b'})$, invertible
on the space of degree at most $kd_m$ because differentiation is
nilpotent and the constant is nonzero. The surviving polynomial
is therefore zero. Repeating this for every pair proves independence.

Thus the full homogeneous ideal of the original primary curve is
exactly $\mathfrak t_m$: a homogeneous polynomial restricts to zero
if and only if its coefficient sums at every target monomial OMG9
are zero. In particular
\[
 \boxed{\dim R_{m,k}=(k+1)^2[1+k(m-1)]=q_k,
 \quad\ker\{\mathscr P_{Z,k}\to\operatorname{Hol}(\mathbb C),
                      f\mapsto f(z(w))\}=\mathfrak t_{m,k}.}
 \tag{OMG13}
\]
This statement concerns the homogeneous ideal, equivalently the
projective closure of $[z(w)]$. Every one of its degrees has been
computed, including degree zero. No assertion about a different
nonhomogeneous affine ideal is used.

\subsection{The projective model and its complete chart inverses}
For $m\ge2$, the morphism
\[
 (\mathbb P^1)^3\longrightarrow\mathbb P^{4m-1},\qquad
 ([X_0:X_1],[Y_0:Y_1],[T_0:T_1])
      \longmapsto[X_aY_bT_0^{d_m-d}T_1^d]_{a,b,d}
 \tag{OMG14}
\]
has multidegree $(1,1,m-1)$. At least one first, second and third
coordinate is nonzero; their corresponding extreme product is
nonzero, so the map is everywhere defined. Its image and scheme
structure are exactly $\operatorname{Proj}R_m$, with complete
closed ideal OMG10. Here are explicit inverses proving this claim.
The eight charts $Z_{a,b,0}\ne0$ and $Z_{a,b,d_m}\ne0$ cover
$\operatorname{Proj}R_m$. Indeed OMG11 gives
\[
 Z_{a,b,d}^{d_m}=Z_{a,b,0}^{d_m-d}Z_{a,b,d_m}^{d}
             \quad\text{in }R_m.
\]
A homogeneous prime containing every extreme variable contains
every variable by this equality, and hence gives no point of Proj.
On $Z_{a,b,0}\ne0$ the three inverse affine coordinates are
\[
 x=Z_{1-a,b,0}/Z_{a,b,0},\quad
 y=Z_{a,1-b,0}/Z_{a,b,0},\quad
 t=Z_{a,b,1}/Z_{a,b,0}.
\]
For every variable its ratio to $Z_{a,b,0}$ is
$x^{\mathbf1_{a'\ne a}}y^{\mathbf1_{b'\ne b}}t^d$, by OMG7
and its proved injective quotient. Consequently the degree-zero
localized ring is $\mathbb C[x,y,t]$. These three coordinates
are algebraically independent, since their images in the parameter
fraction field are $X_{1-a}/X_a,Y_{1-b}/Y_b,T_1/T_0$.
For the chart at $Z_{a,b,d_m}$ replace the last ratio by
$t=Z_{a,b,d_m-1}/Z_{a,b,d_m}$; the general ratio has exponent
$d_m-d$ in $t$. Its coordinate ring is again $\mathbb C[x,y,t]$.
On chart overlaps these formulas invert the same parameter ratios,
so they glue to the inverse of OMG14. This proves the isomorphism
of schemes and its full defining homogeneous ideal. For $m=1$
the third factor has degree zero and the same argument uses the
four first/second charts: the model is the Segre $(\mathbb P^1)^2$
with ideal $Z_{0,0,0}Z_{1,1,0}-Z_{0,1,0}Z_{1,0,0}$.

In the actual Fourier coordinates let
$\mathscr P_x=\mathbb C[x_1,\ldots,x_n]$ and define
\[
 \mathfrak a=\{f(x):f(\mathsf H_mZ)\in\mathfrak t_m\},\qquad
 R=\mathscr P_x/\mathfrak a.
 \tag{OMG15}
\]
The invertible substitution OMG6 proves that $\mathfrak a$ is the
complete homogeneous ideal of $x(w)$ and that $R\simeq R_m$.
Its explicit generators are all OMG10 quadrics with every $Z$
replaced by its component of $\mathsf H_m^{-1}x$. These are the
actual period and unit coefficients of OMG4, and $\mathfrak a$
is prime because the target ring in OMG7 is a domain. Thus the
original curve's projective closure is the precise linear image
of OMG14 under $\mathsf H_m$, with no prescribed substitute matrix.

\subsection{The exact graded source-dual algebra and every factorial}
For each $k\ge0$ retain the original sum algebra
\[
\begin{gathered}
 e_k=1+kd_m,\quad c_k=k/2,\quad
 \beta_{k,a,b}=c_k+(2a-k)\delta+i(2b-k)\gamma,\\
 \chi_k(S)=\prod_{a,b=0}^k(S-\beta_{k,a,b})^{e_k},\quad
 E_k=\mathbb C[S]/(\chi_k),\quad Y_k=(M_S-c_kI)/i.
\end{gathered}\tag{OMG16}
\]
For $k=0$ this is $E_0=\mathbb C[S]/(S)=\mathbb C$. Let
$H_{k,a,b,D}=\varepsilon_{k,a,b}(S-\beta_{k,a,b})^D$,
$0\le D<e_k$, be the complete original primary basis obtained
by the same Taylor inverse as in OMG2. It satisfies
\[
 e^{wY_k}[1]=\sum_{a,b=0}^k e^{\omega_{k,a,b}w}
             \sum_{D=0}^{e_k-1}\frac{i^{-D}w^D}{D!}H_{k,a,b,D},
 \quad\omega_{k,a,b}=(2b-k)\gamma-i(2a-k)\delta.
 \tag{OMG17}
\]
For $f\in\mathscr P_{x,k}$ expand its original function uniquely as
$f(x(w))=\sum c_{a,b,D}(f)w^De^{\omega_{k,a,b}w}$ using OMG12.
Define the complex linear dual map, without a conjugation,
\[
 \Theta_k(f)(H_{k,a,b,D})=i^DD!c_{a,b,D}(f).
 \tag{OMG18}
\]
Then OMG17 proves exactly
$\Theta_k(f)(e^{wY_k}[1])=f(x(w))$.
The independence and surjectivity in OMG9--13 prove that this
map is onto $E_k^\vee$, with kernel $\mathfrak a_k$.
Consequently it induces a specified vector-space isomorphism
\[
 \boxed{\Theta_k:R_k\xrightarrow{\ \sim\ }E_k^\vee.}
 \tag{OMG19}
\]
All phases and factorials are those of OCS17; for a monomial
$x^\nu$ its coefficients are precisely the finite multinomial
OCS15 rule with the entries OMG4.

The collection of these vector spaces has an exact graded algebra
structure coming from the original sum maps. For $k,j\ge0$ put
\[
 \Delta_{k,j}:E_{k+j}\longrightarrow E_k\otimes E_j,
               \qquad[P(S)]\longmapsto[P(S_1+S_2)].
 \tag{OMG20}
\]
On a primary pair the eigenvalue is the displayed sum
$\beta_{k,a,b}+\beta_{j,a',b'}=\beta_{k+j,a+a',b+b'}$.
The sum of the two nilpotents has index $e_k+e_j-1=e_{k+j}$:
its power $e_k+e_j-2$ has nonzero top coefficient
$\binom{e_k+e_j-2}{e_k-1}$ and its next power is zero. Every
pair of total sign counts occurs. Thus the full minimal polynomial
of $S_1+S_2$ is exactly $\chi_{k+j}$, proving both well-definedness
and injectivity of OMG20. It also gives the complete primary formula
\[
 \Delta_{k,j}H_{k+j,A,B,D}
 =\sum_{\substack{a+a'=A,\ b+b'=B}}
   \sum_{\substack{d+d'=D\\0\le d<e_k,\ 0\le d'<e_j}}
       \binom Dd H_{k,a,b,d}\otimes H_{j,a',b',d'}.
 \tag{OMG21}
\]
Indeed the total primary idempotent is one precisely on the
matching eigenvalue pairs, and the nilpotent power expands by the
ordinary binomial theorem there. All other components vanish.
Substitution of the sum of three coordinates proves coassociativity;
interchange of the two coordinates proves cocommutativity. The
$k=0$ maps give the counit. Hence
\[
 (\ell\star\ell')(v)=(\ell\otimes\ell')(\Delta_{k,j}v)
       \quad(\ell\in E_k^\vee,\ \ell'\in E_j^\vee)
 \tag{OMG22}
\]
is an associative commutative graded product with unit in $E_0^\vee$.
For the exact dual basis $\lambda_{k,a,b,D}=i^DD!H_{k,a,b,D}^\vee$,
OMG21 gives
\[
 \lambda_{k,a,b,D}\star\lambda_{j,a',b',D'}
       =\lambda_{k+j,a+a',b+b',D+D'}.
\]
The coefficient is exactly
$\binom{D+D'}D i^DD!\,i^{D'}D'!=i^{D+D'}(D+D')!$;
this records the cancellation before stating the basis product.
Equivalently, OMG20 sends $e^{wY_{k+j}}[1]$ to
$e^{wY_k}[1]\otimes e^{wY_j}[1]$, with
$c_{k+j}=c_k+c_j$. Evaluating OMG18 on this identity proves
\[
 \boxed{\Theta:\ R\xrightarrow{\ \sim\ }
                    \bigoplus_{k\ge0}(E_k^\vee,\star)
                  \text{ is an isomorphism of graded algebras}.}
 \tag{OMG23}
\]
Its source multiplication is the actual polynomial multiplication;
its target multiplication has been constructed from the original sum
coalgebra rather than assigning a multiplication between unrelated
finite-dual vector spaces.

\subsection{The full semisimple and nilpotent action in this geometry}
On $\mathscr P_Z$ define the degree-preserving derivation
\[
 \partial_Z Z_{a,b,d}=\omega_{2a-1,2b-1}Z_{a,b,d}
                   +\begin{cases}d Z_{a,b,d-1},&d>0,\\0,&d=0.
                     \end{cases}
 \tag{OMG24}
\]
On the parameter ring OMG7 it is induced by
$\partial X_a=-i(2a-1)\delta X_a$,
$\partial Y_b=(2b-1)\gamma Y_b$,
$\partial T_0=0$, $\partial T_1=T_0$.
The product rule verifies OMG24 and commutation with $\Phi_m$;
therefore the complete ideal $\mathfrak t_m$ is stable.
The corresponding exact flow on the parameters is
\[
 (X_a,Y_b,T_0,T_1)\longmapsto
 (e^{-i(2a-1)\delta w}X_a,e^{(2b-1)\gamma w}Y_b,T_0,T_1+wT_0).
\]
At $(X_0,X_1,Y_0,Y_1,T_0,T_1)=(1,1,1,1,1,0)$ it gives precisely
the original functions OMG5. This is the entire stated analytic
flow on the explicit algebraic variety, with all polynomial degrees.

If $B_Z$ is the matrix in OMG24 on the column of variables, its
rows have diagonal $\omega_\alpha$ and entry $d$ in column
$(\alpha,d-1)$ of row $(\alpha,d)$. In actual coordinates the
derivation is
\[
 \partial_x x=\mathcal A x,\qquad
 \mathcal A=\mathsf H_mB_Z\mathsf H_m^{-1}.
 \tag{OMG25}
\]
The inverse substitution proves $\partial_x\mathfrak a\subseteq
\mathfrak a$, and its flow differentiates the original $x(w)$.
On the exact degree-$k$ basis of OMG9/OMG18 the induced action is
\[
 \partial_x\lambda_{k,a,b,D}
   =\omega_{k,a,b}\lambda_{k,a,b,D}
             +\begin{cases}D\lambda_{k,a,b,D-1},&D>0,\\0,&D=0.
              \end{cases}\qquad
 \Theta_k\partial_x=Y_k^\vee\Theta_k.
 \tag{OMG26}
\]
The displayed action on the basis uses its identification OMG23.
For a direct verification, on the original primal primary basis
$Y_kH_D=\omega H_D+i^{-1}H_{D+1}$ with the last term zero at
$D=e_k-1$. Applying $i^DD!H_D^\vee$ proves the exact lowering
coefficient $D$ in its dual action. Each frequency block has all
$e_k$ terms, with nonzero top iterated derivative
$(e_k-1)!$. Thus every original primary nilpotent has exactly the
same index $e_k$ in this model. The full annihilating polynomial
is retained with its leading phase:
\[
 \chi_k(c_k+i\partial_x)=0\quad\text{on }R_k,\qquad
 \chi_k(c_k+iT)=i^{q_k}\prod_{a,b=0}^k(T-\omega_{k,a,b})^{e_k}.
 \tag{OMG27}
\]
Distinct frequencies and the nonzero top chain prove that this
degree $q_k$ is minimal. The original $S$ action corresponds to
$c_k+i\partial_x$ on degree $k$, or to
$\tfrac12\mathsf E+i\partial_x$ on the graded algebra, where
$\mathsf E$ is the degree derivation. These equalities specify
every constant, sign and nilpotent action under OMG23.

\subsection{The exact original cyclic observation as invariant restriction}
Let $D=\operatorname{diag}(\zeta,\ldots,\zeta^n)$ and let
$G=\mathbb Z/N$ act on $\mathscr P_x$ by
$(g f)(x)=f(D^{-1}x)$. Thus a degree-$k$ monomial $x^\nu$
has character $\zeta^{-\sum t\nu_t}$. Write
\[
 \mathscr I=\mathscr P_x^G,\qquad
 \mathscr J=\operatorname{im}(\mathscr I\longrightarrow R).
 \tag{OMG28}
\]
These are graded algebras. The coefficient image of the original
observation is exactly their restriction at every degree, as follows.
Let $f_t$ be column $t$ of $F$ and let $S_\nu$ be the sum of
all distinct ordered words in the $f_t$ having multiplicity vector
$\nu$. Disjoint word supports make the $S_\nu$ a basis of the
permutation-invariant tensor space. The expansion and projector are
\[
 (Fx)^{\otimes k}=\sum_{|\nu|=k}x^\nu S_\nu,\qquad
 P_k(Fx)^{\otimes k}=\sum_{\substack{|\nu|=k\\\sum t\nu_t=0\pmod N}}
                    x^\nu S_\nu,\quad
 P_k=\frac1N\sum_{a=0}^{N-1}(C^{-a})^{\otimes k}.
\]
Every ordered word occurs once; the factor $k!/\prod\nu_t!$ is
the number of terms in $S_\nu$, not a missing coefficient.
The original sum injection is the iterated map OMG20 into
$E_h^{\otimes k}$. Hence the original observation
$L_k=P_k\Pi^{\otimes k}U^{\otimes k}\iota_k=\jmath\Lambda_k$
sends $e^{wY_k}[1]$ to the displayed projected curve with $x=x(w)$.
The coefficient covectors of its surviving $S_\nu$ pull back to
exactly $\Theta_k(x^\nu)$. They span the full dual of the image,
since every covector on a subspace of a finite dimensional space
extends to its ambient space. Consequently
\[
\boxed{\begin{gathered}
 \Theta_k(\mathscr J_k)=\operatorname{im}\Lambda_k^\vee
                          =(\ker\Lambda_k)^\perp,\qquad
 \operatorname{rank}\Lambda_k=\dim\mathscr J_k,\\
 R_k/\mathscr J_k\xrightarrow{\ \Theta_k|_{K_k}\ }K_k^\vee
       \text{ is an isomorphism of vector spaces},\qquad
 K_k=\ker\Lambda_k.
\end{gathered}}\tag{OMG29}
\]
For the last statement, a covector on $E_k$ vanishes on $K_k$
exactly when it factors through $E_k/K_k\simeq\operatorname{im}
\Lambda_k$. Restriction to $K_k$ is onto by extension of a basis.
This proves the claimed quotient and its precise type: the
degree-$k$ quotient is a vector-space quotient by $\mathscr J_k$.

\subsection{The actual orbit ideal, including every stabilizer}
The whole ideal $\mathfrak a$ is generated by the explicit quadratic
space in OMG15. Let $B_2$ be its coefficient matrix in the ordered
degree-two monomial frame of $\mathscr P_x$. Repeated or dependent
generator rows are harmless. For $a\in\mathbb Z/N$ let
$D_{2,a}$ be diagonal with entry $\zeta^{-a\sum t\nu_t}$ at
monomial $x^\nu$. The actual stabilizer and its exact finite test are
\[
 H=\{a:g^a\mathfrak a=\mathfrak a\}
   =\left\{a:\operatorname{rank}\begin{bmatrix}B_2\\B_2D_{2,a}\end{bmatrix}
                      =\operatorname{rank}B_2\right\}.
 \tag{OMG30}
\]
Equality of these row spaces is equivalent to equality of their
quadratically generated ideals; their dimensions agree under the
invertible diagonal map. This proves the finite test using the
actual entries OMG4. It prescribes no generic coefficient value.
The set $H$ is a subgroup by composition and inverse of ideal
equalities. Put $s=|H|$ and $d=N/s$. Since $G$ is cyclic, its
cosets have representatives $0,\ldots,d-1$ and
$H=\{0,d,2d,\ldots,(s-1)d\}$. Define all distinct component ideals
and the union algebra by
\[
 \mathfrak a_a=g^a\mathfrak a\ (0\le a<d),\qquad
 \mathfrak u=\bigcap_{a=0}^{d-1}\mathfrak a_a,
 \qquad A=\mathscr P_x/\mathfrak u.
 \tag{OMG31}
\]
These are distinct homogeneous prime ideals with the same Hilbert
function OMG13. Their intersection is a homogeneous radical ideal
stable under the original cyclic action. It is the complete ideal
of their scheme-theoretic union by the definition of the union's
ideal. No product formula for this intersection is substituted.

For any invariant polynomial, membership in one component ideal
implies membership in all its translates, by applying every $g^a$.
Thus the exact invariant restriction and its inverse description are
\[
\boxed{\begin{gathered}
 \mathscr I\cap\mathfrak a=\mathscr I\cap\mathfrak u,\\
 \mathscr I/(\mathscr I\cap\mathfrak u)
      \xrightarrow{\ [f]\mapsto[f]_{\mathfrak u}\ } A^G
      \xrightarrow{\operatorname{res}_0} \mathscr J\subseteq R
      \quad\text{are isomorphisms of graded algebras}.
\end{gathered}}\tag{OMG32}
\]
To prove surjectivity of the first map, average any representative
of an invariant class by $N^{-1}\sum g^a$; its class is unchanged.
Its kernel is exactly the displayed intersection. For the second,
use an invariant representative: a zero restriction lies in
$\mathscr I\cap\mathfrak a=\mathscr I\cap\mathfrak u$ and hence
is already zero in $A^G$. Its image is $\mathscr J$ by definition.
This gives the invariant image for every actual $H$, including
nontrivial stabilizers of a composite group order $4m+1$.

\subsection{The complete component conductor and an explicit finite witness}
Set $R_a=\mathscr P_x/\mathfrak a_a$. The simultaneous restriction
\[
 A\hookrightarrow\widetilde A:=\prod_{a=0}^{d-1}R_a,
 \qquad[f]_{\mathfrak u}\longmapsto([f]_{\mathfrak a_a})_a
 \tag{OMG33}
\]
is injective precisely because its kernel before quotienting is the
intersection OMG31. Its component conductor ideals are
\[
 \mathfrak c_a=
 \frac{(\bigcap_{b\ne a}\mathfrak a_b)+\mathfrak a_a}{\mathfrak a_a}
       \subseteq R_a,
 \qquad\bigcap_{b\ne a}\mathfrak a_b=\mathscr P_x\quad(d=1).
 \tag{OMG34}
\]
The conductor of this specified inclusion is exactly
\[
 \boxed{\{t\in\widetilde A:t\widetilde A\subseteq A\}
                   =\prod_{a=0}^{d-1}\mathfrak c_a\subseteq A.}
 \tag{OMG35}
\]
Indeed a tuple supported solely on component $a$ lies in the image
of $A$ if and only if its value belongs to $\mathfrak c_a$, by
lifting its other zero values to the intersection of the other
ideals. If $t$ is in the conductor, multiplication by each component
idempotent of $\widetilde A$ gives such a supported tuple; thus each
component of $t$ lies in $\mathfrak c_a$. Conversely these are
ideals of $R_a$. Multiplying a tuple with all its components in
the indicated ideals by any element of $\widetilde A$ gives a sum
of such supported tuples, each lying in $A$. This proves both
inclusions and the full formula. In particular this conductor is
also the annihilator of $\widetilde A/A$ in $A$.

The action of $G$ permutes these components with its original
coefficient substitutions. Projection of an invariant tuple to
component zero gives an isomorphism
\[
 \widetilde A^G\xrightarrow{\ \sim\ }R^H,\qquad
 r\in R^H\longmapsto(g^a r)_{0\le a<d}
                       \text{ for its inverse}.
 \tag{OMG36}
\]
The inverse is independent of coset representatives because $r$
is fixed by $H$, and invariance follows by reindexing the components.
The same argument applies to the conductor, whose ideals are
permuted. Equations OMG32--36 consequently prove the exact
inclusions of graded subalgebras and conductor ideals
\[
 \boxed{\mathfrak c_0^H\subseteq\mathscr J\subseteq R^H,
            \qquad\mathfrak c_0^H\,R^H\subseteq\mathscr J.}
 \tag{OMG37}
\]
Here $\mathfrak c_0^H$ denotes invariant elements of the specified
component ideal and is an ideal of $R^H$. An invariant tuple in
the conductor already lies in $A$, so its projection belongs to
$\mathscr J$ by OMG32. This proves the first inclusion; the
ideal property gives the second.

There is a nonzero homogeneous element of this conductor with a
completely bounded degree. For every $a\ne0$ choose from the
explicit transformed quadratic list of $\mathfrak a_a$ a quadratic
$f_a\notin\mathfrak a_0$. Such a choice exists: containment of
the entire quadratic space in that of $\mathfrak a_0$, which has
the same dimension, would make the spaces equal and their
generated ideals equal, contrary to distinctness. Membership is
the finite row-space test in the same degree-two monomial frame.
The product $f=\prod_{a=1}^{d-1}f_a$ lies in all other component
ideals. It is nonzero modulo $\mathfrak a_0$ because that ideal is
prime and none of its factors belongs to it. Thus
$\bar f\in\mathfrak c_0$ is nonzero of degree $2(d-1)$.
The group $H$ preserves $\mathfrak c_0$ and acts on the domain
$R$. Its exact norm
\[
 c_H=\prod_{h\in H}g^h\bar f\in\mathfrak c_0^H,\qquad
 \deg c_H=D_H=2s(d-1)=2(N-s),\qquad c_H\ne0
 \tag{OMG38}
\]
is invariant by permutation of its factors and nonzero in the
domain. For $d=1$ take the empty product $c_H=1$, $D_H=0$;
all statements remain exact. Multiplication by this actual
conductor element is injective and gives the finite dimension
bounds, with negative graded pieces set to zero,
\[
\boxed{\begin{gathered}
 (R^H)_{k-D_H}\xrightarrow{\ \times c_H\ }\mathscr J_k
                    \hookrightarrow(R^H)_k,\\
 \dim(R^H)_{k-D_H}\le\operatorname{rank}\Lambda_k
                                  \le\dim(R^H)_k,\\
 q_k-\dim(R^H)_k\le\dim K_k
                                  \le q_k-\dim(R^H)_{k-D_H}.
\end{gathered}}\tag{OMG39}
\]
Every ideal, group and matrix in these expressions is defined by
the actual original coefficient map. In particular the stabilizer
has its finite equality test OMG30 throughout; no particular value
of it has been inserted into the dimensions. To compute the displayed
invariant dimensions directly, let the actual $H$ substitutions act
on the basis OMG9 through $Z\mapsto\mathsf H_m^{-1}D^{-h}\mathsf H_mZ$
and reduce with OMG10. This is an explicit finite matrix on $q_k$
coordinates. The idempotent average $s^{-1}\sum_{h\in H}g^h$ has
trace equal to its rank, proving
$\dim(R^H)_k=s^{-1}\sum_{h\in H}\operatorname{Tr}(g^h|R_k)$.
Thus the dimensions in OMG39 also have complete original finite
coefficient formulas for every actual stabilizer.

\subsection{The exact action defect and all subsequently observed directions}
The derivation $\partial_x$ of OMG25 acts on $R$ with all primary
directions OMG26. The possible failure of its restriction to the
observed image is the specified finite linear map
\[
 \mathscr J_k\longrightarrow R_k/\mathscr J_k,
                       \qquad j\longmapsto[\partial_xj].
 \tag{OMG40}
\]
Via OMG29 it is exactly the complex linear dual of the original
map $\Lambda_kY_kI_K:K_k\to B_k$. Indeed a $j$ corresponds to
the covector $\lambda\Lambda_k$ on $E_k$ for a unique boundary
covector $\lambda$; evaluating OMG26 on $I_Kv$ gives
$\lambda\Lambda_kY_kI_Kv$, proving the identification and all
domains and codomains. The action on ambient polynomials also has
the exact conjugation formula
\[
 g^a\partial_xg^{-a}=\partial_{D^a\mathcal A D^{-a}},\qquad
 g^a\partial_xg^{-a}-\partial_x
                   =\partial_{D^a\mathcal A D^{-a}-\mathcal A}.
 \tag{OMG41}
\]
For a linear vector field $\partial_M f=(Mx)\cdot\nabla f$,
the chain rule proves this by substituting $D^{-a}x$ after
differentiating $f(D^a x)$. This formula retains every original
period/unit coefficient in $\mathcal A$ and does not require the
original ideal or the observed image to be preserved by extra
ambient actions.

Finally the subspace invisible to all original iterates has the
complete model
\[
\boxed{\begin{gathered}
 \mathcal K_{\infty,k}=\bigcap_{r=0}^{q_k-1}\ker(\Lambda_kY_k^r),\\
 R_k\Big/\sum_{r=0}^{q_k-1}\partial_x^r\mathscr J_k
       \xrightarrow{\ \Theta_k|_{\mathcal K_{\infty,k}}\ }
                          \mathcal K_{\infty,k}^\vee
                          \quad\text{is a vector-space isomorphism}.
\end{gathered}}\tag{OMG42}
\]
In finite dimensional duality the annihilator of an intersection
of kernels is the sum of the images of their dual maps: this
follows by taking the dual of the single stacked map with all
those components, whose kernel is precisely the intersection.
Equations OMG26 and OMG29 identify these images with
$\partial_x^r\mathscr J_k$, proving OMG42. The exact original
polynomial OMG27 expresses every higher iterate as a linear
combination of the first $q_k$, so the same intersection is taken
over all nonnegative iterates and is $Y_k$-invariant. This proves
the full nilpotent and cyclic-action receiver in the new homogeneous
model for every original multiplicity, with its actual observation,
periods and amplitude jets throughout.

% END COMPLETE OMG

\end{document}
